% A.1--A.2 and B.1--B.5.  Numbered statements are explicit for the result register.
\subsection{A.1. The integral quotient and its three-dimensional target}

We write $g,u,w,d$ for the ordered homology basis
$(\widehat\gamma,\widehat u,\widehat w,\widehat\delta)$ in this
section, and $\gamma$ for its first dual covector.  Thus $d$ is a
vector, whereas $\gamma$ is a covector.  The cusp coordinate is
$q=1/t$.  Positive target orientations are the complex orientation
on the base followed by the positive $\gamma$-circle orientation.
The complex orientation on $X$ fixes the source orientation.

\paragraph{Proposition A.1.1 (Proved).}
For the markings and the clockwise root generators in the prompt,
the homological matrices and twisting vectors are
\[
 A_1=\begin{pmatrix}1&0&0&0\\6&0&1&0\\-6&-1&-1&0\\-2&1&0&1\end{pmatrix},
 \qquad
 A_2=\begin{pmatrix}1&0&0&0\\0&0&-1&0\\-6&1&0&0\\3&0&1&1\end{pmatrix},
\]
\[
 M_0=\begin{pmatrix}1&0&0&0\\0&1&0&0\\0&1&1&0\\-1&0&0&1\end{pmatrix},
 \qquad v_1=g+2u-4w,\qquad v_2=-g-3u+3w.
\]
The common invariant vectors and covectors are $\mathbb Zd$ and
$\mathbb Z\gamma$, respectively.  The smallest invariant primitive
sublattice containing the cusp vanishing lattice is
$L=\mathbb Z\langle u,w,d\rangle=\ker\gamma$.
The formula
\[
 [z,\Pi(z)x]\longmapsto([z],\gamma(x)\bmod\mathbb Z)
 \tag{A.1.1}
\]
defines a real-analytic fibrewise group epimorphism
$J\longrightarrow B^\circ\times S^1$ with connected kernel $T_L$.
At the finite fillings it descends to a smooth proper $T^3$-bundle
over a Seifert solid torus, of multiplicities $3$ and $4$ and
translation residues $1$ and $-1$.

\begin{proof}
The first two displayed matrices are the inverse transposes of
the matrices printed on page 2 of \cite{mapsAC26}; their product
has inverse $M_0$.  Direct integral multiplication gives
$A_1^3=A_2^4=A_1A_2M_0=I$ and $A_jv_j=v_j$.
Solving the two fixed-vector equations gives $(0,0,0,d_0)$.
For a row vector $(a,b,c,e)$, the first fixed-covector equation
gives $b=2c,e=3c$ and the second gives $b=c,e=2c$;
therefore $b=c=e=0$.  Also
\[
 (M_0-I)(a,b,c,e)=(0,0,b,-a),\qquad
 A_1w=u-w.
\]
These identities prove the assertions about $L$ and the cusp
lattice $K=\mathbb Z\langle w,d\rangle$.  Their primitivity follows
from the displayed bases.  The flat real marking identifies the
deck action on a fibre with $x\mapsto A_jx$.  Since
$\gamma A_j=\gamma$, (A.1.1) descends.  Its kernel is connected
because $\gamma$ is primitive.

On a root disc the affine generator is
\[
 (s,x)\longmapsto
 (e^{-2\pi i/m}s,A_jx+v_j/m),\qquad (m,\kappa)=(3,1),(4,-1).
\]
It induces $(s,\theta)\mapsto
(e^{-2\pi i/m}s,\theta+\kappa/m)$ on the character quotient.
This action is free since $\gcd(\kappa,m)=1$.  Writing
$a=e^{2\pi i\theta}$ and choosing $r$ with $r\kappa\equiv1\pmod m$,
the invariant coordinates $s a^r$ and $a^m$ identify its quotient
with $D^2\times S^1$.  The base coordinate becomes
$(s a^r)^m(a^m)^{-r}$, and hence has multiplicity $m$ at its
central circle.  Before the finite quotient the map is a trivial
$T_L$-bundle.  The free actions on source and target are
equivariant for that map, so it descends locally to a smooth
$T_L$-bundle.  Properness follows from compactness of $T_L$.
The logarithmic translation used to glue the root chart to $J$
changes the character by its $\gamma$-component.  Applying
$\gamma$ to the actual gluing translation therefore gives the
target gluing, proving agreement on every overlap.
The exact checks, including the exhaustive mod-two fixed-vector
calculation used below, are in \texttt{code/maps\_exact.py}.
\end{proof}


Engel's primitive invariant covector is named $\psi$ in his
Lemma 3.8, labelled \texttt{covect} in the supplied source
\cite{mapsEngel}.  It is not named $\varepsilon$ there.  In his
common marking $\mathcal B_0$, the integral comparison is
\[
 P_E=\begin{pmatrix}0&-2&-1&-2\\2&-1&0&-1\\0&2&1&3\\-1&0&0&0\end{pmatrix},
 \qquad \psi P_E=-\gamma.
 \tag{A.1.2}
\]
Thus the quotient characters agree after reversing the target
circle orientation.  With the inverse root generators the
translation vectors are $a_{j-1}=-P_Ev_j/m_j$.
This is the precise sign convention in the supplied family
return, B.2.2; it explains why equality of the unnamed primitive
covectors without a sign convention would be insufficient.

\paragraph{Proposition A.1.2 (Refuted as stated; corrected).}
The literal character has no continuous extension at either
triple point or at generic points of two of the three double
curves.  It has the constant limit $0$ on the smooth locus of
$W$ and on the third double curve away from the triple points.
In the chart $q=z_0z_1z_2$, $x_1=z_1$, $x_2=z_2$, the extendible
double curve is $z_0=z_1=0$, $z_2\ne0$.  The fixed curve of the
$\mathbb C^*$-action in direction $d=e_2$ is instead
$z_0=z_2=0$, $z_1\ne0$, and is one of the nonextendible curves.
Consequently the parenthesis in A.1 identifying the extendible
curve with the fixed curve is false.

\begin{proof}
Put $q=e^{-2\pi T}$ on a positive radial ray and use the exact
cusp periods
\[
 Z_T=iTB_0+C(q),\qquad
 B_0=\begin{pmatrix}0&1\\-1&0\end{pmatrix},
\]
where $C$ is bounded and holomorphic at $0$.  For a fixed
$a\in\mathbb R^2$ the representative
$x_T(a)=\exp(2\pi iZ_Ta)$ has character $a_1\bmod\mathbb Z$ and
logarithmic position
\[
 \frac{\log|x_T(a)|}{\log|q|}
 =B_0a+T^{-1}\operatorname{Im}C(q)a.
 \tag{A.1.3}
\]
The choices $a=(-1/4,1/4)$ and $a'=(-1/2,1/4)$ have logarithmic
positions tending to $(1/4,1/4)$ and $(1/4,1/2)$.
Both are interior to the triangle with vertices $0,e_1,e_2$,
so all three chart coordinates tend to $0$.  Their character
values are $3/4$ and $1/2$.  This proves nonextension at the
first triple point.  The other triangle type has vertices
$e_1,e_2,e_1+e_2$.  Its coordinates can be taken as
$x_1x_2/q,q/x_2,q/x_1$.  Positions $(3/4,3/4)$ and
$(3/4,1/2)$ lie in its interior, give distinct values of $-y_2$
modulo one, and converge to its origin.  Translation by the
lattice preserves triangle type; these are precisely the two
triple points of $W$.

For arbitrary convergent paths in the first chart, inversion
of the exact logarithmic period matrix gives
\[
 \theta_\gamma=-\frac{\log|z_2|}{\log|q|}+o(1)\pmod{\mathbb Z}.
 \tag{A.1.4}
\]
Here the error tends to zero whenever the logarithmic ratios
stay bounded; they do stay bounded near a fixed point of a
normal-crossing stratum because $q=z_0z_1z_2$ and the three
coordinates are bounded.  If $z_0,z_1\to0$ and $z_2\to c\ne0$,
the right side tends to zero.  If $z_0,z_2\to0$ and $z_1=c\ne0$,
the paths $z_2=q^a,z_0=q^{1-a}/c$, $0<a<1$, have limits
$-a\pmod1$ at the same double point.  The same construction
with $z_0=c$ proves nonextension on $z_1=z_2=0$.
At a smooth point of a component the ratio in (A.1.4) tends
to $0$ or $1$, giving zero modulo one.  Integer translates
of these charts give the stated assertion on all of $W$.

The torus action in direction $e_2=d$ is
\[
 \lambda\cdot(z_0,z_1,z_2)
       =(\lambda^{-1}z_0,z_1,\lambda z_2).
\]
Its fixed curve is $z_0=z_2=0$.  This directly locates the
incorrect identification in the prompt.  Edge direction,
torus stabilizer, and coordinate on the quotient polygon must
not be interchanged.
\end{proof}

The nonextension above is local and pointwise.  It does not
contradict extension of the cohomology class over a cusp
neighbourhood.  The toric covering has rays $(v,1)$ for all
$v\in\mathbb Z^2$.  Meridians about its toric divisors kill
the whole cocharacter lattice: $(0,1)$ is a ray and
$(e_i,1)-(0,1)=(e_i,0)$.  The toric cover is therefore simply
connected.  Its free period deck group is $\overline\Lambda$
and hence
\[
 \pi_1(N_0)=\overline\Lambda=\Lambda/K.
 \tag{A.1.5}
\]
For the punctured cusp the map to this group kills $K$ and
the base meridian.  The latter bounds the disc given by a
point of the dense orbit on a central component.  Both
long-period covectors extend in integral $H^1$, since they
kill precisely these relations.  In contrast, an extension
of $\theta_\gamma$ to all of $X\simeq S^6$ would induce a
nonzero map $H_1(F)\to\mathbb Z$ factoring through
$H_1(X)=0$, which is impossible.  This last argument is the
global cohomological obstruction K5; it says nothing about
pointwise cusp limits.

\paragraph{Proposition A.1.3 (Proved).}
For the original data the quotient over the complement of
the cusp disc is the exterior of the $(3,4)$ torus knot,
with the specified Seifert structure.  For general admissible
$(l_0,l_1,l_2)$ the filled quotient is the lens space determined
by the following meridian slopes in one oriented integral
torus basis $(\rho_1,h)$:
\[
 \mu_1=(3,-l_1),\qquad
 \mu_2=(-4,4l_0-l_2),\qquad
 \det(\mu_1,\mu_2)=p=12l_0-4l_1-3l_2.
 \tag{A.1.6}
\]
In the lens-space convention that the second meridian is
$-q_0\mu_1+p\lambda_1$, where
$\det(\mu_1,\lambda_1)=1$, the lens parameter satisfies
$3q_0\equiv4\pmod p$.  In particular the target depends only
on the signed integer $p$, and is a three-sphere exactly
when $|p|=1$.

\begin{proof}
Let $P$ be the three-holed sphere, with oriented boundary
meridians $\rho_0+\rho_1+\rho_2=0$ and positive fibre $h$.
The two finite meridian relations are
$3\rho_1-l_1h=0$ and $4\rho_2-l_2h=0$; the cusp solid-torus
meridian is $\rho_0+l_0h$.  Eliminating $\rho_0$ gives
$\rho_2=l_0h-\rho_1$ on the annulus joining the finite solid
tori, which gives exactly (A.1.6).  Since both slopes are
primitive, this is a genus-one Heegaard splitting.  Choose
$a,b\in\mathbb Z$ with $al_1+3b=1$ and put
$\lambda_1=(a,b)$.  Direct multiplication gives
\[
 \mu_2=-q_0\mu_1+p\lambda_1,
 \quad q_0=a(4l_0-l_2)+4b,
 \quad 3q_0-ap=4.
 \tag{A.1.7}
\]
This proves the claimed lens parameter in a convention fixed
by actual meridians, including orientations.  The fundamental
group of this gluing is $\mathbb Z/|p|$.  Admissibility implies
$\gcd(p,6)=1$, so the modular inverse of $3$ used above exists.

For $(l_0,l_1,l_2)=(0,1,-1)$ the two exceptional-fibre data
are $(3,1),(4,-1)$.  They are realized by the circle action
\[
 s\cdot(z_1,z_2)=(e^{4is}z_1,e^{3is}z_2)
 \quad\text{on } S^3\subset\mathbb C^2.
\]
Its invariant map is $[z_1^3:z_2^4]$; the local circle
translation residues are $4^{-1}=1\pmod3$ and
$3^{-1}=-1\pmod4$.  Removing a regular orbit gives the
same two Seifert solid tori and their product annulus as
above.  A regular orbit is the $(3,4)$ torus knot, up to
the harmless exchange of the two coordinates.  This
identifies the manifold with its peripheral structure;
the abstract presentation $\langle x,y\mid x^3=y^{-4}\rangle$
alone would not have justified that identification.
\end{proof}

\subsection{A.2. The repaired map and the scope of its uniqueness}

\paragraph{Theorem A.2.1 (Proved).}
For every admissible member, there is a smooth map
$e:X\to Q$ over $B$ whose restriction outside an arbitrarily
small cusp neighbourhood is the literal circle quotient.
For fixed quotient gluing, its homotopy class is independent
of the cusp replacement, provided the replacement maps the
cusp neighbourhood into its designated solid torus and is
kept fixed on its outer collar.  In particular there is a
well-defined class for the original member $X\simeq S^6$.

\begin{proof}
On the cusp boundary use the product coordinates
$(q,\theta_\gamma)$ after the specified collar shear.
By (A.1.5), the boundary circle class extends to $N_0$.
Since $S^1=K(\mathbb Z,1)$, it has a continuous representative
$\theta':N_0\to S^1$.  Its restriction to the boundary is
homotopic to the given one: the two maps induce the same
homomorphism on $\pi_1$ of the boundary mapping torus.
The homotopy-extension property for a manifold with boundary
changes $\theta'$ on a collar so that it agrees exactly with
$\theta_\gamma$ there.  Relative smooth approximation gives
a smooth representative, still equal on a smaller collar.
Then $(f,\theta')$ is the required map into $D^2\times S^1$,
and agrees with the finite-filling maps on the overlap.
This proves smoothness even at the critical points of $f$:
the product map is smooth without any submersion assertion.

The inclusion $\partial N_0\to N_0$ surjects on $\pi_1$,
as (A.1.5) explicitly shows.  It follows that
$H^1(N_0;\mathbb Z)\to H^1(\partial N_0;\mathbb Z)$
is injective.  Since both spaces are connected, the long
exact sequence gives $H^1(N_0,\partial N_0;\mathbb Z)=0$.
Two circle replacements agreeing on the boundary therefore
have quotient whose relative $H^1$ class is zero, hence
have a relative homotopy.  Keeping the base coordinate
fixed gives the required homotopy of maps into the solid
torus.  Different choices of the collar size are compared
on a common smaller collar.  This proves the assertion.
\end{proof}

\paragraph{Proposition A.2.2 (Refuted as an unrestricted uniqueness assertion).}
Agreement with the literal quotient outside a neighbourhood
of $W$, without a condition on the image of that neighbourhood,
does not determine a homotopy class $S^6\to S^3$.  All twelve
classes occur among maps satisfying that bare agreement
condition.

\begin{proof}
Take a closed embedded six-ball in the interior of the cusp
neighbourhood.  Homotope a given extension on a smaller ball
to be constant there, keeping its boundary and the complement
fixed; a target coordinate ball around the image of its centre
provides the contraction.  Pinch the smaller ball to
$S^6\vee D^6$ relative to its boundary and insert any based
map $\alpha:S^6\to S^3$ on the new sphere.  The resulting
global map agrees with the original map off that ball, and
represents $[e]+\alpha$ under $[S^6,S^3]=\pi_6(S^3)$.
Distinct $\alpha$ give distinct classes.  Smooth relative
approximation preserves this conclusion.  The insertion
cannot be performed within the cusp solid torus, since that
solid torus has $\pi_6=0$; this is exactly the distinction
in Theorem A.2.1.
\end{proof}

The corrected content of \cite[Remark 8.3]{mapsEngel} is
therefore a quotient over the regular part and the two finite
fillings, followed by a replacement within its circle homotopy
class over a cusp collar.  The literal character does not
extend.  The replacement gives a canonical homotopy class
when the solid-torus image condition is retained.

\subsection{B.1. The framed three-torus and the class of $e$}

We use $\pi_3^{\mathrm S}=\mathbb Z/24\langle\nu\rangle$, with
$e_{\mathbb R}(\nu)=1/24$, and
$\pi_6(S^3)=\mathbb Z/12\langle\nu'\rangle$, where the
stabilization of $\nu'$ is $2\nu$.  We write
$\eta_n:S^{n+1}\to S^n$ for the $(n-2)$-fold suspension of
the complex Hopf map.  The order-two classes below are
unaffected by reversal of either source or target orientation.
An orientation-preserving diffeomorphism to the standard sphere
transports its bounding stable framing to $TX\oplus\mathbb R$.
Its homotopy class does not depend on that diffeomorphism:
two oriented stable trivializations differ by a map
$S^6\to\mathrm{SO}$, and $\pi_6(\mathrm{SO})=0$
by \cite[Corollary to Theorem II and the table on p.~315]{mapsBott59}.
We use outward stabilization and the convention that the normal
frame precedes the tangential frame; a constant permutation is
corrected once to preserve orientation.

\paragraph{Lemma B.1.1 (Proved).}
Let $F$ be a regular fibre of the original fibration.  Compare
the ambient stable framing of the oriented sphere $X$ with
the framing on $TX|_F$ consisting of four translation-invariant
vertical vector fields and an oriented lifted base frame.
The restriction of their stable difference to any coset of
$T_L\subset F$ is null-homotopic.

\begin{proof}
After a constant orientation adjustment the comparison is
a map $D:F\to\mathrm{SO}$, where stabilization is understood.
Its homotopy class is invariant under both monodromies.
Indeed, transport $F$ around a based loop in $B^\circ$ using
the flat real marking.  The translated vertical frame varies
continuously on $F\times[0,1]$; at the endpoint it changes by
the constant real matrix of the affine monodromy.  An oriented
base frame can be chosen periodic along the loop.  The space
of lifts of that frame to $TX$ is affine, so changing the lift
does not affect the comparison class.  The ambient stable
framing is already defined along this entire cylinder.  Thus
the comparison at its endpoints is $D$ and $D\circ A_j$
times a constant positive matrix.  Every such constant
matrix is path connected to the identity.  This proves
$[D]=A_j^*[D]$.  An affine translational part is isotopic
to the identity and has no effect.

The map on fundamental groups defines
$\sigma(D)\in H^1(F;\mathbb Z/2)$.  The exhaustive calculation
in \texttt{maps\_exact.py} gives
\[
 H^1(F;\mathbb Z/2)^{\langle A_1,A_2\rangle}
      =\{0,\gamma\bmod2\}.
\]
Both classes vanish on $L$, so $D|_{T_L}$ lifts to
$\mathrm{Spin}$.  Up to dimension four the latter has
$\pi_1=\pi_2=\pi_4=0$, $\pi_3=\mathbb Z$.
Obstruction theory consequently identifies maps from a
three- or four-dimensional torus to $\mathrm{Spin}$ with
$H^3$ of that torus.  The two possible lifts differ by
multiplication by the central element $-1$, which is
homotopic to the identity, so this identification has no
unrecorded lifting ambiguity.

The doubled stable difference $2[D]$ has zero $\sigma$,
and therefore corresponds to a monodromy-invariant class
in $H^3(F;\mathbb Z)$.  Integral Poincare duality, natural
for these determinant-one monodromies, identifies the
invariants in $H^3(F;\mathbb Z)$ with the fixed vectors
$\mathbb Zd\subset H_1(F;\mathbb Z)$.  Such a class
evaluates to zero on $T_L$, since $d\in L$:
its evaluation is the determinant with columns $d,u,w,d$.
Hence $2[D|_{T_L}]=0$ in $[T_L,\mathrm{Spin}]=\mathbb Z$.
This group has no torsion; thus $[D|_{T_L}]=0$.
The same argument applies to a coset since translation is
isotopic to the identity.
\end{proof}

\paragraph{Theorem B.1.2 (Proved).}
For the repaired quotient of Theorem A.2.1 with $|p|=1$,
\[
 [e]=6\nu'=\eta_3\eta_4\eta_5\ne0
       \quad\text{in }\pi_6(S^3)=\mathbb Z/12.
 \tag{B.1.1}
\]
Its stable class is $12\nu$ and its real Adams invariant
is $1/2\in\mathbb Q/\mathbb Z$.

\begin{proof}
Choose a point of $Q$ over a regular base value outside
the modified cusp disc and outside the two finite root
discs.  Its entire inverse image is one coset of $T_L$.
Order a target frame as the positive real base frame followed
by the positive $\gamma$-direction.  Lifting this frame
under $de$ gives a normal three-framing.  Complete it by
the constant Lie frame of $L$ with the unique constant
orientation adjustment making the full frame positive.
This is precisely the restriction of the product frame
in Lemma B.1.1, up to a constant positive change of basis.
That lemma proves that the stable tangential framing of
the inverse image is the Lie framing of $T^3$.
Pontryagin--Thom therefore gives
\[
 \operatorname{st}[e]=[T^3,\mathrm{Lie}]
 = [S^1,\mathrm{Lie}]^3=\eta^3=12\nu.
 \tag{B.1.2}
\]
The product equality follows by taking the product of the
three framed circles and their tubular collapse maps.
The Lie-framed circle is the nonbounding framed circle;
its class is the stabilized Hopf element $\eta$.

Here is an independent evaluation of the framed bordism
class, which does not use the ring relation $\eta^3=12\nu$.
Put the flat product metric on $\mathbb R^3/\mathbb Z^3$,
with the spin structure determined by its Lie framing.
Its spinors are periodic $\mathbb C^2$-valued functions.
On the Fourier mode $e^{2\pi i k\cdot x}$, $k\in\mathbb Z^3$,
the Dirac operator is the Hermitian matrix
$2\pi\sum k_j\sigma_j$, with eigenvalues
$\pm2\pi|k|$ for $k\ne0$ and a two-dimensional kernel
for $k=0$.  Thus its eta function is identically zero in
its convergent half-plane and hence at zero, while $h=2$.
The Levi-Civita connection and the connection induced
by the framing coincide, so their transgression term
vanishes.  The real formula of
\cite[Theorem (4.14), p.~423]{mapsAPS2}, with its factor
$\epsilon(1)=1/2$, gives
\[
 e_{\mathbb R}(T^3,\mathrm{Lie})
       =\frac12\frac{0+2}{2}=\frac12\pmod{\mathbb Z}.
 \tag{B.1.3}
\]
This detects $12\nu$ in $\mathbb Z/24$.  The complex
$e$-invariant would be zero here; omitting the real
factor would therefore lose the desired detection.

Finally stabilization is injective on $\pi_6(S^3)$.
For the exact unstable input we use
\cite[Theorem 4.1 and Corollaries 4.2--4.3]{mapsToda52}:
the suspended generator is the torsion term in
$[\iota_4,\iota_4]=2\nu_4-E\nu'$, and stabilization
sends it to $2\nu$.  In the group of order twelve this
is injective.  Therefore the class with stable image
$12\nu$ is $6\nu'$.  The composite
$\eta_3\eta_4\eta_5$ has the same stable image, so the
same injection identifies it with $6\nu'$.
\end{proof}

\subsection{B.2. The holomorphic map to the base}

\paragraph{Theorem B.2.1 (Proved).}
The map $f:X\to\mathbb P^1$ of the original sphere
construction is null-homotopic.  This conclusion does not
require the choice of $e$.

\begin{proof}
The quotient projection $h:Q\to B$ is, after the
orientation-preserving identifications fixed in A.1.3,
the weighted Hopf map
$[z_1^3:z_2^4]$, followed by a Mobius transformation
moving its two exceptional values to $0,1$.
It is the ordinary Hopf map after the self-map
\[
 P_{3,4}(z_1,z_2)=
 \frac{(z_1^3,z_2^4)}{\sqrt{|z_1|^6+|z_2|^8}}
 \quad\text{of }S^3.
 \tag{B.2.1}
\]
At a value with both coordinates nonzero, its phase
map has $3\cdot4=12$ preimages, all with positive local
degree; its radial coordinate is strictly increasing.
Consequently its degree is twelve.  A degree-$k$ map
of $S^3$ is homotopic to the $k$th power map of the
topological group $S^3$.  On $\pi_6(S^3)$ the latter
is multiplication by $k$, since pointwise multiplication
and the homotopy-group addition agree.  Thus
$h_* = (\eta_2)_*\circ12=0$ on $\pi_6(S^3)$.
The construction of $e$ is over $B$, so $f=h\circ e$
as actual maps.  This proves $[f]=0$.

For a second proof, let $S_1,S_2$ be the reduced multiple
fibres.  The line bundles $\mathcal O_X(S_j)$ are
topologically trivial because $H^2(X;\mathbb Z)=0$.
The coordinate sections of $f^*\mathcal O(1)$ have
divisors $3S_1$ and $4S_2$ after a Mobius change of
base coordinate.  Choose a smooth frame of
$f^*\mathcal O(1)$.  Since $H^1(X;\mathbb Z)=0$, every
map $X\to\mathbb C^*$ admits a logarithm.  We can
therefore choose frames of $\mathcal O_X(S_1)$ and
$\mathcal O_X(S_2)$ whose third and fourth powers,
under the given divisor isomorphisms, equal the chosen
frame.  The canonical divisor sections then become
smooth functions $a,b$ with no common zero and
$f=[a^3:b^4]$.  Put
$G=(a,b)/\sqrt{|a|^2+|b|^2}:X\to S^3$.
Writing $r=\sqrt{|a|^2+|b|^2}>0$, the explicit homotopy
$[a^3/r^{3s}:b^4/r^{4s}]$, $0\leq s\leq1$, gives
$f\simeq\eta_2P_{3,4}G$.  It is defined everywhere because
$a$ and $b$ have no common zero.  The latter map is zero
by the same degree calculation and the exponent twelve
of $\pi_6(S^3)$.  This proof uses neither the cusp
replacement nor the computed class of $e$.
\end{proof}

\subsection{B.3. The four workbench maps}

The exact maps, including their suspension conventions, are
needed here.  Put $\overline T=(\overline\Lambda\otimes\mathbb R)
/\overline\Lambda$ and let $\overline\pi:N_0\setminus W
\to\overline T$ be the long-period quotient.  On the punctured
cusp write $q=re^{i\phi}$.  The map $q_2$ takes $x$ to
$[r,e^{i\phi},\overline\pi(x)]$ in the unreduced suspension
of $S^1\times\overline T$, sends $W$ to its lower vertex,
and sends the complement of the cusp neighbourhood to its
upper vertex.  Collapsing the two-skeleton of
$S^1\times\overline T=T^3$ and then taking the reduced
suspension gives $H:X\to S^4$.  Collapsing the one-skeleton
of $\overline T=T^2$ and forgetting $\phi$ gives
$\widehat q_3:X\to S^3$.  Retaining only $\gamma$ and
forgetting $\phi$ gives $\widehat q_\gamma:X\to S^2$.
These are the maps in the supplied workbench, not choices
of arbitrary sphere quotients of the same targets.

\paragraph{Theorem B.3.1 (Proved).}
For the exact retained-angle map just defined,
\[
 [H]=\eta_4\eta_5\ne0\in\pi_6(S^4)=\mathbb Z/2.
 \tag{B.3.1}
\]
Its regular inverse image is the native vanishing
two-torus with its Lie stable framing and Arf invariant one.

\begin{proof}
Take $u_0=(-1/4,1/4)$, $q_0=r_0e^{i\pi}$, and
$p_0=B_0u_0=(1/4,1/4)$.  For $r_0$ sufficiently small
the full inverse image of the corresponding open-cell
value lies in the original chart of A.1.2.  To see this
without losing components, a lift of any point of that
inverse image has long coordinate $u_0+n$ for a unique
$n\in\mathbb Z^2$; applying the actual period deck
transformation $-n$ changes it to $u_0$.  Distinct lifts
with this coordinate cannot be deck equivalent.
Put $h_j(q)=\exp(-2\pi i(C(q)u_0)_j)$ and make the
holomorphic coordinate change
\[
 (w_0,w_1,w_2)
   =(z_0/(h_1h_2),h_1z_1,h_2z_2).
\]
It preserves the product $q$.  The full inverse image is
\[
 w(a,b)=
 (r_0^{1/2}e^{i(\pi-a-b)},
  r_0^{1/4}e^{ia},r_0^{1/4}e^{ib}),
 \quad a,b\in\mathbb R/2\pi\mathbb Z.
 \tag{B.3.2}
\]
All coordinates tend uniformly to zero as $r_0\to0$.
The period quotient is injective on a neighbourhood of
the chart origin, so this torus is contained in a
contractible coordinate ball of $X$.
The map $H$ is smooth on an open neighbourhood of this
torus, although its suspension presentation is only
continuous elsewhere.  Choose a smaller open-cell target
ball whose full inverse image lies in that neighbourhood;
compactness permits this because the displayed torus is
the full inverse image of its centre.  Relative smooth
approximation, uniformly small off that neighbourhood,
then gives a smooth representative with exactly the same
inverse image and differential there.  Pontryagin--Thom
therefore applies to the framing computed below.

We compute its framing before discarding the bounded
period terms.  Write
\[
 \ell=\log|q|,\quad
 R(q)=-2\pi\operatorname{Im}C(q),\quad
 A(q)=\ell B_0+R(q),\quad
 v=(\log|w_1|,\log|w_2|)^t.
\]
The exact coordinate identity is
$v=\ell B_0u+R(q)(u-u_0)$, and on (B.3.2) it gives
\[
 du=A(q_0)^{-1}(dv-p_0\,d\ell).
 \tag{B.3.3}
\]
Thus no derivative of a bounded correction has been
ignored.  Let $L_j$ and $Q_j$ denote the real tangent
vectors $w_j\partial_{w_j}$ and $iw_j\partial_{w_j}$
in the real identification $T\mathbb C^3=\mathbb C^3$.
With $A_0=A(q_0)$, dual lifts of the target coordinate
order $(\ell,\phi,u_1,u_2)$ are
\[
 N_\ell=\tfrac12L_0+\tfrac14L_1+\tfrac14L_2,
 \qquad N_\phi=Q_0,
\]
\[
 N_i=-((A_0)_{1i}+(A_0)_{2i})L_0
            +(A_0)_{1i}L_1+(A_0)_{2i}L_2,
 \quad i=1,2.
 \tag{B.3.4}
\]
The tangent vectors are $T_a=-Q_0+Q_1$ and
$T_b=-Q_0+Q_2$.  Formula (B.3.3) checks that the
four normal vectors map to the four standard coordinate
vectors and both tangent vectors map to zero.
The full determinant is
$-\det(A_0)|q_0|^2<0$, so the positive tangent order is
$(-T_a,T_b)$.  The inequality holds for small $r_0$
because $A_0$ is joined to $\ell_0B_0$ through
$\ell_0B_0+sR(q_0)$, all of which are invertible.

Every displayed vector is a constant linear combination
of $L_j,Q_j$.  Its dependence on $(a,b)$ is therefore
exactly
\[
 D(a,b)=\operatorname{diag}(e^{-i(a+b)},e^{ia},e^{ib})
            \in\mathrm{SU}(3).
\]
This map is null-homotopic.  An explicit periodic
contraction uses
\[
 Q_s(\theta)=
 \begin{pmatrix}
 \sin(\pi s/2)e^{-i\theta}&-\cos(\pi s/2)\\
 \cos(\pi s/2)&\sin(\pi s/2)e^{i\theta}
 \end{pmatrix},
 \qquad K_s(\theta)=Q_s(\theta)Q_s(0)^*.
\]
Each matrix lies in $\mathrm{SU}(2)$,
$K_0=I$, and $K_1(\theta)=\operatorname{diag}
(e^{-i\theta},e^{i\theta})$.  The product of its
embeddings in coordinate planes $(0,1)$ and $(0,2)$,
evaluated at $a$ and $b$, contracts $D$.
Consequently the induced stable tangential frame of
(B.3.2) is the constant Lie frame.  The ambient framing
comparison on the containing ball is null-homotopic,
so the same conclusion holds in the framed sphere.

For the induced spin structure the two coordinate
circles are nonbounding.  Its quadratic form on
$H_1(T^2;\mathbb F_2)$ is
$q(xa+yb)=x+xy+y$, and its Arf invariant is
$q(a)q(b)=1$.  Hence its framed class is
$\eta^2\in\pi_2^{\mathrm S}$.  Freudenthal suspension
is an isomorphism for $\pi_6(S^4)$, so the unstable
class is exactly (B.3.1).
\end{proof}

\paragraph{Theorem B.3.2 (Proved).}
The angle-forgetting workbench maps satisfy
$[\widehat q_3]=0$ and $[\widehat q_\gamma]=0$.

\begin{proof}
Both components of $\overline\pi$ extend as integral
$H^1$ classes over $N_0$, by (A.1.5).  Choose a map
$\pi':N_0\to T^2$ representing them and equal to
$\overline\pi$ on an outer collar; the construction
and relative uniqueness are as in A.2.1.  On the
punctured cusp the quotient $\overline\pi-\pi'$ has
a lift to $\mathbb R^2$ that is zero on the collar.
Indeed the punctured cusp retracts onto its boundary
mapping torus, and the difference is zero in relative
$H^1$.  Multiplying that lift by a homotopy parameter
gives a homotopy from $\overline\pi$ to $\pi'$.
After passage to either suspension target this
homotopy is continuous at $W$, uniformly in the
homotopy parameter: the whole compact torus of
angles is collapsed as $r\to0$.  It is also constant
on the outer boundary and extends by the basepoint
outside $N_0$.

After this homotopy the three-sphere map factors as
\[
 X\longrightarrow
 \frac{[0,\epsilon]\times T^2}{\{\epsilon\}\times T^2}
 \longrightarrow\Sigma(T^2/T^{2,(1)}),
 \qquad x\longmapsto[r(x),\pi'(x)]
\]
on $N_0$, with the complement mapping to the vertex.
The middle space is a cone and is contractible.
For the two-sphere map use the first component
$\theta':N_0\to S^1$ and the cone on $S^1$.
These give actual null-homotopies of the two maps.
This argument cannot discard the retained base
angle in $H$, whose intermediate space is instead
$(D^2\times T^2)/(\partial D^2\times T^2)$.
\end{proof}

\paragraph{Theorem B.3.3 (Proved).}
Let $C=M\cup\operatorname{Cone}(\partial M)$ and let
$q_1:X\to C$ be the workbench map retaining the
Seifert quotient away from the cusp and collapsing
$W$ to the cone point.  Under the cofiber splitting
$C\simeq S^3\vee S^2$ determined by the cusp core,
\[
 [q_1]=(j_3)_*(6\nu'),
 \tag{B.3.5}
\]
with no $S^2$ or mixed Whitehead component.

\begin{proof}
Write $Q=M\cup(D^2\times S^1)$ and let $K$ be the
core $\{0\}\times S^1$ of its cusp solid torus.
Collapsing this core turns the solid torus into
the cone on its boundary, by its radial coordinate.
Thus $C=Q/K$ as a quotient space.  The same
homotopy used in B.3.2, now retaining the base angle,
shows $q_1\simeq\kappa e$, where $\kappa:Q\to Q/K$
is the quotient map.  The inclusion $K=S^1\to S^3$
is null-homotopic, and its cofiber is
$S^3\vee\Sigma S^1$.  We choose this equivalence
to restrict on $Q$ to $j_3$.  A null-homotopy of
the circle exists since $\pi_1(S^3)=0$; its
ambiguity is $\pi_2(S^3)=0$.  Hence the indicated
splitting is fixed in the only respect used here,
and B.1.2 proves (B.3.5).
\end{proof}

\subsection{B.4. The other named sphere fibrations}

\paragraph{Theorem B.4.1 (Proved).}
For the sphere members of Viaclovsky's two-parameter family,
the Liu--Rollenske family, and every $X_{3k,4l}$ with
$\gcd(3k,4l)=1$, the primitive invariant-circle quotient
admits a canonical replacement over the degeneration
neighbourhoods as in A.2.1.  Its target is $S^3$ and
its class is $6\nu'$.  In each case the map to
$\mathbb P^1$ is null-homotopic.  The retained-angle
map at the primitive rank-two cusp is $\eta_4\eta_5$,
and the two corresponding angle-forgetting maps are
null-homotopic.  Thus the quotient class $[e]$ does
not distinguish these fibrations.

\begin{proof}
We supply the integral lattices, the local extension
comparisons, and the target calculation; a list of
Betti numbers would not imply this statement.

For $X_{3k,4l}$ the regular homological representation
is still $(A_1,A_2,M_0)$ of A.1.1.  We use K1 with
its original finite-root convention.  In Viaclovsky's
marking $(\delta,a,b,d)$ the comparison is
\[
 P_V=\begin{pmatrix}-2&1&0&1\\-6&0&-1&0\\0&-1&0&0\\-1&0&0&0\end{pmatrix}.
\]
Substitution in the matrices of
\cite[Lemma 2.1 and equations (2.11)--(2.12)]{mapsViaLog}
gives $V_iP_V=P_VA_{i+1}$, $\det P_V=-1$,
$P_Vd=\delta$, $P_Vv_1=-u_0$, and
$P_V(-v_2)=\delta-u_1$.  This checks the step of
K1 that we use, without redoing its analytic
classification.  With $m=3k,n=4l$, the actual
finite translations are the vectors in K1 and
the circle transgression is
$b=nN_0+mN_1=\pm1$ for a sphere member.
The two target fillings have multiplicities $m,n$;
their genus-one meridian determinant is $b$.
Thus their target is $S^3$.  Its map to the base
is a weighted Hopf map of degree product $mn=12kl$.
The common lattices and their mod-two counterparts
are those in B.1.1, and all the arguments of
A.2.1 and B.1.2 apply without change.

For the two-parameter family use the homology basis
$(\widehat\psi,\widehat u,\widehat w,\widehat\delta)$
of K11.  Its actual four matrices are
\[
 H_* =\begin{pmatrix}1&0&0&0\\-2&0&-1&0\\0&1&0&0\\1&1&0&1\end{pmatrix},
 \qquad
 H_1 =\begin{pmatrix}1&0&0&0\\0&1&0&0\\-2&-1&1&0\\-2&-1&0&1\end{pmatrix},
\]
\[
 H_2 =\begin{pmatrix}1&0&0&0\\0&1&1&0\\0&0&1&0\\0&0&0&1\end{pmatrix},
 \qquad
 H_3 =\begin{pmatrix}1&0&0&0\\0&1&0&0\\0&-1&1&0\\1&0&0&1\end{pmatrix}.
\]
They are inverse transposes of the matrices in
\cite[Proposition 4.2]{mapsVia2}.  The direct
calculation gives
\[
 \Lambda^\Gamma=\mathbb Z\widehat\delta,
 \quad (\Lambda^*)^\Gamma=\mathbb Z\psi,
 \quad (\Lambda^*\otimes\mathbb F_2)^\Gamma
          =\mathbb F_2\psi,
 \quad L=\ker\psi.
 \tag{B.4.1}
\]
The finite root action has order four and translation
$\ell/4$, where $\ell=(1,-1,-1,0)^t$; hence
$\psi(\ell)=1$.  The three primitive vanishing
lattices are
\[
 K_1=\mathbb Z(\widehat w+\widehat\delta),\qquad
 K_2=\mathbb Z\widehat u,\qquad
 K_3=\mathbb Z\langle\widehat w,\widehat\delta\rangle.
 \tag{B.4.2}
\]
They all lie in $L$.  To check that the needed
class really extends, the two rank-one local
models in the source are the quotient of an
infinite chain of ruled elliptic surfaces, with
the adjacent sections identified by a translation.
Their nearby fundamental group surjects to the
central one with precisely $K_i$ killed; the
base meridian bounds the extending section.
This follows directly by filling the one
vanishing circle with its nodal thimble, and
retaining the other three circle generators.
It is also the degree-one part of the source's
local-specialization proposition, labelled
\texttt{mum}.  At the rank-two cusp the same
assertion is (A.1.5).  Thus $\psi$ extends
over every reduced filling, and its boundary
restriction is injective on $H^1$.  A.2.1
constructs and compares the replacements.

There is a single multiple target circle of
multiplicity four.  The other three circle
fillings are ordinary solid tori, with no
meridian shift because the chosen section
extends there.  The finite logarithmic
shift has primitive circle component $1/4$.
Equivalently its integral transgression is
$d_2(4\psi)=1$, as computed directly in
\cite[Proposition 5.4]{mapsVia2}.
Hence the genus-one target has meridian
determinant one and is $S^3$.

For Liu--Rollenske use their common homology
basis $(e_1,e_2,e_3,e_4)$.  The two finite
monodromies are
\[
 R_0=\begin{pmatrix}-1&-1&-1&0\\2&1&1&0\\0&0&1&0\\0&0&0&1\end{pmatrix},
 \qquad
 R_1=\begin{pmatrix}1&1&0&1\\-2&0&-1&-1\\0&-1&1&-1\\0&0&0&1\end{pmatrix}.
\]
Their source uses $R_\infty=(R_1R_0)^{-1}$.
They have respective orders four and six,
but the fibre over $1$ is reduced: finite
monodromy does not mean it is a smooth
multiple fibre.  The common invariants are
\[
 \Lambda^\Gamma=\mathbb Z(-e_1+2e_3),
 \quad(\Lambda^*)^\Gamma=\mathbb Z e_4^*,
 \quad(\Lambda^*\otimes\mathbb F_2)^\Gamma
            =\mathbb F_2 e_4^*.
 \tag{B.4.3}
\]
Here $\delta=-e_1+2e_3$ and $\xi=e_4^*$,
so again the fixed vector belongs to
$L=\ker\xi$.
The local product identification of
\cite[Lemma 3.10 and Proposition 3.11]{mapsLR},
$Y_{\Delta_1}\cong S_{\Delta_1}\times T^2$,
preserves the two circle classes $f_3,f_4$,
with $\xi=f_4^*$.  Consequently $\xi$
extends over this reduced type-II fibre
and is surjected onto by the nearby
fundamental group.  At infinity the
primitive vanishing lattice is
$\mathbb Z\langle h_1,h_3\rangle$ and
$\xi=h_4^*$, so (A.1.5) gives its
extension there.  The only logarithmic
shift is $a=e_4/4$ over $0$.
Its primitive transgression is
$d_2(4\xi)=1$ by \cite[Proposition 4.8]{mapsLR}.
Thus the same one-exceptional-circle
target is $S^3$.  These are applications
of the exact local product and integral
monodromy statements in
\cite[Propositions 3.12--3.13 and Corollary 3.16]{mapsLR},
with the comparison of classes written
explicitly; rational invariant cycles
alone would not suffice.

For both one-multiple-fibre families,
Lemma B.1.1 applies verbatim with
$(\gamma,d)$ replaced by $(\psi,\widehat\delta)$
or $(\xi,\delta)$.  Equations (B.4.1)
and (B.4.3) are all of its hypotheses.
The chosen regular inverse image of
the quotient map is a single $T_L$,
because every replacement is over its
own degeneration disc.  Therefore
Theorem B.1.2 gives $[e]=6\nu'$ in
each family.  The target projection
is the weighted Hopf map of weights
$(4,1)$, whose power-map factor has
degree four; it sends $6\nu'$ to zero.
For $X_{3k,4l}$ the degree is $12kl$,
which also sends it to zero.

Finally each named family has a
primitive hexagonal rank-two cusp.
The proof of B.3.1 is local in its
normal-crossing chart, independent
of the other fibres and of the
bounded period matrix.  It gives
$H=\eta_4\eta_5$ with the same target
coordinate convention.  The two
long-period classes extend over
that cusp, so B.3.2 proves the
two angle-forgetting assertions.
All kernel and mod-two calculations
above are executed in
\texttt{code/maps\_exact.py}.
\end{proof}

\subsection{B.5. A cell presentation from the three fillings}

We give a topological cell construction, retaining the actual
boundary maps before making cellular approximations.  The
smaller cellular complex used for the integer calculation is
homotopy equivalent to this construction; its cell counts are
not asserted to be the counts of a smooth handle decomposition.
Throughout this subsection $p=-1$, so the finite twists are
$v_1=g+2u-4w$ and $v_2=-g-3u+3w$ and the cusp gluing has
zero translation drift.

For a matrix $A$ on $F=\mathbb R^4/\Lambda$, our mapping-torus
convention is
\[
 M_A=(F\times[0,1])/((x,1)\sim(A^{-1}x,0)).
 \tag{B.5.1}
\]
The positive interval loop $\rho$ satisfies
$\rho t_\lambda\rho^{-1}=t_{A\lambda}$.  This fixes the
inverse in (B.5.1) and agrees with the clockwise meridians
used in A.1.

\paragraph{Theorem B.5.1 (Proved).}
The underlying topological manifold of $X$ is obtained from
the flat $T^4$-bundle over a ribbon bouquet of two circles,
with monodromies $A_1,A_2$, and the three spaces $S_1,S_2,W$
by the three double mapping cylinders specified below.
Their maps admit finite polyhedral cell structures.  They
therefore give a CW structure adapted to the pair of pants
and the three special fibres, with all attaching maps
specified by integral affine formulae and primitive circle
quotients.  Replacing the maps by cellular maps on the smaller
torus cell structures gives a finite CW homotopy model with
cell counts
\[
 (c_0,c_1,c_2,c_3,c_4,c_5,c_6)
       =(5,20,45,56,42,17,3).
 \tag{B.5.2}
\]

\begin{proof}
Give a bouquet with oriented loops $\rho_1,\rho_2$ the ribbon
order in which the two ends of each loop are consecutive.
Its oriented thickening is a pair of pants.  Collapse the
widths of its two bands to the edges and the remaining
vertex disc to the vertex.  The induced boundary maps are
piecewise linear paths with words
$\rho_1,\rho_2,(\rho_1\rho_2)^{-1}$; intervals between
successive letters are mapped to the vertex.  The thickening
is the mapping cylinder of this boundary map.  This assertion
can be checked on each band, where the cylinder consists of
the two half-bands, and on the vertex disc, where the cones
on the intervening boundary intervals are its sectors.
Lift this construction to the flat $F$-bundle.  Write $U_*$
for the resulting bundle over the bouquet.  Its boundary
spaces are $M_{A_1},M_{A_2},M_{M_0}$.  Their maps
$b_j:M_{A_j}\to U_*$ and $b_0:M_{M_0}\to U_*$ are the
lifted ribbon boundary maps.  They are the identity on a
marked fibre and send their interval loops respectively to
$\rho_1,\rho_2,\rho_2^{-1}\rho_1^{-1}$.
The equality $M_0=(A_1A_2)^{-1}$ verifies the last map
including its fibre monodromy.

At a finite fibre put
\[
 h_j(x)=A_jx+v_j/m_j,\qquad S_j=F/\langle h_j\rangle,
 \qquad
 r_j([x,\theta])=[x+\theta v_j/m_j].
 \tag{B.5.3}
\]
The last formula is a well-defined map $M_{A_j}\to S_j$:
at $\theta=1$ its value agrees with $[A_j^{-1}x]$, because
$h_j(A_j^{-1}x)=x+v_j/m_j$.  The finite neighbourhood is
the associated disc bundle
$(F\times D^2)/\langle(x,s)\mapsto(h_jx,e^{-2\pi i/m_j}s)\rangle$.
Its boundary point corresponding to $[x,\theta]$ is
\[
 [x+\theta v_j/m_j,\ e^{-2\pi i\theta/m_j}].
\]
Radial collapse is exactly (B.5.3); its mapping cylinder is
the disc bundle.  In particular this is a statement about
the underlying smooth neighbourhood and does not replace
the scheme-theoretic fibre $m_jS_j$ by a reduced fibre in
the holomorphic map.  The logarithmic gluing translation is
$\log(s)\Pi v_j/(2\pi i)$.  Its clockwise drift is precisely
$v_j/m_j$.  Subtracting that displayed drift leaves a
periodic translation with a lift to $\mathbb R^4$; linear
interpolation of this lift straightens it on a collar.
Thus the actual gluing is identified with (B.5.3), without
discarding a loop of translations of nonzero lattice drift.
This also corrects the first paragraph on page 42 of
\cite{mapsAC26}: multiplicity does not prevent this
neighbourhood from being a smooth normal disc bundle.
We make no assertion of a holomorphic tubular projection.
The meridian relation proved in that paragraph remains
valid; its meridian is the normal circle winding $m_j$
times around the base, rather than a once-around base
loop.

Here is the complete cusp boundary map.  Put
$a=(a_1,a_2)\in\mathbb R^2/\mathbb Z^2$ for the $(g,u)$
coordinates and $b\in T_f=\mathbb R^2/\mathbb Z^2$ for
the $(w,d)$ coordinates.  Let $\mathcal H$ be the hexagonal
honeycomb dual to the $A_2$ triangulation.  For $B_0a$ in
the interior of the hexagon centred at $v\in\mathbb Z^2$,
write $v(a)=v$.  The central fibre has the explicit quotient
description
\[
 W_{\mathcal H}=(T_B\times T_f)/\sim,
 \qquad T_B=\mathbb R^2/\mathbb Z^2,
 \tag{B.5.4}
\]
where over a honeycomb edge separating centres $v,v'$ one
collapses the primitive circle on $v'-v$, and over each
honeycomb vertex one collapses all of $T_f$.  No circle is
collapsed over a hexagon interior.  The three oriented edge
directions can be chosen as
\[
 n_1=w,\qquad n_2=d-w,\qquad n_3=-d.
 \tag{B.5.5}
\]
The quotient is homeomorphic to the actual $W$: the positive
part of each toric component is its hexagon, the phase
stabilizers are exactly these primitive circles, and the
period action identifies hexagons by $B_0\mathbb Z^2$.
The constant phase cocycle is removed by the linear phase
gauge coming from $\operatorname{Re}C(0)a$.

With the interval convention (B.5.1), the cusp attachment is
\[
 r_0([a,b,\theta])
   =[a,b-\theta(B_0a-v(a))]\ \in W_{\mathcal H}.
 \tag{B.5.6}
\]
This formula includes its seam data.  On crossing an edge
the two expressions differ by $\theta(v-v')$, which lies
in the circle collapsed on that edge.  At a vertex every
phase is collapsed.  Replacing $a$ by $a+\lambda$ changes
$v(a)$ by $B_0\lambda$, so the expression is periodic.
At $\theta=1$, the integer vector $v(a)$ disappears modulo
$\mathbb Z^2$, leaving $[a,b-B_0a]$; this is the endpoint
identification by $M_0^{-1}$.  Hence (B.5.6) is continuous
on the entire mapping torus, not only on its open strata.

We justify that the cusp neighbourhood is its actual
mapping cylinder.  Straighten $C(q)$ to $C_0=C(0)$ by
the map on logarithmic coordinates
\[
 \xi\longmapsto
 \xi-(C(q)-C_0)(\operatorname{Im}Z_q)^{-1}
                       \operatorname{Im}\xi.
 \tag{B.5.7}
\]
Writing $\xi=Z_qa+b$ proves that this sends $Z_qa+b$ to
$Z_q^0a+b$, fixes the real lattice, and conjugates the
period actions.  In a bounded-position toric chart the
correction multiplying $x$ tends to one as $q\to0$;
the same holds for its inverse.  Thus it extends by the
identity on $W$ to a homeomorphism of the closed cusp tube.

Let $A$ be the nonnegative real part of the frozen toric
cover over $|q|\leq\eta$, and $A_0$ its central part.
The local orthant descriptions make $A$ a three-manifold
with boundary; the corners may be rounded.  Its interior
has coordinates $(y,\rho)\in\mathbb R^2\times(0,\eta)$,
$y=\log|x|/\log\rho$.  The positive period action is
$y\mapsto y+B_\rho^+\lambda$, where
$B_\rho^+=B_0-2\pi\operatorname{Im}C_0/\log\rho$.
For small $\eta$ these matrices are invertible.  Thus the
interior of the positive quotient $Q^+$ is explicitly
$T^2\times(0,\eta)$, using $a=(B_\rho^+)^{-1}y$.
The bounded fundamental parallelogram meets only finitely
many of the toric bounded-position charts, including at
$\rho=0$, so $Q^+$ is compact.  Its two boundary components
are tori.  The inner one is the honeycomb quotient, and
the outer one is the positive regular-fibre torus.

Both boundary inclusions induce the identity on the period
deck group $\mathbb Z^2$.  Indeed $A_0$ is the honeycomb
plane and hence contractible; $A$ is contractible because
its interior is contractible and pushing a boundary collar
inward gives a homotopy equivalence with that interior.
The quotient is irreducible.  An embedded locally flat
two-sphere lies in its interior $T^2\times(0,\eta)$.
Lift it to $\mathbb R^2\times\mathbb R$.  Its bounded
ball is disjoint from all its nonzero lattice translates:
disjoint boundary spheres would otherwise give nested
balls, impossible for translates of equal finite volume.
The ball therefore descends.  This also proves
irreducibility of the product target $T^2\times I$.
The boundary tori are incompressible because their
fundamental groups inject.

There is a map of pairs $Q^+\to T^2\times I$ inducing
the identity on $\pi_1$ and restricting to the prescribed
homeomorphisms on the two boundaries.  For its first
coordinate, extend the classifying map to $T^2=K(\mathbb
Z^2,1)$ and use the homotopy extension property on the
boundary; use a collar height for the second coordinate.
Apply \cite[Theorem 6.1, p.~77]{mapsWaldhausen68} to this
map.  The source is sufficiently large, since a boundary
torus pushed into its interior is incompressible.  The
product target has the same property, and all the other
hypotheses were just verified.  The homotopy is fixed on
the boundary.  Its alternative of mapping the connected
source into the boundary is impossible, since the two
source boundaries map to distinct target components.
It therefore ends at a covering map.  The induced
fundamental-group isomorphism makes that covering
one-sheeted.  We obtain a product homeomorphism relative
to both prescribed boundary markings.

Lift this product to $A$ and spread it by compact phases.
The full period action on a phase $(\phi,\vartheta)$ is
$(\phi,\vartheta)\mapsto
(u_\lambda\phi\vartheta^{B_0\lambda},\vartheta)$;
it does not depend on the positive radial point.  The
lifted product is therefore equivariant for the full
action.  Phase stabilizers occur only on $A_0$ and are
exactly those of (B.5.4).  Consequently the resulting
quotient is the mapping cylinder of the boundary phase
map.  On the component with centre $v$, the $q$-phase
acts on $T_f$ by $-v$ times that phase, since $(v,1)$ is
the stabilizer ray.  Substituting the clockwise interval
phase and the marked period $-\theta B_0a$ gives precisely
(B.5.6).

Gluing the three mapping cylinders of $r_j$ to the
mapping cylinder of the ribbon boundary maps gives
\[
 U_*\ \cup\ \coprod_{j=0}^2(M_j\times[0,1])
       \ \cup\ (W\sqcup S_1\sqcup S_2),
 \tag{B.5.8}
\]
with $(x,0)\sim b_j(x)$ and $(x,1)\sim r_j(x)$.
The preceding constructions identify (B.5.8)
homeomorphically with $X$.  The comparison need not
preserve the radial value of $f$ inside the cusp, but
it preserves the special strata and all three outer
collars; this is the sense of adaptation used here.

For a finite cell construction use rational hexagons,
unit cubes for tori, and the affine quotient (B.5.3).
For (B.5.6) use toric phase coordinates
$b'=b-\theta B_0a$.  In these coordinates its expression
on a hexagon is $[a,b'+\theta v]$, and its face
identifications are affine:
$a\mapsto a+\lambda$ changes $b'$ by
$-\theta B_0\lambda$.  There are finitely many faces
in the fundamental polyhedra.  Cut their images at the
integer phase hyperplanes and subdivide; this makes
every face map a finite piecewise affine map.  Primitive
edge-circle quotients are the integral linear maps
$b\mapsto\det(n_k,b)\bmod\mathbb Z$.
Triangulate each occurrence of a torus separately and
interpolate through a common subdivision along each
mapping cylinder.  This does not require a finite
triangulation invariant under the infinite monodromy
group.  Products of simplices with intervals are
triangulated by their ordered vertex prisms.  These
rules supply finite characteristic maps for (B.5.8),
and pulling the cells back gives the asserted CW
structure of $X$.

For the smaller homotopy model put the usual product
cell structure on each torus and apply cellular
approximation to the displayed maps.  This changes
the double mapping cylinder by homotopy equivalence;
equivalently use \cite[Propositions 4G.1--4G.2]{mapsHatcher02}.
The two finite surfaces are mapping tori of $T_L$,
because $m_j\gamma$ is a primitive circle coordinate
on $S_j$; their linear monodromies are
$B_j=A_j|_L$ (an inverse convention gives the same
counts).  Their cell counts are $(1,4,6,4,1)$.
The counts for $U_*$ are $(1,6,14,16,9,2)$, and
each $M_j$ has counts $(1,5,10,10,5,1)$.
In (B.5.4) the quotient honeycomb has two vertices,
three edges, and one hexagon.  The surviving phase
spaces are respectively a point, a circle, and a
two-torus, giving counts $(2,3,4,2,1)$ for $W$.
Adding these cells and one prism cell for each
boundary cell gives (B.5.2).
\end{proof}

\paragraph{Theorem B.5.2 (Proved).}
The attaching data of Theorem B.5.1 give
$\pi_1(X)=1$ and integral homology
$H_k(X;\mathbb Z)=\mathbb Z$ for $k=0,6$ and zero
otherwise.  They give an explicit recognition of $X$
as a homotopy sphere, and exhibit the normally framed
three-torus used in B.1.

\begin{proof}
The graph piece has presentation
\[
 \pi_1(U_*)=
 \langle\Lambda,\rho_1,\rho_2\mid
 [\Lambda,\Lambda]=1,\quad
 \rho_jt_\lambda\rho_j^{-1}=t_{A_j\lambda}\rangle.
\]
Formula (B.5.3) imposes $\rho_j^{m_j}=t_{v_j}$.
Formula (B.5.6) kills $w,d$ and the cusp meridian
$a_0=\rho_2^{-1}\rho_1^{-1}$.  Since $A_1w=u-w$,
the normal closure also kills $u$.  Writing $c$ for
the image of $g$, the remaining presentation is
\[
 \langle c,x,y\mid c\text{ central},\quad
 xy=1,\quad x^3=c,\quad y^4=c^{-1}\rangle.
 \tag{B.5.9}
\]
Substituting $y=x^{-1}$ gives $x^3=x^4$, so $x=c=y=1$.
This calculates the nonabelian fundamental group,
not only its abelianization.

We record the integral homology maps in primitive
bases.  Write $\gamma,u^*,w^*,d^*$ for the cohomology
basis dual to $g,u,w,d$, and put
\[
 \eta_1=2u^*+w^*+3d^*,\quad
 \eta_2=u^*+w^*+2d^*,\quad
 q_2=u^*w^*+6\gamma d^*.
 \tag{B.5.10}
\]
Products here are exterior products.  First attach
the two finite pieces over a disc containing $0,1$;
call the result $X^\circ$.  Its homotopy type is the
double mapping cylinder $S_1\leftarrow F\to S_2$.
The degree-one
map $(\pi_{1*},-\pi_{2*})$ is
\[
 \alpha_1=
 \begin{pmatrix}
 3&0&0&0\\0&2&1&3\\4&0&0&0\\0&-1&-1&-2
 \end{pmatrix}.
 \tag{B.5.11}
\]
To verify its primitive target bases, abelianize
$\pi_1(S_j)=\langle\Lambda,\rho_j\mid
\rho_j\lambda\rho_j^{-1}=A_j\lambda,
\rho_j^{m_j}=t_{v_j}\rangle$.
For $S_1$ the images are
$g=3\rho_1,u=2w,d=3w$; for $S_2$ they are
$g=-4\rho_2,u=w,d=2w$.  Thus
$(\rho_1,w)$ and $(\rho_2,w)$ are free integral bases.
This proves (B.5.11) without a rational invariant-lattice
identification.  Its Smith form is $(1,1,1,0)$,
its kernel is generated by
$\nu_0=-u-w+d$, and its primitive left kernel is
$(4,0,-3,0)$.  Consequently
$H_1(X^\circ)=\mathbb Z\zeta$, normalized by
\[
 g\mapsto12\zeta,\quad
 \rho_1\mapsto4\zeta,\quad
 \rho_2\mapsto-3\zeta,\quad
 a_0\mapsto-\zeta.
 \tag{B.5.12}
\]

For completeness we derive the finite degree-two
pullback lattices.  The mapping-torus matrices on
$L$ are
\[
 B_1=\begin{pmatrix}0&1&0\\-1&-1&0\\1&0&1\end{pmatrix},
 \qquad
 B_2=\begin{pmatrix}0&-1&0\\1&0&0\\0&1&1\end{pmatrix}.
\]
The exterior-power mapping-cone complexes show that
$H_*(S_j)$ is free, with ranks $(1,2,2,2,1)$:
the nonzero Smith entries of
$\bigwedge^r B_j-I$ are all one.  In particular the
intersection lattice in degree two is unimodular.
Direct exterior multiplication gives the primitive
invariant lattices
$I_1=\langle\gamma\eta_1,q_2\rangle$ and
$I_2=\langle\gamma\eta_2,q_2\rangle$.
Their intersection determinants have absolute values
$9$ and $4$.  The pullback under a covering of degree
$m_j$ multiplies the intersection pairing by $m_j$.
Its image therefore has determinant $m_j^2$ and index
one in $I_1$, respectively index two in $I_2$.
The embedded torus generated by $v_2,d$ is invariant
under $h_2$ and covers its quotient with degree four.
The class $a\gamma\eta_2+bq_2$ evaluates on it as
$-2a-6b$.  A pullback class must therefore satisfy
$a+b\equiv0\pmod2$.  This is already an index-two
condition and hence determines the image exactly.
Thus
\[
 \pi_1^*H^2(S_1)=\langle\gamma\eta_1,q_2\rangle,
 \qquad
 \pi_2^*H^2(S_2)=\langle2\gamma\eta_2,
                                  \gamma\eta_2+q_2\rangle.
 \tag{B.5.13}
\]

Use bases with pullbacks
$(\gamma\eta_1,q_2-2\gamma\eta_1)$ and
$(2\gamma\eta_2,q_2-3\gamma\eta_2)$, and their
Kronecker-dual homology bases.  Evaluating on
$(gu,gw,gd,uw,ud,wd)$ gives
\[
 \alpha_2=
 \begin{pmatrix}
 2&1&3&0&0&0\\-4&-2&0&1&0&0\\
 -2&-2&-4&0&0&0\\3&3&0&-1&0&0
 \end{pmatrix}.
 \tag{B.5.14}
\]
Its Smith form is again $(1,1,1,0)$ and its primitive
left kernel is $\Phi=(4,2,3,2)$.  Thus its cokernel is
$\mathbb Z e$, where the cohomology class $\mathcal A$
given by $\Phi$ satisfies
$\langle\mathcal A,e\rangle=1$ and restricts on $F$
to $2q_2$.  The exact sequence of this double mapping
cylinder gives
\[
 0\longrightarrow\mathbb Z e\longrightarrow
 H_2(X^\circ)\longrightarrow\mathbb Z\nu_0
 \longrightarrow0.
 \tag{B.5.15}
\]

The central quotient (B.5.4) has cellular boundary
$\partial_1(e_k)=Q-P$ and all higher cellular
boundaries zero.  Indeed the hexagon attaching word
can be chosen as
$e_1e_2^{-1}e_3e_1^{-1}e_2e_3^{-1}$; every edge occurs
with both signs.  Crossing with a surviving phase
circle preserves this cancellation, and the torus
cell boundary is the commutator.  Hence
\[
 H_*(W)=(\mathbb Z,\mathbb Z^2,\mathbb Z^4,
                       \mathbb Z^2,\mathbb Z).
\]
At $\theta=0$ the cusp map is the quotient
$c(a,b)=[a,b]$.  On $H_1$ it is projection to
$\langle g,u\rangle$.  With degree-two target basis
$(T_B,\overline C_1,\overline C_2,\overline C_3)$,
its matrix is
\[
 c_2=\begin{pmatrix}
 1&0&0&0&0&0\\0&0&-1&0&0&0\\
 0&-1&-1&1&1&0\\0&0&0&1&0&0
 \end{pmatrix}.
 \tag{B.5.16}
\]
For a mixed torus $\beta v$ this follows from the
degree formula
$\sum_k\det(n_k,B_0\beta)\det(n_k,v)[\overline C_k]$.
The first factor is its oriented honeycomb edge
coefficient, and the second is the degree of the
primitive circle quotient.  The map sends $gu$ to
$T_B$ and $wd$ to zero.  In degree three it sends
$guw,gud$ to a basis and kills $gwd,uwd$; in degree
four it has degree one on the open hexagon times
$T_f$.  All these maps are integrally surjective.
Their kernels are exactly
$\operatorname{im}(\bigwedge^r M_0-I)$: inclusion
follows by transport around the boundary, and the
displayed kernels and monodromy images have the
same primitive bases, as checked in
\texttt{code/maps\_cw.py}.

Let $K_i$ be the kernel of $H_i(M_0)\to H_i(W)$.
The mapping-torus sequence and the preceding
surjections give
\[
 K_1=\mathbb Z a_0,\qquad
 K_2=\mathbb Z\xi_w\oplus\mathbb Z\xi_d,
 \qquad \xi_v=[S^1_v\times S^1_{a_0}].
 \tag{B.5.17}
\]
These are literal tori because $M_0$ fixes $w,d$.
Under the mapping-torus projection they go to $w,d$,
so their span is primitive.  In (B.5.6), over a
hexagon centre the $a_0$ coordinate is collapsed;
thus their images in $H_2(W)$ vanish.  The ranks
from the mapping-torus sequence show that this is
the whole kernel.

Orient the connecting map in (B.5.15) so that
\[
 \partial j_*\xi_v=(I-A_2)v
                   =(I-A_1^{-1})v\quad(v\in\langle w,d\rangle).
\]
One verifies this on the two rectangles obtained
by cutting the swept torus along the finite band:
their unmatched boundary circles are $v$ and
$-A_2v$.  Since $(I-A_2)w=-\nu_0$ and
$(I-A_2)d=0$, the class $z=j_*\xi_w$ together with
$e$ is a basis of $H_2(X^\circ)$.
To determine $j_*\xi_d$, use the circle translation
in direction $d$ on $X^\circ$.  For this action the
integral sweep operator on cohomology restricts on
$F$ to contraction by $d$.  The class
$\mathcal A|_F=2q_2$ satisfies
\[
 \iota_d(2q_2)=-12\gamma.
\]
Restriction $H^1(X^\circ)\to H^1(F)$ is injective
by (B.5.12), with image $\mathbb Z\,12\gamma$.
Thus $\iota_d\mathcal A=-\Theta$, where
$\Theta(\zeta)=1$.  Evaluating the swept torus and
using $\Theta(a_0)=-1=p$ gives
\[
 j_*\xi_d=-p e=e,
 \qquad
 [j_*|_{K_2}]_{(\xi_d,\xi_w)}^{(e,z)}
       =\begin{pmatrix}1&0\\0&1\end{pmatrix}.
 \tag{B.5.18}
\]
This computes the attaching coefficient that cannot
be recovered from ranks alone.

In degree one the full cusp attachment map in bases
$(g,u,a_0)$ and $(\bar g,\bar u,\zeta)$ is
\[
 D_1=\begin{pmatrix}1&0&0\\0&1&0\\-12&0&1\end{pmatrix}.
\]
In degree two, split the surjection onto $H_2(W)$;
the remaining block is exactly (B.5.18).  Thus the
full map is integrally equivalent to $I_4\oplus I_2$.
Both are unimodular.  Mayer--Vietoris proves
$H_1(X)=H_2(X)=0$, with no torsion.  The cell count
(B.5.2) gives $\chi(X)=2$.  Since $X$ is a closed
oriented six-manifold, duality and universal
coefficients give $H_4=H_5=0$ and make $H_3$ free:
$\operatorname{Ext}(H_3,\mathbb Z)$ injects into
$H^4(X)=H_2(X)=0$.  The Euler characteristic then
gives $\operatorname{rank}H_3=0$.  This proves all
the asserted integral groups.  The program checks
every matrix, Smith form, primitive coefficient
and cell count just used, using exact arithmetic.

Starting from (B.5.9), successive applications of
the Hurewicz theorem make $\pi_i(X)=0$ for
$1\leq i\leq5$.  The same theorem produces a map
$S^6\to X$ carrying the fundamental class to the
chosen generator.  The mapping-cylinder pair of
this map is simply connected and has zero relative
homology.  If it had a first nonzero relative
homotopy group, the relative Hurewicz theorem
would make its relative homology nonzero.  All
relative homotopy groups therefore vanish, and
the CW Whitehead theorem makes the map a homotopy
equivalence.  These
are \cite[Theorems 4.32 and 4.5]{mapsHatcher02}
applied to the explicitly computed complex.

Finally keep a small regular disc in the pair of
pants uncollapsed while forming its ribbon spine.
The same construction relative to that disc retains
the block $F\times D^2$ with its flat vertical
frame and base frame.  The three-torus
$\{\gamma=\theta_0\}\subset F$ has normal frame
given by the two base directions and the
$\gamma$ direction.  This is exactly the regular
fibre and normal framing of B.1.  The CW
presentation therefore retains the geometric
framed fibre; no smooth framing of hypothetical
cancelled handles has been inferred from the
homology calculation.
\end{proof}

\subsection{B.6. Smooth monodromy: every requested class is zero}

This subsection settles B.6 for the original family and every
sphere-producing higher-multiplicity family of K1. We use the complex
orientation of $X$, the positive loop in the $w$-plane, and the generator
of $\Theta_7$ with Eells--Kuiper invariant $1/28$, as in the conventions
and the supplied result K3. Write $a=[\Phi]\in\Z/28$. The mapping torus
has the oriented diffeomorphism type $(S^1\times S^6)\#\Sigma_\Phi$.
The reduction from this mapping torus to the glued sphere, and the
independence of its Eells--Kuiper invariant from the two spin structures,
are those of K3; no different convention is introduced here.

\paragraph{Proposition B.6.1 (Proved; precise scope of K7).}
If the monodromy of the sixfold parameter pullback has computed class
$b=6a\in\Z/28$, then $b$ is even and, writing $b=2r$, the original
class satisfies
\[
 a\equiv5r\pmod {14}.
\]
There are exactly two residues modulo $28$ satisfying this condition,
and they differ by $14$. Without a value of $b$, this assertion excludes
none of the $28$ possibilities for $a$.
\begin{proof}
Pulling a mapping torus back by a degree-six map of its oriented base
circle replaces its monodromy by $\Phi^6$. Thus its class is $6a$.
The kernel of multiplication by six on $\Z/28$ is $\{0,14\}$ and
its image consists of the even residues. Dividing $6a=2r$ by two gives
$3a=r$ modulo $14$; multiplication by five inverts three modulo $14$.
Every $a$ has an image under this multiplication, so the conditional
statement does not itself restrict $a$. This is the exact arithmetic
content of K7, not a computation of the sixfold monodromy.
\end{proof}

\paragraph{Theorem B.6.2 (Proved: complete answer to B.6).}
The original class is $a=0$, and the corresponding class for
every sphere-producing family $X_{3k,4l}$ is also zero for
the positive loop with period $2kl$.  The original computation
is B.6.21; the actual higher-family replacement and cup
calculation are B.6.22--B.6.23, and the complete spin-trace
comparison proving the assertion is B.6.24.  Corollary B.6.25
then gives an explicit spin filling of every resulting
mapping torus, with one zero-handle and one one-handle.
Propositions B.6.3--B.6.6 give a separate analytic conifold
reduction for every pair; they do not by themselves evaluate
or compare the numerical classes.  The following records retain
the earlier unsuccessful parameter-filling approaches, followed
by the completed smooth calculation for all requested pairs.

The strongest explicit closing datum available is K8. With
$\ell=(2\pi i)^{-1}\log(F/G)$, $F^3=t$, $G^4=t-1$, its coherent
even-shift comparison has translation
\[
 b_{h,c}=-\frac h{12}\Pi_c\widehat\gamma+
          \left(h\ell-\frac{h^2}{24}\right)\widehat\delta,
 \qquad h\in2\Z.
\]
For the sixfold closing map, $h=12$, subtracting the integral period
$-\Pi_c\widehat\gamma$ produces the scalar $t^4/(t-1)^3$ on the
ordinary locus. At the cusp the same comparison is a holomorphic
toric unit. These representatives differ by one
$\widehat\gamma$ deck transformation. This proves equality of the
endpoint maps, but scaling the logarithms separately does not give
an isotopy on the overlap: an intermediate comparison would require
a nonintegral multiple of that deck transformation. Such a multiple
is not an identification in the lattice quotient. We have therefore
not replaced this mismatch by a purported global scalar gauge.

The CW model of B.5 gives the actual underlying topological attachment
and the integral maps needed for the homotopy-sphere calculation. Its
use of the positive toric quotient and a relative product
homeomorphism does not supply the missing smooth framing of the
parameter-dependent attachment. Carrying its cellular boundary
matrices around the parameter circle would give homology information,
but would not distinguish the $28$ smooth mapping classes.

We also checked the real representatives and their cap transport in
A.3. For the positive-imaginary parameter, the actual order-three
transport is $C=(A_2^{-1}A_1^{-1}A_2)^2$: the two reflection matrices
have product $B_1=A_2^{-1}A_1^{-1}A_2$, and conjugating one reflection
to the other requires $2k=1$ modulo three. Thus $k=2$ is forced.
All four reflection lifts at the multiplicity-four fibre were checked
independently in \texttt{maps\_real\_audit.py}. These verifications
fix the real attachment markings but do not determine the components
of the resulting orientation-reversing diffeomorphisms in the full
diffeomorphism group. In particular, their fixed-set topology alone
does not evaluate their product in $\Theta_7$.
The precise comparison theorem is
\cite[Theorem 1.3]{mapsOlbermann10}: a fixed three-sphere has even
Haefliger knot class, and each even knot supports four involutions up
to equivariant diffeomorphism relative to that sphere. Thus even the
ambient knot class, if computed, would leave additional involution data.

A further possible comparison is the explicit smooth plumbing
relation of Krylov \cite[Theorem 3]{mapsKrylov03}, used in
\cite[Theorem 4.1]{mapsRandalWilliams25}. On $S^3\times S^3$, the
quaternionic maps $Y(u,v)=(u,uv)$ and $U(u,v)=(v^{-1}u,v)$ satisfy
$(YUY)^4=\Sigma_{\rm Milnor}^{-1}$. Their relation detects a smooth
sphere class that their homology matrices do not detect. To apply
this relation here requires an actual relative smooth identification
of our parameter transport with a word in such plumbing maps; no such
identification has been proved. We also read the finite semistable
models in the supplied checkpoint, B.1--B.3 of
\texttt{S6\_02\_PUSH\_AB\_push.tex}. Their component monodromy
obstruction prevents gluing them to the plain cusp models and hence
prevents using those particular models as a bounding manifold.
These explain why those earlier approaches did not evaluate $a$.
The completed fixed-circle calculation below does not require those
unproved identifications.

\paragraph{Proposition B.6.3 (Proved; the smoothing monodromy is Torelli).}
Fix the sphere twist data and a sufficiently small nonzero smoothing
parameter $s$ in D.3. The even-shift comparisons extend to the joint
moving-zero family. Its monodromy $\Psi_s$ around the positive $w$
loop acts as the identity on $H_3(Y_{u,s};\Z)\simeq\Z^2$.
Thus the conifold route places this monodromy in the Torelli subgroup
of the mapping class group of $S^3\times S^3$; its homology action
does not determine its remaining smooth class.
\begin{proof}
The moving section multiplies the linearization by $1-st$. On an
additive period cover set
\[
 L_s=(2\pi i)^{-1}\log(1-st),\qquad
 \beta_{c,s}=\beta_c-L_s,\qquad
 \ell_s=\ell-L_s/12.
\]
The sign follows directly from the multiplicative period
$\exp(-2\pi i\beta_c)$ of the $\widehat\gamma$ generator.
The same formula for the even-shift comparison is
\[
 b_{h,c,s}=-\frac h{12}\Pi_{c,s}\widehat\gamma+
             \left(h\ell_s-\frac{h^2}{24}\right)\widehat\delta.
\]
The two added $L_s$ terms cancel, so this additive vector equals
$b_{h,c}$ on the common cover. Near the finite fibres $1-st$ is a
holomorphic unit and the checked finite collar proof is unchanged.
At the new zero $t=1/s$, the logarithmic long period cancels the
logarithm in $\ell_s$, giving an additive translation holomorphic
on the transverse elliptic factor and a holomorphic unit in the
vanishing toric coordinate. At infinity the removed logarithm in
$\beta_{c,s}$ is matched by the removed logarithm in $\ell_s$;
the same assertion holds for the other rank-one degeneration.
These are translations preserving the primitive rank-one fans, so
the maps and their inverses extend on their toric charts. The finite
affine defects and the composition law are the same as before because
the additive vector is unchanged. This constructs the even-shift
comparisons on the moving-zero family, not just on its ordinary locus.
For $h=12$ the scalar representative is now
\[
 \frac{t^4}{(t-1)^3(1-st)},
 \qquad \operatorname{div}=4[0]-3[1]-[1/s].
 \tag{B.6.1}
\]
In particular its pole has moved from infinity to the new zero.

Write $E=d\otimes\gamma$. The four fibre monodromies, in order
with product one, are
\[
 A_1,\quad A_2,\quad C_w=I+w\otimes u^*,\quad C_d=I-E.
 \tag{B.6.2}
\]
Indeed the cusp period formerly containing $-\log q/(2\pi i)$
now contains $-\log(q-s)/(2\pi i)$, whereas the other nodal period
still contains $\log q/(2\pi i)$. Thus the two primitive rank-one
monodromies are separated, and $C_wC_d=M_0$. The parameter return
has linear part $S=I-2E$; its affine translation acts trivially on
fibre cohomology. All four matrices commute with $S$.

Let $j$ include the complement of these four values in $B$, and let
$V$ be the rational degree-one cohomology local system. For each
$q=0,\ldots,4$ we have
\[
 R^q g_*\Q=j_*\bigwedge^qV.
 \tag{B.6.3}
\]
At a finite fibre this follows by transfer from its finite free torus
cover: the central and nearby invariant groups agree. At a primitive
rank-one fibre, its normalization is a ruled surface over an elliptic
curve, with degree-zero normal line bundle and two sections glued by
a translation. Both the line bundle and that translation are
topologically trivial after choosing a trivialization and a translation
isotopy. Its specialization is therefore the product calculation for
a nodal two-torus and a smooth two-torus. In the nodal factor the
groups in degrees zero, one and two map isomorphically to the
monodromy-invariant groups; the product calculation proves (B.6.3)
in all degrees. This argument only asserts the rational comparison
needed below, and includes both separated rank-one fibres.

Here is an explicit calculation of the middle sheaf cohomology. Put
$G_i=\bigwedge^2(A_i^{-t})$ at the two finite values, and use
$\bigwedge^2(C_w^{-t}),\bigwedge^2(C_d^{-t})$ at the other two.
Parabolic cocycles are tuples $(z_i)$ with
$z_i\in\operatorname{im}(G_i-I)$ and
\[
 z_1+G_1z_2+G_1G_2z_3+G_1G_2G_3z_4=0,
\]
modulo tuples $((G_i-I)v)_i$. This computes
$H^1(B,j_*\bigwedge^2V)$: the local condition is exactly zero
restriction to each punctured-disc cohomology group.
In the exterior basis of covectors
$(\gamma u^*,\gamma w^*,\gamma d^*,u^*w^*,u^*d^*,w^*d^*)$,
where juxtaposition denotes exterior product, two representatives
are the columns of the following four matrices:
\[
 P_1=\begin{pmatrix}1&-8\\2&-7\\0&-6\\0&1\\-2&1\\-1&-1\end{pmatrix},
 \quad
 P_2=\begin{pmatrix}6&-3\\3&0\\6&0\\-1&0\\-1&-1\\1&-1\end{pmatrix},
 \quad
 P_3=\begin{pmatrix}0&0\\0&0\\0&0\\0&0\\0&-1\\0&0\end{pmatrix},
 \quad P_4=0.
 \tag{B.6.4}
\]
The four local image ranks are $4,4,2,2$. The relation matrix
has rank five, and the coboundary matrix has rank five; the two
displayed cocycles are independent modulo its image. Hence the
parabolic group has dimension two and these form a basis. Direct
multiplication gives, for its two basis cocycles $p_1,p_2$,
\[
 (\bigwedge^2S^{-t}-I)p_1=d(-2\gamma u^*),\qquad
 (\bigwedge^2S^{-t}-I)p_2=d(-2\gamma w^*).
 \tag{B.6.5}
\]
Its induced parameter action is therefore the identity. All the
matrices and rank assertions, in all exterior degrees, are checked
exactly in \texttt{maps\_smoothing\_monodromy.py}.

The middle parabolic term is $E_2^{1,2}$ in the Leray spectral
sequence and survives to $E_\infty$: the base has real cohomological
dimension two. It already has dimension two, which is the full
$b_3(Y)$ proved in D.3.2. Consequently the other filtration quotients
of $H^3(Y;\Q)$ vanish and the natural filtration identifies this
group with $E_\infty^{1,2}$. Naturality identifies its parameter
action with (B.6.5). Since $H^3(Y;\Z)$ is torsion free, identity
over $\Q$ implies identity over $\Z$, and duality gives the claim
for homology.
\end{proof}

\paragraph{Proposition B.6.4 (Proved; the normal-framing obstruction vanishes).}
For the monodromy $\Psi_s$ in B.6.3, let $M_s$ be its oriented mapping
torus. Then $p_1(M_s)=0$ in $H^4(M_s;\Z)$. Its Torelli
normal-framing invariant is zero. Consequently $[\Psi_s]$ belongs
to the injected subgroup $\Theta_7$ of
$\pi_0\operatorname{Diff}^+(S^3\times S^3)$.
This statement does not evaluate that central class or identify it
with $[\Phi]$ through the conifold surgery without further relative
framing information.
\begin{proof}
We first compute the two swept critical-curve classes. Use $k$ for
the parameter generator and $a_j$ for the finite base generators in
the boundary markings of B.5. Their conjugation defects are
\[
 k a_j k^{-1}=t_{d_j}a_j,\qquad
 d_1=(0,1,-1,-1)^t,\quad d_2=(0,0,-1,1)^t.
 \tag{B.6.6}
\]
Indeed the finite return is $x\mapsto Sx+\beta_j$, where the
noninvariant part of $\beta_j$ is $-g/6$. The remaining part is
fixed by $A_j$: it is the real inverse image of a complex multiple
of the invariant vector $d$. Conjugating the affine finite generator
$x\mapsto A_jx+v_j/m_j$ therefore gives
\[
 d_j=(S-I)v_j/m_j-(I-A_j)g/6.
\]
Substitution gives (B.6.6) and $d_1+A_1d_2=0$. At the two
rank-one values the holomorphic unit representatives have
single-valued additive lifts in the toric and transverse coordinates;
the compatible zero-section boundary marking gives zero $k$-defect.

The regular five-dimensional fibre $N$ of $M_s\to B$ has fundamental
group $\Lambda\rtimes_S\Z$. It is the product of a three-dimensional
nilmanifold and the $(u,w)$ two-torus. Its invariant one-forms can be
written
\[
 \gamma,\upsilon,\omega,\vartheta,\kappa,
 \qquad d\vartheta=2\kappa\wedge\gamma,
 \qquad d\gamma=d\upsilon=d\omega=d\kappa=0.
 \tag{B.6.7}
\]
The first three closed forms are dual to $g,u,w$, and $\kappa$ is
the parameter class. These formulas follow either from
$[k,g]=-2d$ or from the connection form of the central circle bundle.
The finite base actions on its Lie lattice are
\[
 \widetilde A_j=
 \begin{pmatrix}A_j&-d_j\\0&1\end{pmatrix};
\]
the two other actions are the corresponding block lifts of $C_w,C_d$.
For the indicated cohomology classes this gives
\[
 (C_w^*-I)[\gamma\omega]=-[\gamma\upsilon],\qquad
 (\widetilde A_2^*-I)[\gamma\upsilon]
       =[\gamma\omega]-[\gamma\upsilon].
 \tag{B.6.8}
\]
Here the star denotes cohomological transport, namely inverse
transpose on the one-forms in the chosen homology marking.
In the second identity the additional form $-\gamma\kappa$ is
exact by (B.6.7). Thus both $[\gamma\upsilon]$ and
$[\gamma\omega]$ vanish in the global coinvariants of $H^2(N;\Q)$.
The full rational cohomology and these exact actions are computed in
\texttt{maps\_mapping\_torus\_cohom.py}; only the two identities
(B.6.8) are needed for the argument.

Let $p_s=1/s$. The central surface of the rank-one filling there is
the degree-zero ruled surface over the $(u,w)$ elliptic curve with
its sections identified by translation. The parameter return acts
on its normalization by a translation and a constant fibre scaling.
Both admit paths of such maps commuting with the conductor
identification. Its local mapping torus consequently has the
cohomology of this central surface times a circle. In degree three
it contains two independent classes
\[
 z_u=\kappa\gamma\upsilon,\qquad
 z_w=\kappa\gamma\omega.
 \tag{B.6.9}
\]
Here $\gamma$ is the nodal graph-circle class. A toroidal
neighbourhood retracts to the central surface, by radial collapse
in its node and core charts; choosing the collapse continuously
around the parameter circle identifies these with classes of the
whole local mapping torus. Their restrictions to a regular $N$
are zero by (B.6.7). They therefore define independent skyscraper
kernel sections in the third direct image sheaf for $M_s\to B$,
supported at $p_s$.

We compute their Leray differential, including its boundary term.
On the boundary of this filling, add the closed base-meridian form
$\alpha$. The two commuting monodromies $C_d$ and $S=C_d^2$ give
\[
 d\vartheta=(\alpha+2\kappa)\wedge\gamma.
\]
Consequently the restrictions of (B.6.9), in the Wang summand for
the base boundary circle, are respectively
\[
 -\tfrac12[\gamma\upsilon],\qquad
 -\tfrac12[\gamma\omega].
 \tag{B.6.10}
\]
For example
$\kappa\gamma\upsilon=-\alpha\gamma\upsilon/2+
d(\vartheta\upsilon)/2$ on that boundary. All statements here are
over $\Q$, so the division by two is legitimate.

For clarity, the comparison of this Wang term with the global
$d_2$ is the following cochain construction. Cover $B$ by a small
disc at $p_s$ and its complement, and then refine the complement
by the other three discs and a graph. A local degree-three class
with zero regular-fibre restriction restricts on the annulus to
its degree-two Wang class. The connecting homomorphism of the
base cover sends this to $H^2(B,R^2)$, which is exactly its Leray
$d_2$. The natural generization $R^2\to j_*H^2(N;\Q)$ has kernel
and cokernel supported at finitely many points. Their positive
cohomology vanishes, so
\[
 H^2(B,R^2)=H^2(B,j_*H^2(N;\Q))
          =H^2(N;\Q)_{\pi_1(B^*)}.
\]
In the graph-and-discs cochain model the connecting map is the
natural map from the local coinvariants to these global
coinvariants, with the sign fixed by the chosen boundary
orientation. Equations (B.6.8)--(B.6.10) prove $d_2z_u=d_2z_w=0$.
There is no later differential because the base has cohomological
dimension two.

By B.6.3 and the Wang sequence, $b_3(M_s)=2$. The two surviving
classes just constructed therefore exhaust its rational third
cohomology. Their local restrictions are zero at infinity. At
$p_s$ they are (B.6.9), whose restrictions to the swept critical
elliptic curve vanish: the graph-circle class is zero on each
elliptic fibre, so its restriction on the swept three-manifold is
a multiple of $\kappa$, and wedging with $\kappa$ kills it.
Thus both swept critical three-manifolds pair to zero with all of
$H^3(M_s;\Q)$. Their classes in $H_3(M_s;\Q)$ are zero.

It remains to relate these classes to $p_1$. Let $\mathcal E$ be
the complex rank-three tangent bundle along the fibres of the
parameter family, restricted to $M_s$. Factor the differential
of the map to $B$ by its divisorial zeros at the two multiple
fibres; its target is then a line bundle $\mathcal L$. The
saturated vertical tangent sheaf $\mathcal V$ is a rank-two
bundle. At either rank-one critical curve the local equation is
$xy=q$ times an elliptic coordinate, and its kernel is freely
generated by
$x\partial_x-y\partial_y$ and the elliptic tangent direction.
The invariant vector field $d$ is nowhere zero as a section of
$\mathcal V$: it is the first generator at $p_s$ and the second
at infinity. Hence $\mathcal V$ is an extension of a line bundle
by the trivial line, and $c_2(\mathcal V)=0$.

Write $\mathcal C$ for the disjoint union of the two swept critical
curves. The local ideal of each is $(x,y)$, giving the global
exact sequence, or equivalently its restricted topological $K$-class,
\[
 0\longrightarrow\mathcal V\longrightarrow\mathcal E
 \longrightarrow\mathcal L\otimes I_{\mathcal C}\longrightarrow0,
 \qquad
 [\mathcal E]=[\mathcal V]+[\mathcal L]-[\mathcal L|_{\mathcal C}].
\]
The Wang sequence and D.3.2 give $H^2(M_s;\Z)=0$. All the line
first Chern classes in this identity vanish. The degree-four
Chern-character term of the final sheaf is the oriented dual of
$\mathcal C$, as follows from its codimension-two Koszul resolution.
It follows that
\[
 c_2(\mathcal E)=\operatorname{PD}[\mathcal C],\qquad
 p_1(M_s)=-2\operatorname{PD}[\mathcal C]=0
       \quad\hbox{in }H^4(M_s;\Q).
\]
Again by the Wang sequence $H^4(M_s;\Z)\simeq\Z^2$, so this
vanishing is integral.

Finally apply the Torelli exact sequence quoted from Kreck in
\cite[Section 2.1, equation (2)]{mapsKrylov03}. For $S^3\times S^3$
its left term is the injected $\Theta_7$, because this manifold
bounds the framed manifold $S^3\times D^4$. The right term is
the homomorphism from $H_3$ to the image of
$\pi_3SO(3)\to\pi_3SO(4)$. Its mapping-torus characteristic class
is detected by $p_1$: the stable clutching bundle of a generator
of $\pi_3SO(3)$ has first Pontryagin number of absolute value four,
as computed in \cite[proof of Theorem 2]{mapsKrylov03}. Thus the
zero $p_1$ class forces the entire integral normal-framing
homomorphism to vanish. Exactness gives the assertion.
\end{proof}

\paragraph{Proposition B.6.5 (Proved; the conifold comparison preserves the sphere class).}
Let $a\in\Theta_7$ be the original six-sphere monodromy class,
and let $b\in\Theta_7$ be the central class of the smoothing
monodromy supplied by B.6.4. With the same positive parameter
loop and the complex orientations, $a=b$. Equivalently, if
$M_\Phi$ and $M_s$ are their seven-dimensional mapping tori, then
\[
 \mu(M_\Phi)=\mu(M_s)=\frac{a}{28}\quad\hbox{in }\Q/\Z.
 \tag{B.6.11}
\]
Thus no additional unknown framing correction is introduced by
the reverse conifold surgery. This identifies the remaining
invariant; it does not evaluate it.

\begin{proof}
We first construct the surgery trace around the whole parameter
circle, including its gluing. The even-shift comparisons in B.6.3
preserve the actual smoothing parameter $s$, not merely its zero
set. In the node chart, the single-valued substitution $x'=x/u$
for $u\ne0$ makes the equation
\[
 x'y-zv=-s.
\]
The relative holomorphic Morse coordinates consequently have
transition germs preserving this fixed quadratic function.
There is no conformal scalar multiplying the quadratic form.
Such a germ $F$ retracts to its derivative by
$F_r(z)=r^{-1}F(rz)$, $0<r\leq1$, with $F_0=dF_0$;
each germ in this deformation still preserves the quadratic
function. The chosen small resolution is part of the family,
so these derivatives preserve its ruling. They therefore belong
to $SO(4,\mathbb C)$ rather than the component which exchanges
the two rulings. Reduce this group to its maximal compact
$SO(4)$ and choose small invariant collars. The deformation of
germs gives the same boundary bundles up to isotopy, and collar
isotopies extend to their adjacent pieces.

Here is an explicit equivariant handle trace for this reduction.
Write $D^3\subset\operatorname{Im}\mathbb H$ and
$D^4\subset\mathbb H$. The action
\[
 (g,h)(v,q)=(gvg^{-1},gqh^{-1}),\qquad
 (g,h)\in SU(2)\times SU(2),
 \tag{B.6.12}
\]
on $D^3\times D^4$ factors through $SO(4)$, since $(-1,-1)$
acts trivially. Its boundary faces are the two local surgery
models. On $S^2\times D^4$, choose the identification of the
twistor two-sphere with the exceptional $\mathbb{CP}^1$ for
which its stabilizer line is the tautological line of degree
minus one. With this convention the fibrewise complex structure
$J_v(q)=vq$ identifies the real normal bundle with
$\mathcal O(-1)\oplus\mathcal O(-1)$ over $\mathbb{CP}^1$:
at $v=i$ the stabilizer $U(1)$ acts by scalar multiplication
on the complex two-dimensional fibre, and its associated
tautological line has degree minus one. On $D^3\times S^3$,
the map $(v,q)\mapsto vq\in T_qS^3$ identifies the same action
with that on the cotangent-disc model of the Milnor fibre.
Their common face is $S^2\times S^3$, with the common action
induced on the quadric link. Identify both faces with the
given complex-oriented local models along that link; reversing
the twistor convention would require reversing its sphere
identification too. The outward normal conventions
are those of replacing $S^2\times D^4$ by $D^3\times S^3$.
Thus (B.6.12), together with the unchanged complement, builds
a bundle of oriented seven-dimensional surgery traces over
the parameter circle. Denote its total space by $Z$. Its
boundary is $M_s\sqcup(-M_\Phi)$.

For one trace $W$, the handle relative to $S^6\times I$ has
index three. Its cellular description gives
\[
 H^3(W;\Z)=\Z,\qquad H^4(W;\Z)=0,
 \qquad H^6(W;\Z)=\Z.
\]
The restriction $H^3(W;\Z)\to H^3(S^3\times S^3;\Z)$
is a primitive injection. Indeed the reverse trace has one
index-four handle, so its homology sequence gives a surjection
$H_3(S^3\times S^3)\to H_3(W)=\Z$ with primitive rank-one
kernel; dualize this sequence. B.6.3 then implies that trace
monodromy acts as the identity on $H^3(W;\Z)$. The Wang
sequence, naturally with respect to its smoothing boundary,
now gives a commutative square
\[
 \begin{array}{ccc}
 H^4(Z;\Z)&\longrightarrow&H^4(M_s;\Z)\\
 \big\downarrow\scriptstyle\simeq&&
       \big\downarrow\scriptstyle\simeq\\
 H^3(W;\Z)&\longrightarrow&H^3(S^3\times S^3;\Z).
 \end{array}
 \tag{B.6.13}
\]
The top restriction is therefore an injection of a rank-one
free group. Its value on $p_1(Z)$ is $p_1(M_s)=0$, by B.6.4,
so $p_1(Z)=0$. Moreover the exact cohomology sequence of
$(Z,\partial Z)$ shows that
$H^4(Z,\partial Z;\Q)\to H^4(Z;\Q)$ has zero image, since
restriction even to the single component $M_s$ is injective.
The intersection form defining the signature has consequently
zero rank, and $\sigma(Z)=0$.

The trace $W$ is simply connected and spin: the index-three
attachment extends the unique spin structure, as its attaching
two-sphere has no first cohomology. Trace monodromy preserves
this unique structure. Either of its lifts makes $Z$ spin,
with compatible induced boundary spin structures. For a spin
seven-manifold with rationally zero first Pontryagin class,
the Eells--Kuiper invariant is
\[
 \mu(M)=\frac{p_1(V)^2-4\sigma(V)}{896}\pmod{\Z},
 \tag{B.6.14}
\]
using a spin coboundary $V$ and a rational relative lift of
$p_1(V)$. The square is independent of the lift. This definition
does not require $H^4(M;\Q)=0$; see
\cite[Section 1.4, equations (2)--(3)]{mapsCrowleyNordstrom19},
whose normalization is $28$ times (B.6.14). Applying its gluing
formula to $Z$, with the zero lift of $p_1(Z)$, proves
$\mu(M_s)-\mu(M_\Phi)=0$.

Finally, an oriented diffeomorphism supported in a six-disc
changes its mapping torus from the product by connected sum
with its corresponding homotopy seven-sphere. This follows
directly by cutting out that disc times a small parameter
interval and regluing its two faces. Every class of
$\pi_0\operatorname{Diff}^+(S^6)$ has such a representative,
and B.6.4 provides such a representative for $\Psi_s$ as well.
The product mapping tori have zero $\mu$: they bound respectively
$D^7\times S^1$ and $S^3\times D^4\times S^1$, with zero
Pontryagin square and signature. These bounds work for either
circle spin structure. Additivity of (B.6.14), and its
identification $\Theta_7\simeq\frac1{28}\Z/\Z$, prove
(B.6.11) and $a=b$.
\end{proof}

\paragraph{Proposition B.6.6 (Proved; the reduction for all coprime multiplicities).}
For every sphere member $X_{3k,4l}$ of K1, the monodromy associated
to the positive generator $2kl$ has the same class in $\Theta_7$
as a central monodromy on its conifold smoothing $S^3\times S^3$.
The smoothing monodromy acts trivially on integral $H_3$, and
its seven-dimensional mapping torus has $p_1=0$. The numerical
values of these central classes, and any relation between
different pairs $(k,l)$, remain undetermined here.

\begin{proof}
Put $h=2kl$. The D.3 construction applies with the finite
replacements of orders $3k,4l$: their affine actions are still
free, and multiplying the linearization by $1-st$ is a unit
change on each of their fixed root discs. The node and its
small resolution are unchanged. The one-surgery proof of D.3.2
therefore identifies the smoothing with $S^3\times S^3$.

We spell out why the comparison from K1 survives the separated
cusps. Use the fixed-data case of its proof, so the first
coordinates of its rational finite translation vectors $b_j$,
integral defects $\ell_j$, and connector vectors $d_j$ are all
zero. Changing $\beta$ to $\beta-L_s$ changes only the
$\widehat\gamma$ column of the period matrix. It consequently
leaves these finite translation formulas unchanged. In K1 the
remaining cusp defect has already been adjusted to belong to
$\mathbb Z w\oplus\mathbb Z d$; write it as $aw+bd$.
Assign $aw$ to the old elliptic node and $bd$ to the new
$\delta$ node. Their product is the old defect, since
$C_wd=d$. Each is an allowed integral winding in its primitive
rank-one toric direction. Realize it locally by the appropriate
integer power of the base coordinate, times a holomorphic unit;
the transverse elliptic translation is holomorphic. These
are precisely the extending rank-one translations and they
agree with the unsplit two-direction prescription when the
two zeros coincide.

The differences of these local prescriptions have the same
coarse vertical-vector sheaf as in K1:
\[
 0\longrightarrow\mathcal O_B\longrightarrow\mathcal V_B
 \longrightarrow\mathcal O_B(-p_0)\longrightarrow0.
\]
To check this assertion locally, quotient by the invariant
$\delta$ vector. The quotient is the Lie line of the unchanged
elliptic surface; its finite-root exponents and its simple
pole at the old elliptic node are unchanged. The new
$\delta$-direction node has a nonvanishing logarithmic
$\delta$ generator in the abstract vertical bundle and a
regular transverse elliptic generator, so it adds no zero
or pole to that quotient. The extra subgroups in the enlarged
finite roots act trivially on vector coefficients, just as
in K1. Both end terms have zero $H^1$, hence the Cech
correction glues the new local prescriptions. On a small
parameter disc $\Delta_s$, including the coalescence, this is
the relative exact sequence
\[
 0\longrightarrow\operatorname{pr}_B^*\mathcal O_B
 \longrightarrow\mathcal V_{B\times\Delta_s}
 \longrightarrow\operatorname{pr}_B^*\mathcal O_B(-p_0)
 \longrightarrow0.
\]
The local logarithmic generators described above also define
this sequence at $s=0$. The two-chart Cech complex for either
end term has zero first cohomology, so
$R^1\operatorname{pr}_{\Delta_s*}\mathcal V_{B\times\Delta_s}=0$.
The same holds after adjoining a small Stein $c$-disc. Solving
that relative complex supplies a correction holomorphic in
$c,s$, through the coalescence. The integral windings are constant
and the local units vary holomorphically; at the coalescence
their two toric directions are the original ones. The
comparison thus covers the identity on $s$ and on $B$, with
the chosen small resolution. Different corrections differ
by the connected vertical translation group, so they have
the same smooth monodromy class.

The four homological matrices are still (B.6.2). Finite transfer
still gives (B.6.3), since the root actions are free even when
their multiplicity exceeds the linear order. Replacing
$S=I-2E$ by $I-hE$ multiplies the two coboundaries in (B.6.5)
by $h/2$. Thus the same rank-two parabolic group proves the
integral Torelli assertion.

For the characteristic-class calculation the regular nilpotent
fibre now has
\[
 d\vartheta=h\kappa\gamma.
\]
Changing the finite twist data changes the affine parameter
defects but not their first coordinate, which remains zero.
The base actions therefore still fix $\gamma$; on wedging
their action on $\upsilon,\omega$ with $\gamma$, the extra
terms are multiples of $\gamma\kappa$, which are exact.
The two coinvariant vanishings in (B.6.8) follow unchanged.
Choose the local parameter lift at the new rank-one point
by subtracting its integral toric winding, so its long
$\gamma$ translation has zero representative. Its boundary
equation is then
$d\vartheta=(\alpha+h\kappa)\gamma$.
The two local skyscraper classes $\kappa\gamma\upsilon$ and
$\kappa\gamma\omega$ consequently have Wang components
$-[\gamma\upsilon]/h$ and $-[\gamma\omega]/h$.
They survive and exhaust the mapping-torus $H^3$ exactly as
in B.6.4. Their restrictions to both swept critical curves
are zero. The same saturated-vertical tangent sequence proves
$p_1=0$ integrally, and the Kreck sequence places the smoothing
class in the central $\Theta_7$. Finally the fixed-$s$ trace
argument of B.6.5 applies unchanged and identifies this class
with that of the original sphere family.
\end{proof}

\paragraph{Proposition B.6.7 (Refuted; the scalar-shear claim already fails on the ordinary locus).}
The first isotopy assertion of Lead 8 is false for the marked
return of the original family. On $X^\circ=f^{-1}(B^\circ)$,
that return is not homotopic over $B^\circ$ to
\[
 x\longmapsto\rho(-2\theta(x))x
\]
for a globally defined monodromy-invariant $\gamma$ coordinate
$\theta:X^\circ\to\mathbb R/\mathbb Z$.
In fact, after allowing arbitrary additional circle-valued
scalar gauges on the base, a quotient of the affine defect
has exact order six. Its vanishing under the sixth iterate
does not compute the sixth iterate in $\Theta_7$.

\begin{proof}
Take a compact pair-of-pants core $P\subset B^\circ$ and the
marked zero section on its ordinary torus bundle. A connection
for the compact family may be chosen to equal the flat real
period connection over this core and then extended over the
three collars. Consequently its induced return on the core
fundamental group is the one calculated from the closing map
in K8. With base generators $a_1,a_2$ and fibre translations
$t_\lambda$, that group and its return are
\[
 \pi_1(f^{-1}(P))=\Lambda\rtimes F(a_1,a_2),\qquad
 \Phi_*t_\lambda=t_{S\lambda},\qquad
 \Phi_*a_j=t_{d_j}a_j,
\]
where $S=I-2E$ and $d_1,d_2$ are exactly (B.6.6).
Different choices of the path from a basepoint to its image
conjugate this homomorphism by an integral fibre translation.

The pure scalar shear has the same linear part $S$.
For the zero-section coordinate its two affine defects are
zero. Changing the coordinate by a map $P\to S^1$ adds
defects in $\mathbb Z d$; allowing arbitrary such defects
only weakens the proposed comparison. A homotopy over $P$
tracks the basepoint inside its fibre. It would therefore
require, after quotienting by the invariant $\mathbb Z d$,
an integral vector $\bar\lambda\in\bar\Lambda:=\Lambda/\mathbb Z d$
with
\[
 (\bar A_j-I)\bar\lambda=\bar d_j,\qquad j=1,2.
 \tag{B.6.15}
\]
The possible sign from the direction of basepoint conjugation
has been absorbed into $\bar\lambda$.
Writing $\bar\lambda=(x,y,z)$ in the $(g,u,w)$ basis,
the nonzero equations are
\[
 6x-y+z=1,\qquad -6x-y-2z=-1,\qquad
 -y-z=0,\qquad -6x+y-z=-1.
\]
Their unique rational solution is
\[
 (x,y,z)=(1/6,0,0).
 \tag{B.6.16}
\]
It is not integral, proving the stated obstruction.

This is independent of the auxiliary ordinary marking.
Changing the basepoint path changes $(\bar d_j)$ by one
simultaneous coboundary. Replacing the lifted base generators
by $t_{r_j}a_j$ changes $d_j$ by $(S-I)r_j$, which lies in
$\mathbb Z d$ and disappears in the quotient. An integral
linear change of marking transports both the local system
and the displayed class. Thus none of these changes removes
the nonzero class in $H^1(P,\bar\Lambda)$.

Moreover $6\bar d_j=(\bar A_j-I)g$, and (B.6.16) shows
that no smaller positive multiple is a coboundary. Its
order is exactly six; the three nonzero Smith factors of
the stacked quotient coboundary matrix are $1,1,6$.
Since $Sd_j=d_j$, the defect of the $r$th iterate is $rd_j$,
so the sixth iterate is the first one for which this
particular obstruction vanishes. In the full lattice,
subtracting that integral coboundary gives
\[
 6d_1-(A_1-I)g=-4d,\qquad
 6d_2-(A_2-I)g=3d.
 \tag{B.6.17}
\]
These are the two scalar windings seen in K8, with the
inverse sign if one reverses its comparison direction.
They still have to be matched to the finite and cusp collars
to compute a diffeomorphism class of the compact sphere.

The arithmetic also explains a distinction in terminology.
For a shift $h=2r$, a literal shear $-h\theta$ with a global
$\gamma$ coordinate would require the base windings of
$\theta$ to be $1/3,-1/4$, after solving the quotient
equations. Those are not integers. A scalar gauge whose
$h$th root is not globally defined is a different map; the
sixfold scalar representative in K8 is of that latter kind.
On a simply connected ordinary base chart the translation
has an additive lift and can be interpolated, which proves
the local shear statement there. It does not give the
asserted isotopy over all of $B^\circ$.
Every rank, Smith factor and rational solution used here
is checked in \texttt{maps\_ordinary\_shear\_obstruction.py}.
\end{proof}

\paragraph{Corollary B.6.8 (Proved; every sphere pair has the same order-six affine obstruction).}
For the positive-generator return of every $X_{3k,4l}$, the
quotient affine class in $H^1(P,\Lambda/\mathbb Z d)$ has exact
order six. It obstructs an ordinary-locus isotopy to a pure
scalar shear, even after arbitrary base circle gauges.
This assertion does not identify the different pairs'
classes in $\Theta_7$ with one another.

\begin{proof}
Use the K1 coordinates with $b=4lN_0+3kN_1=\pm1$ and
$h=2kl$. The fixed-data repeated-norm solutions in its
comparison proof are
\[
 \ell_1=(0,0,-2lN_0,0)^t,\qquad
 \ell_2=(0,-kN_1,0,0)^t.
\]
Their norms are respectively $-hN_0d$ and $-hN_1d$, as
required by the shift of the affine vectors. The remaining
cusp defect is
\[
 -M_0(\ell_1+A_1\ell_2)
       =(0,0,2lN_0-kN_1,kN_1)^t,
\]
which is already in $\mathbb Zw\oplus\mathbb Zd$; no
additional adjustment of its long coordinates is needed.
Thus these are the integral defects of the allowed fixed-data
closing comparison. Its quotient cocycle has the unique
rational coboundary
\[
 \bar\lambda=
 \left(\frac{4lN_0-3kN_1}{6},\,
       2lN_0-kN_1,\,2kN_1-2lN_0\right)^t.
\]
The last two coordinates are integral and the first is
$b/6-kN_1$. Since $b=\pm1$, the least positive integer
making all three coordinates integral is exactly six.
The same basepoint and marking argument as in B.6.7 proves
both the order assertion and the obstruction to a scalar
shear. For $(k,l,N_0,N_1)=(1,1,1,-1)$ the quotient defects
here and those of (B.6.6) differ by the integral coboundary
of $(-1,-3,4)^t$, checking their comparison convention.
The symbolic identities are verified in the same exact
program as B.6.7.
\end{proof}

\paragraph{Proposition B.6.9 (Known formula; a complex filling is sufficient).}
Let $M_r$ be the oriented mapping torus for $r$ positive parameter
loops, either before or after the conifold surgery, and suppose
that a compact almost complex eight-manifold $V$ has boundary
identified orientation-preservingly with $M_r$. Then
\[
 r a=\frac{-c_1(V)^4+4c_1(V)^2c_2(V)+3c_2(V)^2
                    -3\operatorname{sign}(V)}{24}
       \quad\hbox{in }\mathbb Z/28.
 \tag{B.6.18}
\]
All characteristic products are evaluated using relative lifts;
their boundary restrictions vanish. A spin condition on $V$ is
unnecessary. The formula is conditional on the actual boundary
identification and on the indicated characteristic numbers;
no such filling or numerical value is inferred from it.

\begin{proof}
By B.6.4--B.6.6, either mapping torus has $H^2(M_r;\mathbb Z)=0$,
torsion-free $H^4(M_r;\mathbb Z)$ and $p_1(M_r)=0$. The original
mapping torus has these properties directly, and the smoothing
mapping torus is $(S^1\times S^3\times S^3)\#\Sigma_{ra}$.
Its two spin structures have the same Eells--Kuiper invariant:
each product spin structure bounds $S^1\times S^3\times D^4$,
with zero invariant, and connected sum adds $ra/28$.
The identical assertion for the original mapping torus is K3.

Spin$^c$ structures form a torsor for $H^2(M_r;\mathbb Z)$,
so the boundary Spin$^c$ structure is unique. Both spin structures
induce it: their difference in $H^1(M_r;\mathbb Z/2)$ is the
reduction of an integral degree-one class. Consequently the
Spin$^c$ structure of an almost complex $V$ has the required
boundary compatibility without an additional parity hypothesis.
Put $z=c_1(V)$ and $p=p_1(V)$. The Spin$^c$ formula in
\cite[Section 2.2, equation (2.1)]{mapsCrowleyNordstrom21},
with $\widehat p=(p+z^2)/2$ and $d=0$, gives in our normalization
\[
 \mu(M_r)=\frac{p^2-4\operatorname{sign}(V)}{896}
             -\frac{p z^2}{192}+\frac{z^4}{384}
             \quad\hbox{in }\mathbb Q/\mathbb Z.
\]
Indeed its right side before division by $28$ is
$(\widehat p^2-\operatorname{sign}(V))/8
-5z^2\widehat p/12+z^4/4$.
Substituting $p=c_1^2-2c_2$ and multiplying by $28$ gives
(B.6.18). The restrictions of $c_1$ and $p_1$ are zero;
torsion-freeness makes the restriction of $c_2$ zero as well.
Thus the relative products required by this formula exist.
The standard relative-pairing definition makes their stated
combination independent of the lifts, as in the cited formula.
\end{proof}

\paragraph{Proposition B.6.10 (Proved: the scalar-return cover and its elliptic surface).}
Fix $p\in\mathbb P^1\setminus\{0,1,\infty\}$ and let $C_p$ be
the smooth projective normalization of
\[
 h^{12}=\frac{t^4}{(t-1)^3(t-p)}.
 \tag{B.C1}
\]
Multiplying the right side by a nonzero constant does not affect
the statements below.  Write $\eta(h)=\zeta_{12}h$ for its deck
generator.  This is a connected genus-three cover, with
ramification indices $3,4,12$ at $0,1,p$, respectively.
Let $S_p\to C_p$ be the relatively minimal elliptic surface obtained
by pulling back Engel's marked surface $S$ of type $IV^*,III,I_1$.
Then $S_p$ has precisely twelve singular fibres, all of type $I_1$,
and
\[
 \chi(\mathcal O_{S_p})=1,\quad q(S_p)=p_g(S_p)=3,\quad
 (b_2^+,b_2^-)=(7,15),\quad \sigma(S_p)=-8.
 \tag{B.C2}
\]
The section induced by $P$ has height two, meets the zero section
with intersection number zero, and satisfies $(P-O)^2=-2$.
For $k=0,\ldots,11$, the dimensions of the $\zeta_{12}^k$
eigenspaces of $\eta^*$ on $H^2(S_p;\mathbb C)$ and on the
complexification of its positive subspace are, respectively,
\[
 \begin{array}{c|rrrrrrrrrrrr}
 k&0&1&2&3&4&5&6&7&8&9&10&11\\\hline
 \dim H^2_k&3&1&2&2&2&1&3&1&2&2&2&1\\
 \dim H^+_k&1&0&1&1&0&0&2&0&0&1&1&0 .
 \end{array}
 \tag{B.C3}
\]
In particular the invariant intersection form has rank three
and signature minus one.  This proposition computes the cover
of the scalar return gauge in B.6.3.  It does not identify that
cover with a complete semistable component normalization.

\begin{proof}
The three valuations of the function in (B.C1) are $4,-3,-1$.
Their greatest common divisor with twelve is one, which proves
connectedness.  At these values there are respectively $4,3,1$
points above the base, with the displayed ramification indices.
Riemann--Hurwitz gives
\[
 2g(C_p)-2=-24+(12-4)+(12-3)+(12-1)=4.
\]
There are twelve unramified points over infinity.

The order-three and order-four base changes at $0$ and $1$
are exactly the good-reduction charts used in D.3.3:
$t=z^3,x=z^4X,y=z^6Y$ at zero and
$t-1=z^4,x=z^2X,y=z^3Y$ at one.  Thus the relatively
minimal pullback is smooth over those points.  The fibre at
$p$ was already smooth, so ramified base change there remains
a smooth elliptic fibration.  Over infinity are twelve
unramified copies of $I_1$.  Euler localization gives
$c_2(S_p)=12$.  The canonical bundle of an elliptic surface
with a section and no multiple fibres is pulled back from
its base, so $c_1^2=0$.  Noether's formula yields
$\chi(\mathcal O)=1$.

For completeness the usual direct-image calculation is enough
to determine the irregularity: $\pi_*\mathcal O=\mathcal O_{C_p}$
and $R^1\pi_*\mathcal O=\mathcal L^{-1}$, where
$\deg\mathcal L=\chi(\mathcal O)=1$.  The negative-degree
line bundle has no section, so $q=g(C_p)=3$; hence $p_g=3$.
The Euler number gives $b_2=22$, and the K\"ahler decomposition
gives $b_2^+=1+2p_g=7$.  This proves (B.C2).
Shioda's base-change identity
\cite[Proposition 8.12]{mapsShioda90} multiplies the original
height $1/6$ by twelve.  Since all fibres are irreducible,
his height formula \cite[Theorem 8.6]{mapsShioda90} becomes
$\langle P,P\rangle=2+2(P\cdot O)$; it gives $P\cdot O=0$.
Both section self-intersections are $-\chi=-1$, proving the
last intersection assertion.

We compute the Hodge characters, including their signs.
Choose the rational two-form $\Omega_S$ whose divisor on the
original rational elliptic surface is minus the fibre at infinity.
Its pullback to the minimal surface has orders $0$ at each of the
four points above zero, $2$ at each of the three points above one,
$11$ at the point above $p$, and $-1$ at the twelve points above
infinity.  The first two values follow by substituting the two
good-reduction charts in $dt\wedge dx/y$; the third is the
ramification factor of $dt$, and the last is unramified pullback.
At every other fibre the form is regular and nonvanishing as
a canonical generator.  Its divisor on the base therefore has
degree $6+11-12=5$, as also follows from $K_{C_p}\otimes\mathcal L$.

Every eigenform is a rational function of $t$ times
$h^j\Omega_S$, for some $0\leq j\leq11$.  Its least allowed
factors $t^a(t-1)^b(t-p)^c$ satisfy
\[
 a=-\lfloor j/3\rfloor,\qquad
 b=\lceil(j-2)/4\rceil,\qquad
 c=\lceil(j-11)/12\rceil=0.
\]
Regularity at infinity allows an additional polynomial of
degree at most $-a-b-c-1$.  This degree is nonnegative exactly
for $j=6,9,10$, and is zero in each case.  A basis is consequently
\[
 \frac{h^6(t-1)}{t^2}\Omega_S,\qquad
 \frac{h^9(t-1)^2}{t^3}\Omega_S,\qquad
 \frac{h^{10}(t-1)^2}{t^3}\Omega_S .
 \tag{B.C4}
\]
The invariant K\"ahler class and these forms together with their
conjugates give the positive row of (B.C3).

We determine the full second-cohomology characters by fixed points.
For $1\leq j\leq11$, the number of fixed points of $\eta^j$ on
$C_p$ is
$1+4\mathbf 1_{4\mid j}+3\mathbf 1_{3\mid j}$.
Thus its first-cohomology trace is two minus this number.
Pullback identifies $H^1(C_p)$ with $H^1(S_p)$: it is injective
by the zero section and both groups have dimension six.
The traces on $H^1$ and $H^3$ agree by duality and conjugation.
Above $p$ the whole smooth elliptic fibre is fixed and has Euler
number zero.  At a point above zero, a nontrivial order-three
elliptic automorphism fixes three points.  At a point above one,
an order-four automorphism fixes two points, and its square fixes
four.  There are no other fixed points.  The local Lefschetz signs
are positive, since the normal actions are complex rotations.
Subtracting the contributions of $H^0,H^1,H^3,H^4$ gives
\[
 \operatorname{Tr}(\eta^{j*}\mid H^2)=
 \begin{cases}
 22&j=0,\\
 4&j=4,8,\\
 6&j=6,\\
 0&\text{otherwise}.
 \end{cases}
\]
The finite Fourier transform gives the first row of (B.C3).

The program \texttt{fourfolds\_B6\_cover.py} performs these twelve
character calculations in $\mathbb Q[x]/(x^4-x^2+1)$ and checks
every displayed dimension exactly; no numerical eigenvalues are used.
\end{proof}

\paragraph{Proposition B.6.11 (Proved: the untwisted plumbing has signature two).}
Let $W^\pm$ be a sufficiently small closed parameter-disc
neighbourhood in either smooth semistable total fourfold of
K13, A.2.2, with its complex orientation.  Its boundary is the
mapping torus of the \emph{untwisted} family $X(u)$ under one
positive $u$ loop, and this boundary is not spin.  One has
\[
 \sigma(W^\pm)=2,\qquad c_1(W^\pm)^2=0
       \quad\text{in }H^4(W^\pm;\mathbb Z).
 \tag{B.P1}
\]
The signature form modulo its radical is rationally represented by
$\left(\begin{smallmatrix}2&1\\1&6\end{smallmatrix}\right)$.
In particular the $E_8$ summand of the rational elliptic surface
does not give an $E_8$ signature for this parameter filling.
These are not the finite logarithmic replacements used for the
sphere; no smooth monodromy value for those replacements follows.

\begin{proof}
We first give the local-to-global intersection calculation.
Write $Q$ for the irreducible central normal-crossings fibre and
$\nu:R\to Q\subset W$ for its smooth normalization; here $R$ means
$R^+$ or $R^-$.  Its conductor sections $i_0,i_\infty:S\to R$
are glued by $T=t_{6P}$.  A sufficiently small semistable tube
retracts to its central fibre: in the product charts $xy=u$ shrink
the normal coordinates toward their zero loci, fixing that locus,
and use positive radial normal fields on the single-branch charts.
Collars make these contractions agree, and compactness gives a
finite choice of charts.  The normalization mapping cylinder
therefore computes the homology of $W$.  In particular
\[
 H_4(S;\mathbb Q)\longrightarrow H_4(R;\mathbb Q)
 \longrightarrow H_4(W;\mathbb Q)\longrightarrow H_3(S;\mathbb Q)=0.
 \tag{B.P2}
\]
Thus all middle classes are pushed forward from the normalization.

The normalization is a complex codimension-one immersion with
transverse double locus.  Let $D_0,D_\infty$ be its two conductor
divisors.  Its normal line has
\[
 c_1(N_\nu)=-[D_0]-[D_\infty].
\]
Indeed the parameter differential in $u=xy$, restricted to either
branch, is the normal differential multiplied by the coordinate
vanishing on the conductor.  Since the parameter divisor is
principal, these local identities identify the normal line with
$\mathcal O_R(-D_0-D_\infty)$.

For $a,b\in H^2(R;\mathbb Q)$, put $a_i=i_i^*a$ and $b_i=i_i^*b$.
Push the four-cycles Poincar\'e dual to $a,b$ into $W$, and perturb
one transversely.  Their same-branch intersections are counted by
the Euler class of $N_\nu$; their two other intersections lie on
the two branches of the conductor.  The complex orientations give
positive signs to both cross terms.  Thus their intersection is
\[
 -a_0b_0-a_\infty b_\infty
       +a_0T^*b_\infty+b_0T^*a_\infty
   =-\bigl(a_0-T^*a_\infty\bigr)
         \bigl(b_0-T^*b_\infty\bigr),
 \tag{B.P3}
\]
where products on the right are evaluated on $[S]$.
This also follows by doing the perturbation in each $xy=0$ chart;
the two unequal branch intersections are the two last terms.
Define $U(a)=a_0-T^*a_\infty$.
Equation (B.P2) shows that the signature is precisely the negative
of the surface intersection form restricted to $\operatorname{im}U$.
The relation coming from $H_4(S)$ causes no discrepancy:
its dual divisor is $D_0-D_\infty$, and its image under $U$ is zero
because the two conductor normal bundles have trivial product.

We compute this image using the actual surface classes
\[
 F=[S_t],\qquad m=[P]-[O],\qquad
 v=[6P]-[O]-3F.
\]
They satisfy
\[
 F^2=Fm=Fv=0,
 \qquad m^2=-2,\quad mv=-1,\quad v^2=-6.
 \tag{B.P4}
\]
These equalities follow from the section intersections
$P\cdot O=0$, $[6P]\cdot O=P\cdot[6P]=2$ and the
self-intersection $-1$ of every section.  The two latter
intersection numbers follow directly from height $1/6$ and
the narrow section $6P$, as in \cite[Lemma 2.1]{mapsEngel}.
Translation acts by the integral isometry
\[
 T_*x=x+(xF)v-(xv-3xF)F.
\]
It is checked on $F,O,v$ and the seven component-root classes,
which $t_{6P}$ fixes.  In particular
$\operatorname{im}(I-T^*)_{\mathbb Q}=\mathbb QF+\mathbb Qv$.
On $\mathbb P(\mathcal O\oplus L)$, the relative hyperplane
class has two section restrictions differing by $m$; hence its
image under $U$ is $m$ modulo $\mathbb QF+\mathbb Qv$.
The blow-up in K13 is along the fibre $S_\infty$ in one section.
Its exceptional divisor restricts to $F$ on that section and zero
on the other.  The small resolution introduces no further divisor
class: the normalization homology, or the two displayed conifold
charts, gives one additional second-cohomology generator in total.
Consequently
\[
 \operatorname{im}U=\mathbb QF+\mathbb Qv+\mathbb Qm.
\]
Its radical is $\mathbb QF$.  The matrix on $(m,v)$ is negative
definite, with leading entry $-2$ and determinant eleven.
Negating it in (B.P3) gives the positive matrix in the statement,
and proves $\sigma(W)=2$.  The signs and determinant are checked
exactly in \texttt{fourfolds\_B6\_plumbing.py}.

Finally on the initial projective bundle,
$c_1(TR)=\pi_R^*c_1(S)+D_0+D_\infty$.
Blowing up the conductor curve subtracts the exceptional divisor
both from $c_1(TR)$ and from $D_0+D_\infty$; the subsequent small
conifold resolution is crepant.  The tangent-normal sequence
therefore gives
\[
 \nu^*c_1(W)=\pi_R^*c_1(S)=\pi_R^*F.
\]
Its square is zero.  The cohomological normalization sequence,
with $H^3(S;\mathbb Z)=0$, makes
$H^4(W;\mathbb Z)\to H^4(R;\mathbb Z)$ injective.
This proves the integral assertion in (B.P1).
We check the boundary assertion integrally as well.
The normalization sequence gives $H_6(W;\mathbb Z)=\mathbb Z$,
generated by the reduced irreducible central fibre, and
$H_5(W;\mathbb Z)=0$.  For the latter, $H_5(R)=0$ and the next
map $H_4(S)\to H_4(R)$ is nonzero: after pushing through the
normalization modification its dual is $m\ne0$.
Poincar\'e--Lefschetz duality identifies
$H^2(W,\partial W)\to H^2(W)$ with the class of that central
divisor.  It is zero because the divisor is the zero set of the
global parameter function.  The next relative group is
$H^3(W,\partial W)=H_5(W)=0$.  The exact sequence thus gives
\[
 H^2(W;\mathbb Z)\xrightarrow{\ \simeq\ }
 H^2(\partial W;\mathbb Z).
\]
The class $c_1(W)$ is primitive: its pullback $\pi_R^*F$ pairs
to one with the zero-section curve in either conductor copy of
$S$.  Its restriction to the boundary is therefore not divisible
by two.  For a complex tangent bundle $w_2=c_1\bmod2$, and the
stable tangent restriction to the boundary has the same $w_2$.
Consequently $w_2(\partial W)\ne0$, proving that this is not the
spin mapping torus required in B.6.9.

When relative characteristic numbers for a compatible boundary
are defined, the absolute vanishing also kills the terms
$c_1^4$ and $p_1c_1^2$ by pairing a relative lift with the zero
absolute square.  It does not supply the missing relative
$p_1^2$ or replace the finite logarithmic caps.
\end{proof}


\paragraph{Proposition B.6.12 (Proved; a generic cyclic-norm obstruction).}
Use the sectioned elliptic surface $S_p\to C_p$ and deck action
$\eta$ of B.6.10. Outside a countable subset of
$\mathbb P^1\setminus\{0,1,\infty\}$, its Mordell--Weil group
$L_p$ satisfies
\[
 L_p^\eta=\mathbb ZP,\qquad
 (1+\eta+\cdots+\eta^{11})L_p=12\mathbb ZP.
 \tag{B.6.19}
\]
In particular $6P$ represents a nonzero class of order two
in the invariant group modulo the cyclic norm image.
There is no elliptic-surface automorphism covering $\eta$
whose twelfth power is translation by $6P$ or by $-6P$.
This obstructs a conductor-step construction on this
\emph{sectioned Jacobian pullback}. It does not exclude a
normalization which is an elliptic torsor without a section,
an additional division cover, or a different semistable model.
It does not compute an element of $\Theta_7$.

\begin{proof}
We first determine the two Hodge summands relevant to the norm.
Let $\Phi_d$ denote the $d$th cyclotomic polynomial. In the
notation of B.6.10, the rational $\Phi_2$ summand of $H^2$ has
rank three and signature $(2,1)$; the rational $\Phi_4$ summand
has rank four and signature $(2,2)$. On either of the complex
$\eta$-eigenspaces of order four the Hermitian form has
signature $(1,1)$. We will prove that each associated
one-dimensional period map is nonconstant.

Here is an exact calculation of their monodromy as $p$ moves.
The homological matrices of the original elliptic surface are
\[
 R_1=\begin{pmatrix}0&1\\-1&-1\end{pmatrix},\qquad
 R_2=\begin{pmatrix}0&-1\\1&0\end{pmatrix},\qquad
 R_0=\begin{pmatrix}1&0\\1&1\end{pmatrix},\qquad
 R_1R_2R_0=I.
\]
For the character $k=3$ of the twelve-fold cover, use the
cohomological local system with four matrices
\[
 (G_1,G_2,G_3,G_4)
       =(R_1^{-t},\,iR_2^{-t},\,R_0^{-t},\,-iI).
 \tag{B.6.20}
\]
For $k=6$ replace the character factors $(i,-i)$ by $(-1,-1)$,
giving
$(R_1^{-t},-R_2^{-t},R_0^{-t},-I)$.
Changing the deck-character convention conjugates the
order-four computation and does not change its conclusion.

These systems compute the nontrivial-character pieces of
$H^2(S_p)$ by parabolic cohomology. To check this identification,
push the elliptic fibration through the finite map $C_p\to\mathbb
P^1$. The $R^0$ and $R^2$ terms are the permutation system of
the cover; their global contributions have trivial character.
For $R^1$, good reduction at the ramification points identifies
the stalk with the invariant subspace of the displayed local
matrix. At an $I_1$ point the nodal specialization gives the
same invariant subspace. Thus its character summand is exactly
the middle extension $j_*G$, and its $H^1$ is the stated
parabolic group. The possible Leray differentials have trivial
character targets or sources, so none changes these summands.

Explicitly, a parabolic cocycle is a tuple
$z_j\in\operatorname{im}(G_j-I)$ satisfying
\[
 z_1+G_1z_2+G_1G_2z_3+G_1G_2G_3z_4=0,
\]
modulo the tuples $((G_j-I)v)_j$.
The four local image ranks are $(2,1,1,2)$ for $k=3$ and
$(2,2,1,2)$ for $k=6$. In each case the relation has rank two
and the coboundary space has rank two. The quotient dimensions
are therefore two and three. Representatives for bases are
the columns of the following matrices, grouped in two-row
blocks for $j=1,2,3,4$:
\[
 P_3=\begin{pmatrix}
 -1-i&i\\-1&0\\-1&0\\i&0\\0&-1\\0&0\\0&0\\0&0
 \end{pmatrix},\qquad
 P_6=\begin{pmatrix}
 0&2&-1\\-1&1&0\\-1&1&0\\-1&-1&0\\
 0&0&-1\\0&0&0\\0&0&0\\0&0&0
 \end{pmatrix}.
\]
They satisfy the image and relation conditions, and are
independent modulo the displayed coboundaries.

Let the Artin generator $\sigma_j$ send
$a_j\mapsto a_ja_{j+1}a_j^{-1}$ and
$a_{j+1}\mapsto a_j$. Moving the fourth puncture around the
first two gives the pure braids
\[
 \sigma_3\sigma_2\sigma_1^2\sigma_2^{-1}\sigma_3^{-1},
 \qquad \sigma_3\sigma_2^2\sigma_3^{-1}.
\]
The fourth matrix is scalar, so these braids preserve the
local representation. Evaluate them on the above cocycles
using $z(ab)=z(a)+G(a)z(b)$ and
$z(a^{-1})=-G(a)^{-1}z(a)$, then reduce by coboundaries.
For $k=3$ the resulting two matrices are
\[
 T_1=\begin{pmatrix}-1+i&1\\i&1\end{pmatrix},\qquad
 T_2=\begin{pmatrix}2i&1-2i\\i&1-i\end{pmatrix};
\]
their commutator has characteristic polynomial
$x^2-5x+1$. For $k=6$ they are
\[
 T_1=\begin{pmatrix}1&-2&1\\1&0&0\\1&-1&1\end{pmatrix},
 \qquad
 T_2=\begin{pmatrix}-1&-2&2\\1&0&-1\\-1&-1&2\end{pmatrix},
\]
whose commutator has polynomial $(x-1)(x^2-7x+1)$.
Both have an eigenvalue of absolute value greater than one.
All cocycle, quotient and word identities are checked exactly
in \texttt{maps\_component\_periods.py}; the displayed matrices
also give direct small certificates of noncompact monodromy.
A possible scalar deck factor in identifying a moving-root
basepoint cancels in the commutator.

The minimal elliptic pullbacks form a smooth projective family
as $p$ varies away from $0,1,\infty$: the two good-reduction
coordinate changes in B.6.10 work relatively, and the twelve
nodal fibres remain disjoint. Their period maps are holomorphic
by \cite[Theorem (1.1)]{mapsGriffiths68II}.
If the period in the order-four summand were constant,
monodromy would preserve its positive line and its negative
orthogonal line, hence lie in a compact unitary subgroup.
If the order-two period were constant, monodromy would
preserve the positive two-plane and negative line, again a
compact subgroup. The computed hyperbolic commutators exclude
both possibilities. Each period domain has complex dimension
one. For any fixed nonzero rational vector, being of type
$(1,1)$ is a proper period hyperplane condition, so its
preimage in this nonconstant one-dimensional family is
discrete. On the universal cover take the countable union
over rational vectors and a countable set of marking charts.
Its projection to the parameter curve is countable.
Outside that set, neither rational $\Phi_2$ nor rational
$\Phi_4$ contains a nonzero $(1,1)$ class.

Now let $p$ be outside this countable set.
The Mordell--Weil lattice embeds equivariantly into rational
$(1,1)$ cohomology by the orthogonal projection away from the
fibre and zero-section classes; see
\cite[Lemmas 8.1--8.3 and Theorem 8.9]{mapsShioda90}.
All fibres are irreducible, so its height lattice is integral,
even and torsion-free, with every nonzero vector having norm
at least two, by \cite[Theorems 8.6--8.7]{mapsShioda90}.
Since $P$ has norm two, it is primitive. The invariant
second cohomology has rank three by B.6.10; the fibre and
zero section account for two dimensions and $P$ for the
third. Hence $L_p^\eta=\mathbb ZP$.

The only eigenvalues of an operator of order twelve whose
fourth power is one have orders $1,2,4$. The last two are
absent from $L_p\otimes\mathbb Q$, by the period argument.
Consequently $L_p^{\eta^4}=\mathbb ZP$ as well. For any
$Q\in L_p$, write
\[
 (1+\eta^4+\eta^8)Q=mP.
\]
Pairing with $P$ gives $3\langle P,Q\rangle=2m$.
The integral height pairing forces $m=3r$ for an integer
$r$. Applying $1+\eta+\eta^2+\eta^3$ now gives
$NQ=12rP$. Conversely $N(P)=12P$, proving (B.6.19).
In particular $6P$ is not a norm but twice it is.

Finally any elliptic-surface automorphism covering the base
deck generator is $t_Q\eta$ or $t_Q[-1]\eta$: its value on
the zero section gives $Q$, and the remaining automorphism
fixes zero and is $\pm1$ because the elliptic $j$-function
is nonconstant. The twelfth power in the first case is
$t_{NQ}$, which cannot be $t_{\pm6P}$. In the second case
the alternating norm has image in the $\eta=-1$ eigenspace
of $L_p\otimes\mathbb Q$, which is zero. Torsion-freeness
makes that norm zero integrally. This proves the assertion
about the sectioned conductor step. Replacing the elliptic
surface by a torsor removes the zero-section argument.
B.6.13 constructs such torsor data, while B.6.14 isolates
a further obstruction to using them as the requested filling.
\end{proof}

\paragraph{Proposition B.6.13 (Proved; torsor and conductor-line data exist).}
For every $p\notin\{0,1,\infty\}$ there is a compact smooth
holomorphic genus-one fibration $\pi:T_p\to C_p$, locally
isomorphic over $C_p$ to $S_p$, and an automorphism $F$ covering
$\eta$, with
\[
 F^{12}=t_{6P}.
 \tag{B.6.21}
\]
Its torsor class can be chosen to be annihilated by four.
For $p$ outside the countable exceptional set in B.6.12 it
has no global section. There is also a holomorphic line bundle
$M$ of relative Picard class $P$ such that, for the unique
point $c_p\in C_p$ over $p$,
\[
 M\otimes F^*M^{-1}\simeq\pi^*\mathcal O_{C_p}(c_p).
 \tag{B.6.22}
\]
Thus a linearization with a simple zero at that fibre exists.
These are conductor data, not an identification of a parameter
filling with the prescribed mapping torus. The obstruction in
B.6.14 still applies to the direct one-component plumbing.

\begin{proof}
Let $G=\langle\eta\rangle\simeq\mathbb Z/12$, and let
$\mathcal E$ be the abelian sheaf of holomorphic sections of
the identity Neron group of $S_p\to C_p$ in the analytic
topology. Since all singular fibres are $I_1$, this is the
smooth locus of the elliptic fibration. The analytic
exponential sequence is
\[
 0\longrightarrow\mathcal A\longrightarrow\mathcal V
   \longrightarrow\mathcal E\longrightarrow0,
 \tag{B.6.23}
\]
where $\mathcal V$ is the coherent Lie line and
$\mathcal A=j_*\Lambda$ is the constructible period lattice.
This assertion also holds at an $I_1$ fibre. In the Tate chart
the exponential is $z\mapsto e^{2\pi iz}$; at the central
$\mathbb C^*$ fibre its kernel is the invariant rank-one
lattice. A local nonvanishing section has a holomorphic
logarithm near each value, proving local surjectivity.

We first compute the obstruction group on the quotient
orbifold $\mathcal B=[C_p/G]$, whose coarse curve is
$B=\mathbb P^1$. Write $q:\mathcal B\to B$.
The higher direct images $R^jq_*\mathcal E$, $j>0$, are
skyscraper sheaves at the three branch values and are the
cohomology of the corresponding stabilizers on section germs.
At such a point, choose a small invariant root disc. The
formal subgroup at zero is uniquely divisible and its
logarithm is a complex vector space, so its higher finite
group cohomology vanishes by averaging. Only the action on
the central elliptic curve contributes.

Here are the local calculations, with actual sections.
Engel's generator can be taken as
$P=(0,\sqrt{2}\,t^2(t-1))$. At zero the root chart
$t=z^3,x=z^4X,y=z^6Y$ has limiting curve
$Y^2=X^3+2$, section $(0,-\sqrt2)$, and stabilizer action
$(X,Y)\mapsto(\zeta_3^2X,Y)$. Its fixed group consists of
$O,(0,\sqrt2),(0,-\sqrt2)$ and has order three.
At one, use $t=1+z^4,x=z^2X,y=z^3Y$. The limit is
$Y^2=X^3-3X$, with section $(0,0)$ and action
$(X,Y)\mapsto(-X,iY)$. Its fixed group is
$\{O,(0,0)\}$ of order two. The norms on either of these
nontrivial complex-multiplication elliptic curves are zero;
hence their cyclic $H^2$ groups are respectively
$\mathbb Z/3$ and $\mathbb Z/2$, generated by the
specialization of $P$. At $c_p$, the stabilizer has order
twelve but acts trivially on the elliptic fibre, and
multiplication by twelve is onto, so its $H^2$ is zero.
The nodal fibres are unramified and have trivial stabilizers.

The restrictions of $6P$ vanish in all these local groups.
For an explicit order-three witness use $2P$.
At the order-four point $6P$ specializes to zero; the unique
local inverse of multiplication by four in the formal group
is invariant, so its norm is $6P$. At $c_p$ choose a local
holomorphic half of $P$ on the ordinary downstairs disc;
its pullback has norm $6P$. The root equations, actions and
local witnesses are checked in
\texttt{fourfolds\_B6\_torsor\_local.py}.

There is no additional degree-two obstruction on the coarse
curve. To see this without omitting a component sheaf, take
invariants in (B.6.23). The Lie invariants are coherent and
higher finite-group cohomology of them is zero. The resulting
sequence is
\[
 0\longrightarrow q_*\mathcal A\longrightarrow q_*\mathcal V
  \longrightarrow q_*\mathcal E
  \longrightarrow R^1q_*\mathcal A\longrightarrow0.
\]
The last sheaf is a finite skyscraper sheaf, and the first is
constructible on a real two-dimensional curve. Splitting this
sequence at its middle image, coherent acyclicity in degrees
at least two and $H^j(B,q_*\mathcal A)=0$ for $j>2$ give
\[
 H^2(B,q_*\mathcal E)=H^3(B,q_*\mathcal E)=0.
\]
Skyscraper sheaves have no higher cohomology. The coarse-map
spectral sequence therefore identifies restriction to the
stabilizers with an isomorphism
\[
 H^2(\mathcal B,\mathcal E)\simeq
       \mathbb Z/3\oplus\mathbb Z/2.
 \tag{B.6.24}
\]
In particular the inflation of the cyclic class $[6P]$
vanishes in this group.

The group--Cech double complex for $G$ acting on
$\mathcal E$ gives the five-term exact sequence, whose
relevant part is
\[
 H^1(C_p,\mathcal E)^G\longrightarrow
 H^2(G,H^0(C_p,\mathcal E))\longrightarrow
 H^2(\mathcal B,\mathcal E).
 \tag{B.6.25}
\]
Here $H^0(C_p,\mathcal E)=L_p$, and cyclic group cohomology
identifies the middle group with $L_p^\eta/NL_p$.
Thus $[6P]$ is the transgression of an invariant analytic
torsor class $\alpha$. In Cech terms, invariance supplies
an isomorphism $F$ of the torsor covering $\eta$; its
twelfth power is translation by an invariant section $R$,
and the transgression is $[R]$ modulo norms. Its sign is
immaterial here because $[6P]$ has order at most two.
Replacing $F$ by $t_QF$ adds $NQ$ to $R$ and therefore
arranges (B.6.21) exactly.

We justify the finite-order choice of $\alpha$.
The elliptic monodromy on $C_p$ minus its twelve nodal
points contains both $R_1$ and $R_2$: in the four-punctured
downstairs group, multiply either old loop by a power of
the new-point loop to kill its covering character. The
latter has identity elliptic monodromy and covering
character generating $\mathbb Z/12$. Filling the ramified
points kills only loops with identity elliptic monodromy.
The columns $(1,-2)^t$ of $R_1-I$ and $(-1,1)^t$ of
$R_2-I$ form a unimodular pair. Thus the global elliptic
lattice coinvariant group is zero. The exact sequence
$j_!\Lambda\to j_*\Lambda$ has skyscraper cokernel,
so compact-support duality on the punctured curve gives
\[
 H^2(C_p,\mathcal A)=H^2_c(C_p^\circ,\Lambda)=0.
\]
Consequently $H^1(C_p,\mathcal V)\to H^1(C_p,\mathcal E)$
is onto. By relative Serre duality and B.6.10 the characters
of $H^1(C_p,\mathcal V)$ are $6,3,2$ modulo twelve.
There is no invariant vector, so the group norm on this
space, and hence on $H^1(C_p,\mathcal E)$, is zero.
Every invariant class is therefore annihilated by twelve.
Replace $\alpha$ by $9\alpha$: its transgression is still
$[6P]$, while $4(9\alpha)=0$.

Represent this class by holomorphic section-valued Cech
translations. Gluing the corresponding copies of $S_p$
gives a compact smooth genus-one surface locally isomorphic
to $S_p$, with no multiple fibres. Such translations extend
over $I_1$ because they act by multiplication on the nodal
curve's normalization. Invariance and the preceding norm
adjustment give the required holomorphic $F$. If this torsor
had a section at a generic $p$, it would be the sectioned
model excluded by B.6.12, proving its nontriviality.

It remains to construct and normalize $M$. On a local
trivialization of the torsor use the line bundle
$\mathcal O(P-O)$. On overlaps these lines have the same
relative degree-zero Picard class; their difference is a
line from the base. The scalar obstruction on triple
overlaps lies in $H^2(C_p,\mathcal O^*)=0$, by the ordinary
exponential sequence and the real dimension of the curve.
Thus they glue to a line bundle $M$ of relative class $P$.
Since $F$ fixes this relative class, there is a line $J$ on
$C_p$ with
\[
 M\otimes F^*M^{-1}=\pi^*J,\qquad
 \bigotimes_{i=0}^{11}\eta^{i*}J
       \simeq M\otimes t_{6P}^*M^{-1}
       \quad\hbox{after pullback to }T_p.
\]
We verify the degree on the right for this torsor, rather
than assuming the sectioned height formula applies to it.
A torsor annihilated by four has a finite splitting cover:
lift its class through the local multiplication-by-four
sequence and trivialize the resulting finite torsion torsor
on the ordinary curve. At an $I_1$ point the added division
branches have Tate coordinates equal to a unit times
$q^{j/4}$; they extend after finite base change and minimal
resolution. Thus compactification gives a finite curve
$D\to C_p$, of degree $d$, on whose regular pullback model
the torsor has a section.

On that model, the pulled-back $M$ has degree zero on every
fibre component and differs from $\mathcal O(P_D-O_D)$ by
a base line. All pulled-back sections $P_D,6P_D$ are narrow:
the original $I_1$ sections lie in its smooth group, and
their pullbacks meet the identity components of the
resolved $I_n$ fibres. Restriction to $O_D$ therefore gives
the degree
\[
 \chi+(P_D\cdot O_D)+(6P_D\cdot O_D)-(P_D\cdot6P_D)
       =\langle P_D,6P_D\rangle=12d.
\]
The first equality is the narrow height formula and the
second is base-change functoriality
\cite[Theorem 8.6 and Proposition 8.12]{mapsShioda90}.
Dividing by $d$ shows that the line on the right has degree
twelve. Hence $\deg J=1$.

Finally the deck characters on $H^1(C_p;\mathbb C)$ are
$1,2,5,7,10,11$, each once. They follow either from the
base fixed-point traces in B.6.10 or from the three
holomorphic forms of characters $7,10,11$.
The characteristic polynomial is $\Phi_6\Phi_{12}$, so
\[
 \det(1-\eta^*\mid H^1(C_p;\mathbb Z))
      =\Phi_6(1)\Phi_{12}(1)=1.
\]
Thus $1-\eta^*$ is an automorphism of $\operatorname{Pic}^0(C_p)$.
There is a unique degree-zero line $D_0$ with
\[
 D_0\otimes\eta^*D_0^{-1}
       =\mathcal O(c_p)\otimes J^{-1}.
\]
Replace $M$ by $M\otimes\pi^*D_0$ to obtain (B.6.22).
The canonical section of $\mathcal O(c_p)$ supplies the
claimed simple zero. The lattice coinvariants, base
characters and norm arithmetic are also recorded in
\texttt{maps\_component\_periods.py}.
\end{proof}

\paragraph{Proposition B.6.14 (Proved: the one-component torsor plumbing cannot fill the required mapping torus).}
Let $W\to\Delta_r$ be a proper semistable family with smooth
complex four-dimensional total space, restricted to a small closed
disc.  Suppose that its reduced central fibre is irreducible and
that its normalization is a projective-line bundle, or the stated
conductor blow-up of such a bundle, over a genus-one torsor
$T\to C_p$ with Jacobian $S_p$ from B.6.10.  Suppose the map to
$B=\mathbb P^1$ extends over $W$ and restricts on that normalization
to the composite through $T\to C_p\to B$.
Then $H^2(\partial W;\mathbb Q)\ne0$.
Consequently this model does not fill the mapping torus of any
finite iterate of either the sphere monodromy or its
$S^3\times S^3$ smoothing monodromy.

More concretely, the proposed plumbing with conductor line
$\mathcal O_{C_p}(c_p)$ and its simple zero at $c_p$ produces
a rank-one fibre of multiplicity twelve at $p$.  A free quotient
of this chart retains that multiplicity.  Thus the analytic
torsor existence assertion does not by itself produce the
required primitive rank-one filling.

\begin{proof}
The small semistable total neighbourhood retracts to its central
fibre, as in B.6.11.  Because that fibre is reduced and irreducible,
its top homology is $\mathbb Z$, generated by its fundamental class;
the normal-crossing locus has real codimension two and introduces
no further top-dimensional class.  Therefore $H_6(W;\mathbb Q)$
is generated by this central divisor.
Poincar\'e--Lefschetz duality identifies the map
\[
 H^2(W,\partial W;\mathbb Q)\longrightarrow H^2(W;\mathbb Q)
\]
with its divisor class.  That class is zero: the central divisor
is the reduced zero divisor of the global parameter $r$.
The exact sequence of the pair hence makes
$H^2(W;\mathbb Q)\to H^2(\partial W;\mathbb Q)$ injective.

We show that the extended base class is nonzero on $W$.
The elliptic degree-one local system on $C_p$ has no global
invariants.  Indeed the cover is unramified above the original
$I_1$ point and is a normal cyclic cover.  Its monodromy group
therefore contains the original nodal transvection and all its
conjugates, in particular its conjugate by the order-three
elliptic generator.  The fixed line of the transvection is real
and one-dimensional.  An order-three elliptic generator has no
real invariant line, since its characteristic polynomial is
$T^2+T+1$.  The two fixed lines are distinct, proving the claim.

A genus-one torsor has the same rational degree-one local system
as its Jacobian: local translation changes act trivially on
cohomology.  The local smooth and nodal charts identify its first
direct image with the invariant extension of this system.
In the Leray spectral sequence of $T\to C_p$, the only possible
differential killing $H^2(C_p;\mathbb Q)$ consequently has zero
source $H^0(C_p,R^1)=0$.  It follows that
\[
 H^2(C_p;\mathbb Q)\longrightarrow H^2(T;\mathbb Q)
 \quad\hbox{is injective}.
\]
Pullback from $B$ to $C_p$ multiplies the degree by twelve and is
nonzero.  Pullback to a projective-line bundle is injective by
the bundle formula.  The conductor blow-up and its small
resolution preserve this injectivity: a degree-one proper
modification has a cohomological left inverse given by its
oriented pushforward, or here directly by the unchanged
conductor section.  Thus the class of $\mathcal O_B(1)$ pulls
back nontrivially to the normalization, and hence nontrivially
to $W$.  Its nonzero restriction to $\partial W$ proves the
first assertion.

For any diffeomorphism of $S^6$ or $S^3\times S^3$, the Wang
sequence gives zero second cohomology for its mapping torus,
since both the first and second cohomology of the fibre vanish.
This remains true for every finite iterate.  It contradicts
the preceding necessary property of the proposed filling.

The local multiplicity makes the mismatch visible without
cohomology.  At the unique point $c_p$ choose a coordinate $z$
with $t-p=z^{12}$.  The proposed conductor section has a simple
zero in $z$, so after multiplying coordinates by units its node
chart is
\[
 xy=rz.
\]
For $r\ne0$ this is smooth and the base map is
\[
 t-p=(xy/r)^{12}.
 \tag{B.T1}
\]
The divisor over $p$ is therefore twelve times each reduced local
branch.  An etale quotient preserves the order of vanishing at
a generic divisor point, so a free component-permuting quotient
does not turn (B.T1) into a primitive node.
The required D.3 filling instead has base function a unit times
$xy$ at a generic rank-one chart.  This proves the precise local
failure.  It does not exclude a model with additional central
components, a suitably ramified total-space quotient, or a
different normalization.  Those additional attachments would
require their own boundary comparison and characteristic-number
calculation; blowing up only the central fibre of the incorrect
model leaves its boundary unchanged and is not such a comparison.
The exact root-chart and multiplicity checks are
saved in \texttt{fourfolds\_B6\_torsor\_local.py}.
\end{proof}


\paragraph{Proposition B.6.15 (Proved: the free-circle filling formula).}
Let $Q$ be the oriented mapping torus of a diffeomorphism of
$S^2\times S^3$ preserving orientation and the positive generator
of $H^2$.  Choose generators $x\in H^2(Q;\mathbb Z)$ and
$y\in H^4(Q;\mathbb Z)$ so that $\langle xy,[Q]\rangle=1$, and write
\[
 x^2=ky,\qquad p_1(Q)=\ell y.
\]
Let $M\to Q$ be the principal circle bundle of Euler class $x$,
oriented as the boundary of the associated complex disc bundle
$W=D(L_x)$.  For $k\ne0$ its rational first Pontryagin class vanishes
and its normalized Eells--Kuiper invariant is
\[
 \mu(M)=\frac{3\ell^2/k-8\ell-4k-12\operatorname{sign}(k)}{2688}
       \pmod {\mathbb Z}.
 \tag{B.6.26}
\]
Moreover $\ell=4k+24r$ for an integer $r$, and equivalently
\[
 \mu(M)=\frac{k-\operatorname{sign}(k)+32r+144r^2/k}{224}
       \pmod {\mathbb Z}.
\]
For $k=0$, the condition $p_1(M)=0$ is equivalent to $\ell=0$,
and under this condition $\mu(M)=0$.
These formulas do not specify the integers for the surgery on the
fixed curve of the given family, nor the characteristic correction
of that surgery trace.
\begin{proof}
The fibre diffeomorphism acts trivially on integral cohomology, so
its Wang sequence gives the stated infinite cyclic groups in
degrees two and four.  Their product pairing is unimodular:
$y$ is the degree-three fibre generator swept around the mapping
circle.  The fibre has a unique spin structure, hence its mapping
torus is spin.  The index of its spin Dirac operator twisted by
$L_x$ is the degree-six characteristic number
\[
 \left\langle \widehat A(Q)e^x,[Q]\right\rangle
    =\frac{k}{6}-\frac{\ell}{24}\in\mathbb Z.
\]
Thus $\ell\equiv4k\pmod {24}$.

The tangent bundle of the disc bundle is stably the pullback of
$TQ\oplus L_x$, so $W$ has a spin-c structure of determinant
$L_x$ and
\[
 z=x,\qquad p_1(W)=(\ell+k)y.
\]
On its boundary $L_x$ is canonically trivialized by the unit vector,
so this structure induces a spin structure on $M$.  The Gysin
sequence gives, when $k\ne0$,
\[
 H^4(M;\mathbb Z)\cong\mathbb Z\oplus\mathbb Z/|k|,
 \qquad p_1(M)=\ell\,\pi^*y
\]
with $\pi^*y$ generating the displayed torsion subgroup.  Hence
$p_1(M)$ vanishes rationally.  When $k=0$, the pullback of $y$
is injective and primitive, which proves the asserted equivalence.

Let $U$ be the integral Thom class of the oriented complex line
bundle.  The map $H^4(W,M)\to H^4(W)$ sends $Ux$ to $x^2=ky$.
Its intersection form therefore has matrix $(k)$ and signature
$\operatorname{sign}(k)$ when $k\ne0$.  The rational relative lift
of $p_1(W)$ is $(\ell+k)Ux/k$.  Thom integration gives
\[
 \langle p_1(W)^2,[W,M]\rangle=\frac{(\ell+k)^2}{k},\qquad
 \langle p_1(W)z^2,[W,M]\rangle=\ell+k,\qquad
 \langle z^4,[W,M]\rangle=k.
\]
For the last two products one may use the relative lift $U$ of
$z$, so their definitions do not require $z^2$ to be nonzero.
Insert these values into the spin-c formula in B.6.9 to obtain
(B.6.26); substitution of $\ell=4k+24r$ gives its second form.
If $k=\ell=0$, the middle intersection map is zero, $p_1(W)=0$,
and all the displayed relative products vanish, giving $\mu(M)=0$.
The coefficient identities are checked exactly in
\texttt{code/maps\_free\_circle.py}.

For the proposed fixed-curve surgery, the disc bundle has the actual
free-circle mapping torus as boundary once its quotient bundle is
identified.  Comparison with the original sphere mapping torus
still requires its own parametrized spin surgery trace.  The
analytic conifold surgery of B.6.5 is along a different curve and
supplies no such comparison for this model.  In particular neither
vanishing of the ordinary affine obstruction after six loops nor
homotopy of the quotient return fixes the integer $\ell$.
\end{proof}

\paragraph{Proposition B.6.16 (Proved: a smooth fixed-circle replacement).}
The actual family and its even-shift closing maps admit a smooth
equivariant conifold replacement along the fixed curve of the
$\delta$ circle action.  Its fibres are $S^3\times S^3$, the circle
action on them is free, and the finite logarithmic fillings are
unchanged.  In the original ordered lattice $(g,u,w,d)$ the four
local monodromies of the resulting map to $B$ are
\[
 A_1,\quad A_2,\quad
 C_1=I+w\otimes(u^*-g^*),\quad
 C_2=I+(w-d)\otimes g^*.
 \tag{B.6.27}
\]
Their product is one.  The marked parameter return has linear part
$S=I-2d\otimes g^*$ and the same finite affine defects as in
B.6.4; the two new cusp defects are zero.  On rational $H^3$ it
preserves a two-step filtration with one-dimensional quotients,
acting as the identity on each quotient.  The resulting integral
unipotent shear is not evaluated by these assertions.
\begin{proof}
Use the following short and long bases at the cusp:
\[
 a_1=w,\qquad a_2=w-d,\qquad
 b_1=u,\qquad b_2=g+u.
\]
They form an integral basis of the fibre lattice, and the cusp
nilpotent sends $b_i$ to $a_i$ and kills both $a_i$.  In these
coordinates the normalized period matrix is
\[
 \Omega_{c,0}(q)=\frac{\log q}{2\pi i}I+C_c(q),
 \tag{B.6.28}
\]
where $C_c$ is holomorphic at zero.  The original fan has edge
directions $a_1,a_2,a_1-a_2$; thus it is the square grid with
each square subdivided along the $d=a_1-a_2$ diagonal.
This is precisely the small resolution of the product of two
infinite nodal-chain models with the same degeneration parameter.
There is one orbit of the resolved diagonal curves under the two
long-period translations.  Its quotient is the actual fixed curve
of D.2.14, not the different diagonal curve used in B.6.3.

Here is an explicit construction including its boundary matching.
Fix a small cusp disc and restrict $c$ to a compact interval
representing the parameter loop.  On a smaller disc replace
(B.6.28) by
\[
 \Omega_{c,s}(q)=
  \frac1{2\pi i}
   \begin{pmatrix}\log q&0\\0&\log(q-s)\end{pmatrix}
      +C_c(q),\qquad 0<|s|\ll1.
 \tag{B.6.29}
\]
For clarity this inner filling can be constructed before taking
any lattice quotient.  Take the product of the two infinite
nodal-chain families, with parameters $q$ and $q-s$.  Its node
charts are
\[
 xy=q,\qquad zv=q-s,
 \quad\hbox{or equivalently}\quad xy-zv=s.
 \tag{B.6.30}
\]
The two long generators shift the respective chain indices by one
and multiply the two torus coordinates by the columns of
$\exp(2\pi iC_c(q))$.  These multiplications extend on all the
chain charts, since their factors are nonvanishing.  They commute.
The usual bounded-unit lattice estimate proves that the action is
free and properly discontinuous and that its quotient is proper
over the disc: the two diagonal logarithmic lengths tend to
infinity at their respective nodes, their minimum is uniformly
large for small $q,s$, and the off-diagonal logarithms remain
bounded.  A bounded rectangle in the two chain indices is a
compact fundamental set.  This is the same estimate as the
square-grid special case of the supplied manuscript's
Theorem 4.5, Steps 1--4, with the two diagonal lengths allowed
to differ.  In particular no positivity or symmetry of $C_c$
is required.  At $s=0$ the square-grid quotient has ordinary
double points, and its indicated small resolution recovers the
given cusp, including its unit factors.

On an annulus separated from both $0$ and $s$ the difference
between (B.6.29) and (B.6.28) is the single-valued, small matrix
\[
 \frac{\log(1-s/q)}{2\pi i}E_{22}.
\]
Multiply it by a smooth radial cutoff equal to one on the inner
boundary and zero near the outer boundary.  Real nonsingularity
of the period matrix is open; after reducing $|s|$ this interpolation
therefore defines a smooth marked torus bundle on the annulus.
It agrees exactly with the inner chain quotient and with the
original cusp on their respective collars.  Glue the original
exterior, and retain both finite logarithmic fillings verbatim.
This constructs a smooth global replacement.  Holomorphicity of
the cutoff, or a global holomorphic smoothing, is not asserted.
All changes occur inside the cusp disc, and the outer peripheral
identification is the original marked one.

The construction also retains the actual parameter return.  The
period perturbation just made is independent of $c$.  The change
$c\mapsto c+h$ changes the long periods only by the integral
short vector $h d\otimes g^*$; its exponentiated lattice
multipliers are consequently unchanged, for integral $h$.
Use the existing cusp representative of the closing map, namely
the holomorphic torus unit obtained by exponentiating
\[
 \overline b_{h,c}(q)=
 \left(-\frac h2\mu(q),
 -\frac h{12}C_{c,21}(q)+h\kappa_0(q)-\frac{h^2}{24}\right)
 \tag{B.6.31}
\]
in the original short coordinates $(w,d)$.  Formula (B.6.31) is
the checked cusp formula in K8, displayed immediately after
$\ell=-s_0/12+\kappa_0(q)$ in
\texttt{S6\_02\_PUSH\_C\_closing\_cocycle.tex}; here $s_0$
denotes the additive logarithm and is distinct from the new
smoothing parameter.  Change short coordinates integrally to
$(a_1,a_2)$ and use this same unit throughout the new filling.
On (B.6.30) a torus unit acts by
\[
 (x,y,z,v)\longmapsto
   (\lambda_1x,\lambda_1^{-1}y,
    \lambda_2z,\lambda_2^{-1}v).
\]
It preserves $q$ and the actual parameter $s$, extends over both
nodes and commutes with the two chain shifts.  It also defines the
same affine translation on the interpolated annulus.  Near its
outer boundary it is the original closing map, so the original
finite comparisons and their composition law remain unchanged.
In particular the map has not been replaced by a fractional
long-period translation.  Its additive cusp lift is single-valued
on both new punctures, giving zero new affine defects.

This replacement is the usual local small-resolution-to-smoothing
surgery at the one orbit of nodes.  It is therefore surgery on
the actual null-homotopic fixed $S^2$, with its unique real
framing; D.2.14 identifies its underlying smooth fibre with
$S^3\times S^3$.  In (B.6.30) the $d=a_1-a_2$ circle has weights
$(1,-1,-1,1)$ on $(x,y,z,v)$.  Its only possible fixed point is
the origin, which is absent when $s\ne0$.  At either rank-one
central fibre it has a primitive component in the surviving
elliptic factor, so it is free there as well.  Together with the
unchanged exterior and finite fillings, where it was free, this
proves global freeness.

The primitive monodromy at $q=0$ sends $b_1$ to $b_1+a_1$;
at $q=s$ it sends $b_2$ to $b_2+a_2$.  Expressing these two maps
in the original basis gives (B.6.27), and direct multiplication
gives $C_1C_2=M_0$.  They commute with $S$.  The finite affine
defects remain $(0,1,-1,-1)$ and $(0,0,-1,1)$, as already
computed in B.6.4.

Finally the invariant-specialization calculation of B.6.3 applies
to each of these two primitive rank-one fibres and to the unchanged
finite fibres.  For $j$ the inclusion of the four-value complement
and $V=R^1g_*\mathbb Q$ there is again
$R^qg_*\mathbb Q=j_*\bigwedge^qV$.  The parabolic cocycle
calculation with the new matrices gives, in degrees $q=0,\ldots,4$,
\[
 \dim H^0(B,j_*\bigwedge^qV)=(1,1,0,1,1),\qquad
 \dim H^1(B,j_*\bigwedge^qV)=(0,0,0,0,0).
\]
The degree-one global coinvariant space is one-dimensional, so
$\dim H^2(B,j_*V)=1$.  The exact matrix calculation is recorded
in \texttt{maps\_fixed\_curve\_monodromy.py}.  Thus the only
possible Leray pieces of degree three are the one-dimensional
groups $E_2^{0,3}$ and $E_2^{2,1}$.  Since the total space is
$S^3\times S^3$, both survive and exhaust $H^3$.  The return acts
as the identity on each: on the first by direct invariant-form
calculation, and on the second because $S-I$ has image in the
vanishing coboundaries.  Saturating the rank-one integral
filtration gives an integral upper triangular return with
diagonal entries one.  The off-diagonal integer is an extension
in this filtration.  The earlier analytic smoothing had all its
$H^3$ in $E_2^{1,2}$, so its Torelli conclusion cannot be
transferred to this different surgery.
\end{proof}

\paragraph{Proposition B.6.17 (Proved: the actual quotient Pontryagin class).}
Use the fixed-circle replacement of B.6.16, the positive parameter
loop, and the mapping-torus convention
$(z,1)\sim(\Phi z,0)$ of B.6.1.  Let $M_s$ be its
seven-dimensional mapping torus and let
\[
 \pi:M_s\longrightarrow Q=M_s/S^1_\delta
\]
be the principal circle bundle, with Euler class $x$.
Orient $Q$ by putting the positive $\delta$ direction first
in the orientation of $M_s$.  The bundle over the parameter
circle has fibre $S^2\times S^3$.  The class $x$ is primitive,
and there is a unique $y\in H^4(Q;\mathbb Z)$ with
$\langle xy,[Q]\rangle=1$.  If $Z_1,Z_2\subset Q$ are the
quotients of the two swept critical elliptic curves, oriented
by their complex normal bundles, then
\[
 p_1(Q)=-2\operatorname{PD}_Q(Z_1+Z_2),\qquad
 \langle x,[Z_1]\rangle=2,\quad
 \langle x,[Z_2]\rangle=0.
 \tag{B.Q1}
\]
Consequently
\[
                       p_1(Q)=-4y.
 \tag{B.Q2}
\]
This computes the integer $\ell$ in B.6.15, with these
orientations.  It does not compute $k$ in $x^2=ky$ or the
relative Pontryagin number of the parametrized surgery trace.

\begin{proof}
The quotient of each six-dimensional fibre is
$B_5=S^2\times S^3$, by the free-circle description in
D.2.14 applied to the nullsphere surgery.  Its circle Euler
class is primitive.  One can see its generator geometrically:
the vanishing $S^3$ of the local conifold smoothing is
invariant under the free circle, and its quotient is the new
$S^2$ in $B_5$.  The restricted circle bundle is a Hopf bundle,
so its Euler number has absolute value one.  In particular
this sphere generates $H_2(B_5;\mathbb Z)$.
The return preserves the circle orientation and hence its
Euler class.  The integral Wang sequence, with
$H^1(B_5)=0$, gives
\[
 H^2(Q;\mathbb Z)\xrightarrow{\;\cong\;}H^2(B_5;\mathbb Z).
 \tag{B.Q3}
\]
It also gives $H^4(Q;\mathbb Z)=\mathbb Z$ and no torsion
in these degrees.  Integral duality proves the assertion
about $y$.

We keep the complex structures on all torus fibres and on
the local finite and nodal models.  On the interpolating
annulus in B.6.16, the period matrix is smooth rather than
holomorphic in the base.  It nevertheless defines a smooth
vertical complex structure.  An invariant horizontal
distribution extends it to an almost complex structure on
the six-dimensional fibres, making the map to $B$
pseudoholomorphic wherever it is a submersion and agreeing
with the local complex structures near every critical
fibre.  Such horizontal choices form an affine space off
the critical collars, so they can be patched there and
averaged over $S^1_\delta$.  They are made on the parameter
mapping torus itself.  Thus they respect the actual closing
map, not just a chosen trivialization over an interval.

Let $E$ be the resulting complex rank-three bundle tangent
to the six-dimensional fibres of $M_s\to S^1$.  The circle
action on $E$ is complex linear, and it is free on the base.
It descends to a complex bundle $\overline E$ on $Q$.
The vector field $\delta$ and the parameter direction are
invariant trivial real lines, so the real bundles satisfy
\[
 (\overline E)_{\mathbb R}\oplus\mathbb R
                         \cong TQ\oplus\mathbb R.
 \tag{B.Q4}
\]
In particular $p_1(Q)=p_1(\overline E)$.

Write $F:Q\to B$ for the descended base map.  At a multiple
fibre of order $m_i$, divide the derivative by its known
factor of order $m_i-1$.  Its resulting target line on $Q$
is
\[
 \overline L=F^*TB\otimes
        \mathcal O_Q\!\left(-\sum_i(m_i-1)\overline S_i\right),
 \tag{B.Q5}
\]
where $\overline S_i$ is the quotient of the reduced
multiple-fibre divisor in $M_s$.  The notation
$\mathcal O_Q(\overline S_i)$ means the complex line defined
by its oriented complex normal bundle; no complex structure
on all of $Q$ is required.
The generating sphere just described lies inside the
conifold patch, whose base image is a disc.  Therefore
$F^*[\mathrm{point}]$ evaluates to zero on that sphere.
Equation (B.Q3) proves that it is zero in $H^2(Q;\mathbb Z)$.
The local root-disc model also gives, as an equality of
integral divisor classes on $Q$,
\[
 m_i[\overline S_i]=F^*[\mathrm{point}]=0.
\]
Since $H^2(Q;\mathbb Z)$ is torsion free, every
$[\overline S_i]$ is zero.  Thus
\[
                         c_1(\overline L)=0.
 \tag{B.Q6}
\]

The determinant of $\overline E$ is also trivial in
cohomology.  It suffices by (B.Q3) to test it on the new
sphere.  Its inverse image is the local vanishing $S^3$ in
\[
                         xy-zv=s\subset\mathbb C^4.
\]
The circle weights are $(1,-1,-1,1)$.  The holomorphic
residue volume form of this hypersurface is nowhere zero
and invariant: the numerator has total weight zero and
the defining equation has weight zero.  It trivializes the
equivariant determinant on this $S^3$, hence trivializes
the descended determinant on its Hopf quotient.  Therefore
\[
                         c_1(\overline E)=0.
 \tag{B.Q7}
\]
This checks the possible determinant contribution on the
new sphere; it is not inferred from a determinant computation
only on the ordinary torus locus.

Let $V$ denote the rank-two logarithmic vertical tangent
bundle.  It is the ordinary vertical tangent bundle away
from the critical curves.  Near a rank-one node with
$q=xy$ and elliptic coordinate $z$, it is the locally free
bundle with evaluation generators
$x\partial_x-y\partial_y$ and $\partial_z$.
These local extensions glue to the ordinary kernel off
the critical loci.  They descend to $\overline V$ on $Q$.
The complex vector field generated by $\delta$ is an
invariant nowhere-zero section of $V$: the circle is free,
and at either node it has a primitive component in the
surviving elliptic direction.  Hence
\[
                         c_2(\overline V)=0.
 \tag{B.Q8}
\]

After division at the multiple fibres, the derivative
has no remaining failure of surjectivity except the two
rank-one critical loci.  Near each, the complex Koszul
resolution of $(x,y)$ identifies its cokernel with the
target line restricted to the critical curve.  These local
resolutions and the actual kernel off the critical curves
give the following identity in complex topological
$K$-theory of $Q$:
\[
 [\overline E]=[\overline V]+[\overline L]
       -[\overline L|_{Z_1}]-[\overline L|_{Z_2}].
 \tag{B.Q9}
\]
Here the last two terms mean the supported $K$-classes,
extended by their complex-normal Thom classes.  Their
degree-four Chern characters are
$\operatorname{PD}_Q(Z_1)$ and
$\operatorname{PD}_Q(Z_2)$, respectively.  The identity
is consequently valid for the smooth interpolation as
well: it is exact off the critical curves, and on their
neighbourhoods it is the displayed holomorphic Koszul
calculation.  No integrability on the annulus is used.

The degree-two part of (B.Q9), together with (B.Q6)--(B.Q7),
gives $c_1(\overline V)=0$.  Its degree-four part and
(B.Q8) then give
\[
 c_2(\overline E)=\operatorname{PD}_Q(Z_1+Z_2).
\]
Equations (B.Q4) and (B.Q7) prove the first formula of
(B.Q1) integrally.  Equality first in rational cohomology
suffices here because $H^4(Q;\mathbb Z)$ is torsion free.

It remains to fix the two Euler signs.  Use the short and
long bases of B.6.16 and let $\kappa$ denote the positive
parameter direction.  At the first node the surviving
elliptic basis is $(a_2,b_2)$, its complex orientation is
that order, and
\[
 \delta=-a_2,\qquad S b_2=b_2+2a_2.
\]
With circle generator $a_2$, the critical mapping torus
has unit-line clutching
\[
 (\beta,1,z)\sim(\beta,0,e^{2\pi i\,2\beta}z).
\]
Its Chern number is $+2$ on the torus oriented by
$(b_2,\kappa)$.  For example, a unitary connection with
local form $2\pi i\,2\kappa\,d\beta$ has curvature
$-2\pi i\,2\,d\beta\wedge d\kappa$, and the convention
$c_1=iF/(2\pi)$ gives that number.  The ambient quotient
orientation with circle $\delta=-a_2$ orients $Z_1$ as
$-(b_2,\kappa)$.  Reversing both the circle and the base
orientation leaves the Euler pairing $+2$ unchanged.
At the second node the surviving basis is $(a_1,b_1)$,
$\delta=a_1$, and the return is the identity on that
lattice.  Its Euler pairing is zero.  These are the last
two assertions of (B.Q1).  Pairing its first assertion
with $x$ proves (B.Q2).
\end{proof}


\paragraph{Proposition B.6.18 (Proved: the full quotient class has one integer left).}
Let $\overline\Psi$ be the quotient return on $S^2\times S^3$
in B.6.16--B.6.17.  Let $a$ be the Dehn twist about the primitive
$S^2$, using the generator of $\pi_3SO(3)$ whose associated
three-plane bundle over $S^4$ has Pontryagin number $4$.
Let $b$ be the same twist about an embedded sphere representing
twice that primitive class.  Then there is a unique integer $B_0$
such that
\[
 [\overline\Psi]=a^{-1-2B_0}b^{B_0},\qquad
 k=-1+6B_0,\qquad B_0=\frac{k+1}{6}.
 \tag{B.6.32}
\]
This is an equality of smooth isotopy classes.  In particular,
there is no additional torsion ambiguity in the quotient return.
The integer $B_0$ has not yet been evaluated for the given family,
and this statement does not identify its reverse-surgery image
in $\Theta_7$.
\begin{proof}
The quotient return preserves its primitive circle Euler class
and its orientation, so it acts trivially on the integral
homology of $S^2\times S^3$.  Jin's mapping-torus construction
identifies this Torelli group with
$\Omega^{\operatorname{Spin}}_6(\mathbb CP^\infty)$;
his Proposition 3.1 and Theorem 4.9 identify the latter group
with $\mathbb Z^2$ by the characteristic numbers
\[
 \left(\frac{4k-\ell}{24},\ k\right).
\]
The injectivity is explicitly proved in his Lemma 2.1 using
modified surgery and the pseudo-isotopy theorem, so this is
an isotopy classification, not just a detection of rational
invariants.  His Lemmas 4.4 and 4.8, and the table on
pp.~23--24 of version 2, give
\[
 (k(a),\ell(a))=(1,4),\qquad
 (k(b),\ell(b))=(8,8).
\]
These two twists freely generate the Torelli group
\cite[Corollary 1.3 and Theorem 4.9]{mapsJin26}.
The framing is unique because $\pi_2SO(3)=0$; the indicated
sphere embeddings are determined by their homotopy classes,
as stated immediately before his Theorem 1.2.

Write the unique quotient class as $a^A b^{B_0}$.  Additivity
of the mapping-torus bordism invariant gives
\[
                  k=A+8B_0,\qquad \ell=4A+8B_0.
\]
The actual value $\ell=-4$ from B.6.17 therefore gives
$A+2B_0=-1$, which is exactly (B.6.32).  Conversely these
equations recover $A,B_0$ from the two characteristic numbers,
so no further smooth quotient class has the same pair.
This last arithmetic is checked in
\texttt{maps\_quotient\_classification.py}.
\end{proof}

\paragraph{Proposition B.6.19 (Proved: the fixed-curve trace correction).}
For the actual equivariant surgery in B.6.16, let $Z$ be the
eight-dimensional parametrized surgery trace, with oriented
boundary $M_s\sqcup(-M_\Phi)$.  It is spin.  If $k\ne0$,
its middle relative intersection form has rank one and absolute
determinant $|k|$.  There is a primitive generator
$\alpha\in H^4(Z;\mathbb Z)$ with
\[
                 p_1(Z)=(\ell-2k)\alpha.
 \tag{B.6.33}
\]
In particular, if the actual integer $k$ is $-1$, then
\[
 \langle p_1(Z)^2,[Z,\partial Z]\rangle
                  -4\operatorname{sign}(Z)=0,
 \qquad [\Phi]=0\in\Theta_7.
 \tag{B.6.34}
\]
The last conclusion is conditional here on that still required
integer calculation; it does not follow from $\ell=-4$ alone.
\begin{proof}
First the trace exists with the actual parameter closing.
The local equation (B.6.30) is a fixed quadratic smoothing,
and the closing maps preserve both $q$ and $s$.  Their torus
unit actions preserve its small-resolution ruling.  Retraction
of this ruling-preserving quadratic group to its compact form
gives the usual equivariant $D^3\times D^4$ surgery handle,
as in the explicit $SO(4)$ trace in B.6.5.  In the present
case the $\delta$ circle fixes the $D^3$ factor and acts
freely by Hopf rotations on the punctured $D^4$ factor.
The fixed set of the parametrized trace is therefore a
$D^3$ bundle over the parameter circle.  In particular it
has the homotopy type of $S^1$.  The attaching $S^2$ has
its unique spin structure and its unique stable framing
class; the spin surgery extends the unique spin structure
of each $S^6$.  Since the seven-dimensional trace fibre
is simply connected, its spin structure is unique and
is preserved by the closing map.  It gives a spin
structure on $Z$.  Either resulting spin choice on
$M_\Phi$ has the same invariant by K3.

Let $W$ be the seven-dimensional trace fibre.  It is obtained
by adding one index-three handle to $S^6\times I$ along
a null-homotopic $S^2$.  Its nonzero reduced integral
homology groups are $H_3(W)=\mathbb Z$ and
$H_6(W)=\mathbb Z$.  The parameter return preserves
their generators: the handle orientation is preserved
by the fixed quadratic model.  The Wang sequence gives
\[
 H^1(Z)=\mathbb Z\kappa,\quad H^2(Z)=0,\quad
 H^3(Z)=\mathbb Z\beta,\quad
 H^4(Z)=\mathbb Z\alpha,\qquad \alpha=\kappa\beta.
 \tag{B.6.35}
\]
All groups here have integral coefficients.  The product
assertion follows as well from the two columns of the
spectral sequence of $W\to Z\to S^1$.

We use integral Borel equivariant cohomology for the
$\delta$ action.  Put $u=c_1(\mathbb C_1)$ in
$H^2(BS^1;\mathbb Z)$, with its sign chosen so that
its restriction on the free boundary is the Euler
class $x$ of B.6.17.  The fixed set provides a section
over $BS^1$ and over $BS^1\times S^1$ after choosing
a section of its contractible disc bundle.  Thus
$u^2$ and $u^2\kappa$ inject into equivariant cohomology.
In the Borel spectral sequence the only possible
differential from $\beta$ is $d_4\beta$ to the cyclic
group generated by $u^2$; the preceding injection forces
it to vanish.  The $d_2$ target vanishes because $H^2(Z)=0$.
Choose an equivariant lift of $\beta$ and multiply it by
$\kappa$ to lift $\alpha$.  On the fixed set the lift
of $\beta$ can only restrict to a multiple of $u\kappa$,
so this lift of $\alpha$ restricts to zero.  It follows that
\[
 H^4_{S^1}(Z;\mathbb Z)=\mathbb Zu^2\oplus\mathbb Z\alpha,
 \qquad \alpha|_{Z^{S^1}}=0.
 \tag{B.6.36}
\]
Here and below the same letter denotes the lift of the
ordinary primitive class.  No torsion extension is
possible between the two free displayed summands.

At a fixed point the normal real four-plane has the two
unit circle weights, so its equivariant first Pontryagin
class is $2u^2$.  The fixed set has no ordinary positive
even cohomology, hence supplies no additional normal
or tangent characteristic term.  Consequently
\[
                   p_1^{S^1}(Z)=2u^2+c\alpha
 \tag{B.6.37}
\]
for an integer $c$.  Forgetting the group action gives
$p_1(Z)=c\alpha$.

The relative equivariant handle calculation makes the
outgoing restriction integral and primitive.  Excision
and the equivariant Thom isomorphism for its normal
$D^4$ give
\[
 H^4_{S^1}(Z,M_s)=\mathbb ZU,\qquad
 H^5_{S^1}(Z,M_s)=\mathbb ZU\kappa.
\]
The fixed restriction of the Thom class is its normal
Euler class, namely $\varepsilon u^2$, where
$\varepsilon=\pm1$ allows the chosen real normal
orientation.  Thus $U\kappa$ maps to a class whose
fixed restriction is $\varepsilon u^2\kappa\ne0$.
The map from relative to absolute degree-five
cohomology is injective.  The exact sequence of the
pair therefore gives a surjection
\[
 H^4_{S^1}(Z)\longrightarrow H^4_{S^1}(M_s)
                                  =H^4(Q)=\mathbb Z.
\]
Its kernel is generated by $U$, whose $u^2$ coefficient
is $\varepsilon$.  Hence the restriction of $\alpha$
is primitive.  Choose its sign so that
$\alpha|_Q=y$.  Restricting (B.6.37) then gives
\[
                    \ell y=2x^2+cy=(2k+c)y,
\]
which proves (B.6.33).

Write the image of the relative Thom class in (B.6.36)
as $\varepsilon u^2+m\alpha$.  Its restriction to $Q$
is zero, so $m=-\varepsilon k$.  Forgetting equivariance
shows that the ordinary relative intersection map sends
the primitive Thom generator to $m\alpha$.  Passing
from the pair $(Z,M_s)$ to $(Z,\partial Z)$ changes
neither group in degree four, since
$H^3(M_\Phi)=H^4(M_\Phi)=0$.  Integral
Poincare--Lefschetz duality identifies the pairing of
the two primitive generators with a unit.  Thus the
middle form is a rank-one form of absolute determinant
$|m|=|k|$.  In particular, when $k=-1$ its matrix is
$(\epsilon)$ for some $\epsilon=\pm1$.

In this case B.6.17 gives $\ell=-4$ and hence $c=-2$.
The rational relative lift of $p_1(Z)$ has square
$\epsilon c^2=4\epsilon$, while the signature is
$\epsilon$.  Their difference is zero for either
choice of the common orientation sign.  The spin
cobordism formula therefore gives
$\mu(M_\Phi)=\mu(M_s)$.  Substitution of
$k=-1,\ell=-4$ into B.6.15 gives $\mu(M_s)=0$.
Finally K3 identifies $28\mu(M_\Phi)$ with the
specified class $[\Phi]\in\mathbb Z/28$, proving
the conditional conclusion (B.6.34).
\end{proof}

% Independently audited by root and maps. Integral certificate in companion .py.
\paragraph{Theorem B.6.20 (Proved).}
For the marked fixed-curve surgery of the minimal $(3,4)$ family in
B.6.16, let $Q$ be the six-dimensional free-circle quotient of the
parameter mapping torus. Its primitive Euler class $e$ satisfies
$|\langle e^3,[Q]\rangle|=1$. In the oriented convention of B.6.17,
$\ell=-4$, hence the spin congruence of B.6.15 gives $k=-1$.
Consequently the spin-trace calculation of B.6.19 gives the zero class
in $\Theta_7$ for the original minimal-family monodromy.

\emph{Proof.}
Let $g,u,w,k$ be the integral fibre lattice of the free-circle quotient,
with dual forms $G,U,W,K$.  In this lattice notation the last
letter is the parameter direction, distinct from the integer cup
coefficient $k$ in the statement.  The four homology matrices are
\[
 A_1=\begin{pmatrix}1&0&0&0\\6&0&1&-1\\-6&-1&-1&1\\0&0&0&1\end{pmatrix},\quad
 A_2=\begin{pmatrix}1&0&0&0\\0&0&-1&0\\-6&1&0&1\\0&0&0&1\end{pmatrix},
\]
\[
 C_1=I+w(u^*-g^*),\qquad C_2=I+wg^*,\qquad A_1A_2C_1C_2=I.
\]
The finite caps have orders $3,4$ and power vectors
$v_1=(1,2,-4,0)$, $v_2=(-1,-3,3,0)$.
The other two caps are the ordinary primitive nodal products after
quotienting by the free transverse circle. The exact normal-torus
Mayer--Vietoris calculation in \texttt{manifolds\_mapping\_quotient.py} gives
$H_2(Q)=\mathbb Z$, and the regular $gk$-torus maps to twice its generator.
The primitive degree-two class $x$ therefore has fibre restriction
$E_0=2G\wedge K$ after choosing the sign of $x$.

\paragraph{Closed forms on the exterior.}
Remove four discs from the oriented base sphere and choose three disjoint
arcs from the fourth boundary to the first three.  A flat torus bundle over
this surface retracts over the resulting graph. Fibrewise constant forms
and forms supported in disjoint narrow strips about these arcs give a
multiplicative de Rham model: the fibre part is the exterior algebra of
$G,U,W,K$; the first three meridian parts are one-forms with values in this
algebra, with differential $A_i^T-I$.  This assertion can also be checked
by cutting the base to a disc, using the actual linear torus gluings, and
applying de Rham Mayer--Vietoris. In particular, representatives whose
mixed base terms have disjoint support have no product containing two
such base terms.

The integral normal-torus calculation gives the following representative
of $x$ on the exterior:
\[
 E_0=2GK,\qquad
 \eta_1=-4G-W,\quad\eta_2=-3G+U+K,\quad\eta_3=0.
\]
Here a boundary path with parameter $t$ carries the two-form
$E_0+h(t)dt\wedge\eta_i$, where $h$ is a smooth bump of integral one
vanishing near the chosen seam. The notation with $h=1$ below is its
piecewise polynomial path-model shorthand. Its cube is zero, as $E_0^2=0$ and the mixed
terms can be supported in disjoint strips. Independently, if $\alpha_i$
are the extra $\delta$-rows of the actual lifted monodromy, then
\[
 \alpha_1=-2G+U+K,\quad \alpha_2=3G+W-K,\quad\alpha_3=0,
 \qquad\eta_i=-\alpha_i+(A_i^T-I)(-U).
\]
Thus this representative is the actual principal-circle curvature class
up to the stated convention for its overall sign, not a class inferred
from its vertical part alone.

\paragraph{Geometric cap restrictions.}
On a finite cap the constant one-forms $G,K$ descend rationally to the
finite affine torus quotient. Untwisting the affine path by
$x\mapsto x-tv_i/m_i$ pulls a constant form $\zeta$ back to
$\zeta-dt\wedge\iota_{v_i/m_i}\zeta$. This is an actual pullback of forms,
so it preserves products. Consequently a decomposable representative
of $x$ on each finite cap restricts to
\[
 E_0+dt\wedge\beta_i,
 \qquad \beta_1=-\tfrac23K,\quad\beta_2=\tfrac12K.
\]
This convention uses the linear boundary gluing $(x,1)\sim(A_ix,0)$;
its affine cap map is $[x,t]\mapsto[x-tv_i/m_i]$. Thus its normal
meridian is $m_i\xi+v_i$. Reversing the angular convention changes
the final orientation sign, not the absolute cube. The sign also follows
from evaluating on the exact normal tori
$(-\omega(k,v_i),j_i(N_i k))$ in the companion integral calculation.
On each nodal cap, $G$ and $K$ extend as closed one-forms: in the marking
$a_1=w,a_2=w-\delta,b_1=u,b_2=g+u$, the quotient uses the transverse
primitive circle, and $G=b_2^*$: this is the transverse circle on the first cap and
the surviving longitude of the pinched factor on the second cap. The longitude form can be made constant
near the node by an exact, monodromy-invariant change. Its products with
$K$ give a cap representative restricting to $E_0$, so $\beta_3=\beta_4=0$.
All four cap representatives have square zero.

\paragraph{Boundary primitives.}
Put $b_i=\eta_i-\beta_i$.  In the piecewise polynomial notation, a one-form
$\Theta_i(t)=\lambda_i+t b_i$ is globally defined on the boundary mapping
torus whenever $(A_i^T-I)\lambda_i=b_i$; then
$d\Theta_i=dt\wedge b_i$ is precisely exterior minus cap curvature.
For smooth representatives replace $t\eta_i$ by $H(t)\eta_i$,
where $H'=h$ and $H(0)=0,H(1)=1$, retaining the term $-t\beta_i$.
Since $\beta_i\wedge E_0=0$, this change leaves every cube integral below
unchanged. The endpoint comparison is then smooth as well. One can take
\[
\begin{array}{c|c|c}
i& b_i&\lambda_i\\\hline
1&-4G-W+\tfrac23K&(-U+W)/3\\
2&-3G+U+\tfrac12K&(-U+W)/2\\
3&0&0
\end{array}
\]
On the last boundary the based inverse word gives three successive
transported segments
\[
 -C_2^TC_1^TA_2^T\eta_1=-2G+U+K,\quad
 -C_2^TC_1^T\eta_2=3G-U-K,\quad
 -C_2^T\eta_3=0.
\]
Their sum is $G=(C_2^T-I)W$, so start with $\lambda_4=W$ and add the
successive segment integrals. The only $W$ coefficient of $\Theta_4$
remains $1$, while the two nonzero $U$ increments are $+1,-1$.
The last-boundary contribution below therefore cancels exactly;
it is unchanged if the two nonzero segments are reversed.

\paragraph{Multiplicative collar calculation.}
Across a collar with coordinate $s$ from cap to exterior use the closed
form
\[
 \Omega_i=E_{\rm cap}+d(s\Theta_i)
          =E_{\rm cap}+s\,d\Theta_i+ds\wedge\Theta_i.
\]
At its endpoints it equals the chosen cap and exterior representatives.
Both $E_{\rm cap}^2$ and $(d\Theta_i)^2$ vanish. Hence
\[
 \Omega_i^3=6s\,ds\wedge\Theta_i\wedge E_0\wedge dt\wedge b_i.
\]
Integrating $s$ gives the factor $3$. With the algebraic boundary
orientation which records the coefficient of
$\Theta_i\wedge E_0\wedge b_i$ in $G\wedge U\wedge W\wedge K$,
the companion code uses fibre orientation $G\wedge U\wedge W\wedge K$
and obtains
\[
 3[\lambda_1E_0b_1]=2,\qquad
 3[\lambda_2E_0b_2]=-3,\qquad
 3[\Theta_3E_0b_3]=3[\Theta_4E_0b_4]=0.
\]
Thus the sum has absolute value $1$. For completeness, the common
overall orientation sign can be left unassigned: moving $dt$ past $\Theta_iE_0$ contributes a minus sign if
the collar orientation is $ds\wedge dt\wedge G\wedge U\wedge W\wedge K$.
The raw algebraic sum in the code is $-1$.

\paragraph{Explicit change to the inverse mapping-torus convention.}
The calculation above uses $(x,1)\sim(Ax,0)$. The convention
$(x,1)\sim(A^{-1}x,0)$ is obtained by reversing the peripheral
parameter and changing fibre coordinates by $A$ so that the marked fibre
at the basepoint is unchanged. Explicitly,
$[x,t]\mapsto[Ax,1-t]$ is the comparison and is the identity on the marked
$t=0$ fibre. The inverse-convention cap map
$[y,\theta]\mapsto[y+\theta v/m]$ becomes
$[Ax+(1-t)v/m]=[x-tv/m]$ after applying the inverse affine deck
transformation. Thus it is exactly the negative cap term used above.
The ribbon order and all fourth-boundary segments are reversed together.
The data transform as
\[
 A'=A^{-1},\quad \eta'=-A^{-T}\eta,\quad
 \beta'=-A^{-T}\beta=-\beta,\quad\lambda'=\lambda,
 \qquad b'=-A^{-T}b=(A^{-T}-I)\lambda.
\]
These are simultaneous changes of the actual forms, not merely changes
of the displayed matrices. Since $E_0$ is invariant and $\det A=1$,
\[
 [\lambda E_0b']=-[(A^T\lambda)E_0b]
                =-[(\lambda+b)E_0b]=-[\lambda E_0b].
\]
For the multi-segment last boundary use the full reversed primitive
$\Theta'(t)=A^{-T}\Theta(1-t)$; its initial value is again $\lambda$,
and the whole collar integral changes sign. Thus reversing the common
boundary convention multiplies all four contributions by the same sign.
The absolute value $1$ is independent of this choice.

Finally, $H^2(Q;\mathbb Q)$ is one-dimensional and restriction sends its
primitive generator to $2GK$, so the global closed form constructed
above represents the actual Euler class up to its independently fixed
sign. This proves the asserted absolute cube. The independently oriented
value $\ell=-4$ and the spin congruence $\ell\equiv4k\pmod{24}$ select
$k=-1$ from the two possibilities $\{1,-1\}$. The conclusion about
$\Theta_7$ is now B.6.19.\hfill$\square$


\paragraph{Theorem B.6.21 (Proved: the original smooth monodromy is zero).}
For the original $(3,4)$ family, with its complex orientation,
the positive $w$ loop and the generator of $\Theta_7$ fixed by
$\mu=1/28$, the monodromy satisfies
\[
                         [\Phi]=0\in\mathbb Z/28.
 \tag{B.6.38}
\]
It is therefore smoothly isotopic to the identity.  Its sixfold
class is zero, and among the two resulting K7 residues $0,14$
the original class is $0$.  This statement does not yet assert
the same value for the other multiplicity pairs.
\begin{proof}
B.6.16 constructs the free-circle replacement with the original
peripheral identifications and closing map.  B.6.20 gives
$|k|=1$ by an integral-normalized multiplicative calculation on
those actual finite and nodal caps.  Independently B.6.17 gives
$\ell=-4$.  The spin congruence $\ell\equiv4k\pmod {24}$
in B.6.15 then forces $k=-1$.

The circle disc bundle consequently has invariant zero by
B.6.15.  The actual parametrized spin surgery trace has
$p_1$ coefficient $-4-2(-1)=-2$, and its rank-one middle form
is unimodular by B.6.19.  Thus its relative Pontryagin square
is four times its signature, and its normalized spin
cobordism correction is zero.  It follows that
\[
 \mu(M_\Phi)=\mu(M_s)=0.
\]
K3 compares the oriented mapping torus with the glued
homotopy sphere and identifies this invariant with
$[\Phi]/28$; this proves (B.6.38).  No framed local bubble,
choice of a scalar interpolation, or unknown reverse-surgery
correction is used in this conclusion.  The integral sign
of the intermediate cup calculation is fixed by the
independent spin congruence, so its preliminary collar
orientation convention cannot change the answer.
\end{proof}

\paragraph{Proposition B.6.22 (Proved: the actual higher-multiplicity replacement).}
Let $m=3k,n=4l$, with $k,l>0$ and $\gcd(m,n)=1$, and use
any of the admissible fixed-data sphere families of K1.  Its
positive-generator return, with period $h=2kl$, has an actual
smooth fixed-circle replacement as in B.6.16.  With
\[
 A=lN_0,\qquad B=kN_1,\qquad P=4A+3B=\pm1,
\]
its four marked affine defects can be taken to be
\[
 e_1=(0,0,-2A,0),\quad e_2=(0,-B,0,0),\quad
 t_1=2A a_1,\quad t_2=-B a_2,
 \qquad a_1=w,\ a_2=w-d.
 \tag{B.6.39}
\]
Its linear return is $S_h=I-hd\otimes g^*$.  The free-circle
quotient mapping torus has the oriented invariant
\[
                           \ell=-2h=-4kl.
 \tag{B.6.40}
\]
The spin trace and its equivariant characteristic calculation
in B.6.19 apply to this return as well.
\begin{proof}
The invariant power vectors in the question's marking are,
by the checked dictionary K1,
\[
 v_1=N_0\epsilon-\rho_0d,\qquad
 v_2=N_1\epsilon'-(\rho_1+N_1)d,
\quad \epsilon=g+2u-4w,\quad\epsilon'=g+3u-3w.
\]
The sphere condition is $4lN_0+3kN_1=\pm1$.  In particular
$N_0$ is coprime to $3k$ and $N_1$ is coprime to $4l$.
Applying the invariant covector $g^*$ to a possible finite
stabilizer equation therefore proves that the $d$ circle
is free in both logarithmic caps.  Its only fixed locus
is the same cusp $S^2$, with normal bundle and weights
$\mathcal O(-1)_{+1}\oplus\mathcal O(-1)_{-1}$.

For the fixed-data shift $h=2kl$, the explicit repeated-norm
solutions in K1, E.1.2, give the two defects $e_1,e_2$
in (B.6.39), as already checked in B.6.8.  Their combined
cusp defect is
\[
 e_0=-M_0(e_1+A_1e_2)=(2A-B)w+Bd.
\]
Make the following integral change of lifted parameter
generator, leaving the underlying parameter circle fixed:
\[
 \varkappa' =\varkappa+r,
 \qquad r=(-B,\ 2A-B,\ 2B-2A,\ 0).
 \tag{B.6.41}
\]
Since $(M_0-I)r=e_0$, the new cusp defect is zero.
This is a change of lift by an actual integral fibre
period, not a fractional-period interpolation.  In the
toric cusp cover the K1 map is a holomorphic torus unit
times an integral toric power.  Composing its lift with
this integral long-period deck transformation removes
that power.  Hence the same original quotient closing
map now has a holomorphic unit representative on the
entire cusp.  Both its nonzero limit and its parameter
dependence are retained.

Use this representative in exactly the product nodal-chain
construction (B.6.29)--(B.6.30), with the same period
matrix perturbation independent of $c$.  The unit extends
on the two node charts and preserves the actual smoothing
parameter.  The cutoff on the outer regular annulus
matches it to the unchanged original closing map.
The finite maps of K1 are also unchanged: their new
defects merely record the same change of lifted
parameter generator.  Explicitly they are
\[
 d_1=e_1+(I-A_1)r=(0,P,-P,-2A-B),\qquad
 d_2=e_2+(I-A_2)r=(0,0,-P,2A+B),
 \tag{B.6.42}
\]
and the two nodal defects are zero in this marking.
Returning to the old lift changes these last two
defects to
\[
 (C_1-I)r=2Aa_1,\qquad (C_2-I)r=-Ba_2,
\]
which proves (B.6.39) with the actual marked caps.
This also checks the peripheral product relation,
because $C_1$ fixes $a_2$ and $t_1+C_1t_2=e_0$.

Each original fibre is $S^6$ by the stipulated sphere
theorem.  The same nullsphere surgery and spin-lift
argument as in D.2.14 gives a free circle action on
$S^3\times S^3$ with quotient $S^2\times S^3$.
Thus the geometric generator and torsion-free groups
used in B.6.17 are unchanged.  Its determinant proof
uses the same invariant residue form on the new Hopf
sphere.  Every reduced multiple-fibre divisor still
has zero class, since its positive multiple is the
pullback of a base point and $H^2$ of the quotient
is torsion free.  Consequently its descended Koszul
calculation gives the identical formula
$p_1(Q)=-2\operatorname{PD}(Z_1+Z_2)$.

At the first node the surviving elliptic return is now
$b_2\mapsto b_2+h a_2$ modulo the vanishing $a_1$.
At the second it is the identity on $(a_1,b_1)$.
The preceding integral changes of lift only change
affine translations on these tori; they do not change
these linear returns.  The oriented unit-line clutch
calculation of B.6.17 therefore gives Euler pairings
$h,0$, proving (B.6.40).

Finally the replacement uses the same fixed quadratic
germ and ruling-preserving closing maps, so its
parametrized spin handle is the one in B.6.19.
Its ordinary cohomology and its fixed-set Borel Thom
calculation depend only on that handle, not on the
finite multiplicities.  Thus the trace coefficient
is again $\ell-2k_{\rm cup}$ and its middle form has
absolute determinant $|k_{\rm cup}|$.
\end{proof}

% Geometric marked split is maps B6.22; exact symbolic calculation is accompanying .py.
\paragraph{Theorem B.6.23 (Proved).}
For every higher fixed-data family with $m=3k,n=4l$,
$\gcd(m,n)=1$, and $4lN_0+3kN_1=\pm1$, the oriented cup/shear
invariant of the fixed-curve-surgery mapping torus is
\[
                         k_{\rm cup}=-kl.
\]
The independently oriented tangent invariant of B.6.22 is
$\ell=-4kl$.

\emph{Proof.} Write $A=lN_0$, $B=kN_1$, $P=4A+3B\in\{1,-1\}$,
and $h=2kl$. Distinguish the parameter generator $\varkappa$ from the
positive integer $k$. The ordinary finite matrices are the same $A_1,A_2$
as in B.6.20. The actual four affine defects from B.6.22 are
\[
 e_1=(0,0,-2A,0),\quad e_2=(0,-B,0,0),\quad
 t_1=2A a_1,\quad t_2=-B a_2,
 \qquad a_1=w,\ a_2=w-\delta.
\]
Use the integral change of parameter generator
\[
 \varkappa'=\varkappa+r,\qquad
 r=(-B,\,2A-B,\,2B-2A,\,0)\in\mathbb Z\langle g,u,w,\delta\rangle.
\]
This has determinant one. A defect transforms by
$d_i=e_i+(I-M_i)r$, giving exactly
\[
 d_1=(0,P,-P,-2A-B),\quad d_2=(0,0,-P,2A+B),\qquad
 d_3=d_4=0.
\]
The statement concerns the actual marked cap models from B.6.22;
it does not infer a filling merely from these matrices.
After quotienting by $\delta$, the two finite matrices are those in
B.6.20 with their last columns multiplied by $P$. The nodal matrices
are unchanged, and their full ordered product is the identity.

Let $G,U,W,K$ be the new dual fibre forms. Use the same common
curvature-sign convention as in B.6.20, so the Euler representative
has restriction $E_0=hGK$. Its possible overall sign relative to
the oriented principal-circle Euler class is independent of $h$:
the regular $g$-parameter torus has the same unit-line clutching,
with exponent $-h$, for every positive $h$. The actual $\delta$ rows, with the same exterior gauge $-U$
as in B.6.20, give
\[
 \eta_1=-4G-W+(P-2A-B)K,\qquad
 \eta_2=-3G+U+(2A+B)K,\qquad\eta_3=0.
\]
The power vectors are
$v_1=N_0\epsilon-\rho_0\delta$ and
$v_2=N_1\epsilon'-(\rho_1+N_1)\delta$, with
$\epsilon=g+2u-4w$, $\epsilon'=g+3u-3w$.
Thus the geometric finite cap pullbacks have mixed terms
\[
 \beta_1=-\frac{hN_0}{3k}K=-\frac{2A}{3}K,
 \qquad
 \beta_2=-\frac{hN_1}{4l}K=-\frac B2K.
\]
The central coefficients $\rho_i$ do not occur in these quotient forms.
Using $P=4A+3B$, their differences simplify to
\[
 b_1=\eta_1-\beta_1=-4G-W+\frac{2P}{3}K,
 \qquad
 b_2=\eta_2-\beta_2=-3G+U+\frac P2K.
\]
Exactly the same boundary primitives work:
\[
 \lambda_1=(-U+W)/3,\qquad\lambda_2=(-U+W)/2.
\]
The last boundary has transported segments
\[
 -2G+U+(2A+B)K,\qquad
 3G-U-(2A+B)K,\qquad0.
\]
Their sum is again $G=(C_2^T-I)W$, and the opposite $U$ increments
cancel its cube contribution. All smooth seam and inverse-peripheral
comparisons are exactly those proved in B.6.20. Consequently the raw
four collar contributions are
\[
                         h,\quad-\frac{3h}{2},\quad0,\quad0,
 \qquad \sum=-\frac h2=-kl.
\]
This is a symbolic identity, verified without any multiplicity bound.
The program is \texttt{manifolds\_mapping\_general.py} in the code directory.

For $P=1$, the parameter shear has determinant one and the oriented base,
ordinary torus fibre, free circle and positive parameter retain the same
conventions as in B.6.20. Every collar expression is a positive $kl$
multiple of the minimal expression. The calibration $k_{\rm cup}=-1$
from B.6.20 therefore gives $k_{\rm cup}=-kl$ for every such datum.
For $P=-1$, global elliptic fibre inversion sends every $v_i$ to $-v_i$
and changes $\rho_i$ to $-\rho_i$, reducing to $P=1$ while preserving
the complex orientation of the original total space and the positive
parameter. It reverses the free-circle orientation and the quotient
orientation together. Accordingly $e\mapsto-e$ and $[Q]\mapsto-[Q]$,
so both $\langle e^3,[Q]\rangle$ and
$\langle ep_1(Q),[Q]\rangle$ remain unchanged. This proves the stated
signed formula also for $P=-1$.\hfill$\square$


\paragraph{Theorem B.6.24 (Proved: all requested monodromy classes vanish).}
For every admissible coprime sphere pair $(3k,4l)$ and every
fixed twist datum of K1, the monodromy for the positive
generator $2kl$ of its parameter-period lattice satisfies
\[
                    [\Phi_{3k,4l}]=0\in\Theta_7.
 \tag{B.6.43}
\]
The identifications and orientation are those specified in
the question.  In particular this completes the numerical
part of B.6 for both the original and higher families.
\begin{proof}
Write $d=kl>0$.  B.6.22--B.6.23 compute the actual
free-circle quotient invariants as
\[
              k_{\rm cup}=-d,\qquad \ell=-4d=4k_{\rm cup}.
\]
The trace has $p_1$ coefficient $2k_{\rm cup}$, by
B.6.19 and its extension in B.6.22.  We determine the
sign of its rank-one intersection form where that sign
could affect the answer.
The two choices of spin structure on the trace restrict
to the two choices on $M_s$: restriction on $H^1(-;\mathbb Z/2)$
is an isomorphism, generated by the parameter circle.
Choose the one agreeing with the spin structure induced
by the disc-bundle spin-c filling.  The induced choice
on $M_\Phi$ does not change its invariant, by K3.

Put $k_0=k_{\rm cup}$ temporarily.  The disc bundle
$D(L_x)$ has intersection form $(k_0)$ and boundary
$M_s$, by B.6.15.  Its primitive absolute degree-four
generator $y$ restricts to the torsion class $\pi^*y$
of order $|k_0|$ on $M_s$.  The trace has form $(j)$
with $j=\pm k_0$, and its primitive absolute generator
$\alpha$ restricts to the \emph{same} $\pi^*y$,
by the Borel restriction proof in B.6.19.  Its other
boundary component $M_\Phi$ has no degree-four
cohomology and contributes no torsion linking term.

The boundary torsion linking pairing evaluated on this
fixed generator is the signed reciprocal of the
rank-one intersection coefficient, with one common
boundary-orientation sign for the two fillings.
One can see this directly by lifting the generator
rationally through the relative intersection map:
the relative generator divided by its intersection
coefficient is that lift, and evaluation modulo
$\mathbb Z$ gives the reciprocal.  The outgoing
orientation of $M_s$ is the same for the disc bundle
and the trace.  Hence $j=-k_0$ would imply
\[
                         \frac{2}{k_0}=0\quad\hbox{in }\mathbb Q/\mathbb Z.
\]
For $|k_0|>2$ this is impossible, so $j=k_0$.

In that case the normalized spin trace correction is
\[
 \mu(M_s)-\mu(M_\Phi)
 =\frac{(2k_0)^2/k_0-4\operatorname{sign}(k_0)}{896}
 =\frac{k_0-\operatorname{sign}(k_0)}{224}.
 \tag{B.6.44}
\]
The free-circle formula B.6.15, with $\ell=4k_0$, gives
exactly the same value for $\mu(M_s)$.  Their difference
is therefore $\mu(M_\Phi)=0$.

For $|k_0|=1$ both the free invariant and the trace
correction vanish for either possible sign, as already
proved in B.6.19.  For $|k_0|=2$, suppose instead that
$j=-k_0$.  The trace correction would be the negative
of (B.6.44), and subtraction from the free invariant
would give
\[
 \mu(M_\Phi)=\frac{k_0-\operatorname{sign}(k_0)}{112}
                         =\pm\frac1{112}\pmod{\mathbb Z}.
\]
K3 requires $\mu(M_\Phi)$ to lie in
$(1/28)\mathbb Z/\mathbb Z$, so this is impossible.
The trace sign is thus $j=k_0$ also in this case,
and the same cancellation proves zero.

All positive integers $d$ have now been covered.  K3
identifies the last invariant with the specified smooth
homotopy-sphere class divided by $28$, proving
(B.6.43).  The proof used the actual lifted parameter
generator and all four cap maps, and does not equate
the different multiplicity families before computing
their invariants.  The symbolic characteristic identities
and the exceptional $|k_0|=2$ denominator check are
recorded in \texttt{maps\_trace\_linking.py}.
\end{proof}

\paragraph{Corollary B.6.25 (Proved: an explicit spin filling after the computation).}
Every mapping torus in B.6.24 has a spin eight-manifold
filling diffeomorphic to $D^7\times S^1$, with one
zero-handle and one one-handle, and with
$p_1^2=\operatorname{sign}=0$.
\begin{proof}
The computed zero mapping class supplies a smooth
isotopy $h_t$ from the identity to $\Phi$, chosen
constant near its endpoints.  For the convention
$(z,1)\sim(\Phi z,0)$, the map
\[
 [z,t]\longmapsto(h_t(z),e^{2\pi it})
\]
is a well-defined oriented diffeomorphism from
$M_\Phi$ to $S^6\times S^1$.  Give $D^7\times S^1$
its product orientation; its boundary orientation is
the stipulated positive-parameter orientation of
$S^6\times S^1$.  This product is parallelizable and
has zero middle cohomology, proving the two
characteristic-number assertions.  The usual product
handle presentation has one handle in each of
indices zero and one.  Either spin structure on its
boundary extends across the product.  This filling
is a consequence of the preceding actual cap and
trace computation; its existence was not assumed
to establish the zero monodromy class.
\end{proof}
